\begingroup
\clearpage
\part{Arithmetic refinement and the complete retained class}
\noindent These complete derivations refine the preceding accepted edition.
The arithmetic scalar is evaluated through the logarithmic scale; the
original correlated rows and the full retained-class map are updated.
ARU1--15 prove their exact backward receivers. All preceding mathematical
sources remain byte-recoverable in the three cumulative successors.
\providecommand{\tightlist}{\setlength{\itemsep}{0pt}\setlength{\parskip}{0pt}}
\providecommand{\C}{\mathbb C}\providecommand{\R}{\mathbb R}
\providecommand{\Z}{\mathbb Z}\providecommand{\Q}{\mathbb Q}
\providecommand{\N}{\mathbb N}
\providecommand{\im}{\operatorname{im}}
\providecommand{\diag}{\operatorname{diag}}
\providecommand{\Gr}{\operatorname{Gram}}
\providecommand{\Gram}{\operatorname{Gram}}
\providecommand{\Res}{\operatorname{Res}}
\providecommand{\cond}{\operatorname{cond}}
\providecommand{\ExtensionTok}[1]{#1}\providecommand{\NormalTok}[1]{#1}
\providecommand{\StringTok}[1]{#1}\providecommand{\BuiltInTok}[1]{#1}
\providecommand{\ControlFlowTok}[1]{#1}

\clearpage
\begingroup
\def\R{\mathbb R}
\def\C{\mathbb C}
\def\PP{\mathcal P}
\def\BB{\mathcal B}
\def\dd{\,d}
\def\Tr{\operatorname{Tr}}
\def\Gram{\operatorname{Gram}}
\def\rank{\operatorname{rank}}

\setcounter{equation}{0}
\renewcommand{\theequation}{A02SR.\arabic{equation}}
\section*{Arithmetic correction through the logarithmic scale}
\addcontentsline{toc}{section}{Arithmetic correction through the logarithmic scale}
% BEGIN PRESERVED INPUT 01_SLOW_TAIL_AND_NONLINEAR_RECEIVER.tex SHA256 dd9cefc1ae489a80f796d7cc4a723c8f78d29a35499c6a16b31fa2fac150cc7b

\noindent This is a successor calculation. The sealed signed-arithmetic reader and
its source bytes remain unchanged. The new tail argument is the continuation
supplied in attachment \texttt{ab15652a}; the Gamma-band argument is supplied in
\texttt{Untitled 1775}, final section, and the nonlinear receiving identity in
attachment \texttt{5a756bcd}, first section. We give the full receiving proofs,
repair the tail comparison off its stated annulus, and give finite guards for
the growing-parameter recurrence. The previously proved arithmetic tail and
finite-amplitude central estimates are used at precisely their original domains.
Their full proofs are \texttt{SIGNED\_ARITHMETIC\_INTEGRATED.tex},
\texttt{OUTER\_VALUE\_RECONSTRUCTION.tex}, and \texttt{root\_comparison.tex}
in the preceding accepted edition; they are not new assumptions.

\section{The original objects and proved inputs}
Fix an actual stipulated zero quartet
\[
 \rho=\tfrac12+\delta+i\gamma,\qquad 0<\delta<\tfrac12,\quad\gamma>2,
 \qquad m_0\in\mathbb N,
\]
with each quartet member a zero of $\xi$ of multiplicity at least $m_0$.
Write $h(s)=\prod_{z\in\{\rho,\bar\rho,1-\rho,1-\bar\rho\}}(s-z)^{m_0}$,
$w_h(y)=|(2\xi/h)(1/2+iy)|^2/(2\pi)$ and $m_k=w_h^{*k}$.
For $k=4\ell+1\ge5$ retain
\begin{equation}\tag{SR1}
 e=1+k(m_0-1),\quad q=e(k+1)^2,\quad c=k/2,\quad
 \chi_k(S)=\prod_{u,v=0}^k
 [S-c-(2u-k)\delta-i(2v-k)\gamma]^e.
\end{equation}
The coordinate map is $P(S)\mapsto P(c+iy)$, with inverse substitution
$y=(S-c)/i$. Its monic relation is $Q_k(y)=i^{-q}\chi_k(c+iy)$.
This map is invertible on every degree space and on every remainder quotient;
no root or lower relation is discarded. Set
\[
 \sigma(y)=\frac{|\Gamma(1/4+iy/2)|^2}{2\pi},\qquad
 \int\sigma=\sqrt{2\pi},\qquad \alpha=\pi/2.
\]
For $J_N=\operatorname{rem}_{\chi_k}$ and the full source Gram $H_N[\nu]$,
\begin{equation}\tag{SR2}
 G_N[\nu]=(J_NH_N[\nu]^{-1}J_N^*)^{-1},\qquad
 \BB_k[\nu]=\log\det G_{q-1}+\log\det G_q
             -\log\det G_{2q-1}-\log\det G_{2q}.
\end{equation}
All these quotients have rank $q$. Thus $G_N[A\nu]=AG_N[\nu]$ gives
$\BB_k[A\nu]=\BB_k[\nu]$ for every scalar $A>0$.
In any fixed complement to $Q_k\PP_{N-q}$, completion of the square gives
\begin{equation}\tag{SR3}
 \log\det G_N[\nu]=\log\det H_N[\nu]
 -\log\det\Gram_{Q_k\PP_{N-q}}[\nu]+c_N.
\end{equation}
Here $c_N$ is the logarithmic square of the fixed frame determinant; it cancels
between any two densities. For $N=q-1$ the relation is empty, and for $N=q$ it
has rank one. Formula (SR3) proves the exact map between the full source and
the original attained metric, including their cross blocks.

Put $M(\theta)=\int e^{\theta y}w_h(y)\dd y$, $\mu=(\log M)'$,
$V=\mu'$, $\mathfrak M=M(\alpha)$, $b=\mu(\alpha)$, and
\begin{equation}\tag{SR4}
 I_h=\int_0^\alpha\mu(\theta)^2\dd\theta,\quad d=4m_0-2,\quad
 c_\zeta=2\gamma_{\rm E}-\log(2\pi),\quad
 C_\Gamma=2\mathcal L(2,\pi)-8\log2+4>0.
\end{equation}
All moments through order $9/4$ of the boundary-tilted probability are finite;
in particular $M,M',M''$ extend to the closed tilt interval. The previous
original-source proof establishes $I_h>0$.
Let $\theta(t)=\mu^{-1}(t)$ on $[0,b]$ and
\[
 D(t)=\log\mathfrak M-\log M(\theta(t))-(\alpha-\theta(t))t,
 \qquad W_k(y)=kD(|y|/k)1_{|y|\le kb}.
\]
The exact area is $\int W_k=k^2I_h$. Define
\[
 \omega_k=\sigma\quad(|y|\le kb),\qquad
 \omega_k=\mathfrak M^{-k}m_k\quad(|y|>kb).
\]
The previous finite-amplitude theorem, with its full moving-family contour
proof and the sign correction SAI1, is
\begin{equation}\tag{SR5}
 \left|\BB_k[m_k]-\BB_k[\omega_k]
       -\frac{\log2}{\pi}I_hk^2\right|
 \le E_{\rm cen}:=A_h\left[
 \frac{k^2}{\log^3(q+2)}+k\log^2(q+2)+1+\frac{k^4}{q^2}\right].
\end{equation}
This evaluates a finite change of metric, not only a derivative at its first
endpoint. Its finite Neumann guard and Robin error are SAI7--10.

The proved averaged tail, in the same physical $y$ coordinate, is
\begin{equation}\tag{SR6}
 \int_T^{2T}\left|
 \frac{\mathfrak M^{-k}m_k(y)}
 {\dfrac{k}{\sqrt\pi\mathfrak M}(\log y+c_\zeta)
 (1+y^2)^{-d}\sigma(y)}-1\right|\dd y
 \le C_hT^{199/200},\qquad q^{9/10}\le T\le128q.
\end{equation}
Its proof retains both large-summand complements and the original zeta
mean square, using Ivi\'c's established increment bound \cite{a02sr:Ivic}.
The outer-only pointwise bounds are
\begin{equation}\tag{SR7}
 c_-k^{-1/2}(1+y^2)^{-M_*}\le\omega_k/\sigma
 \le c_+k^{p+1},\qquad M_*=21+4m_0,\quad p=8m_0-9/2.
\end{equation}
These are the proved endpoint/local-limit and full three-factor source bounds,
not a pointwise lower bound for $|\zeta|^2$.

\section{A complete transport of the logarithmic tail}
Write $L=\log q$, $L_c=L+c_\zeta$, $\epsilon=(2L_c)^{-1}$,
$a=d-\epsilon$, and
\begin{equation}\tag{SR8}
 \tau_a=(1+y^2)^{-a}\sigma,\quad
 A_q=\frac{kL_c}{\sqrt\pi\mathfrak M}q^{-1/L_c},\quad
 T=q/L^6,\quad \mathfrak A=\{T\le |y|\le128q\}.
\end{equation}
Finite guards are $L_c>1$, $6\log L+\log128\le L_c/2$,
$T\ge\max(q^{9/10},kb,1)$, and all the prior guards of (SR5)--(SR7).
Each holds eventually for every fixed $h$. For $y\in\mathfrak A$, put
$x=\log(|y|/q)$. Direct cancellation, with all scalars retained, gives
\begin{equation}\tag{SR9}
 s_q(y):=\frac{\dfrac{k}{\sqrt\pi\mathfrak M}
 (\log|y|+c_\zeta)\tau_d(y)}{A_q\tau_a(y)}
 =\left(1+\frac{x}{L_c}\right)
 \exp\left[-\frac{x}{L_c}
 -\frac{\log(1+y^{-2})}{2L_c}\right].
\end{equation}
For $|z|\le1/2$, integration of $-z/(1+z)$ from zero to $z$ proves
$-z^2\le\log(1+z)-z\le0$. Hence
\begin{equation}\tag{SR10}
 e^{-\eta}\le s_q\le1\quad\hbox{on }\mathfrak A,
 \qquad \eta=\frac{(6\log L+\log128)^2}{L_c^2}
                   +\frac1{2L_cT^2}.
\end{equation}
The formula is not positive on the whole real line. The globally positive
comparison needed below is explicitly
\begin{equation}\tag{SR11}
 \bar s_q=s_q1_{\mathfrak A}+1_{\mathfrak A^c},\quad
 \nu_0=A_q\bar s_q\tau_a,\quad
 \nu_1=\omega_k1_{\mathfrak A}+A_q\tau_a1_{\mathfrak A^c}.
\end{equation}
Thus $e^{-\eta}\le\bar s_q\le1$ globally, and
$\nu_1=a_k\nu_0$ with $a_k$ equal to the actual arithmetic ratio in (SR6)
on $\mathfrak A$ and one off it. This repairs the incoming proof's unqualified
global use of $s_q$ without changing any original density.

We detail the finite-dimensional comparisons. The monic Gamma recurrence
has off-diagonal orthonormal entries $\sqrt{n(n-1/2)}\le n$.
Consequently $\|yP\|_\sigma\le2(D+1)\|P\|_\sigma$ on $\PP_D$.
Jensen applied to the probability $|P|^2\sigma/\|P\|^2$ yields
\begin{equation}\tag{SR12}
 \int(1+y^2)^{-M}|P|^2\sigma\dd y
 \ge[1+4(D+1)^2]^{-M}\|P\|_\sigma^2.
\end{equation}
Together with (SR7), $a\in[d-1/2,d]$, and the explicit $A_q$, this proves
that for some fixed $B=B_h$ and all four measures
$\omega_k,\nu_0,\nu_1,A_q\tau_a$,
\begin{equation}\tag{SR13}
 q^{-B}\|P\|_\sigma^2\le\|P\|_\nu^2\le q^B\|P\|_\sigma^2
 \quad(\deg P\le2q),\qquad 0<\nu/\sigma\le q^B.
\end{equation}
Increasing $B$ absorbs fixed constants after a finite $h$-dependent threshold.
The lower bound is an integrated form bound; the upper density bound follows
pointwise because $a>0$. They hold for every original source and relation
subspace and, by taking the minimum over a fixed fibre, for its attained quotient.

The exact generating function, in the original convention, is
$(1+t^2)^{-1/4}e^{z\arctan t}$ \cite{a02sr:DLMF}. Cauchy's coefficient inequality
on $|t|=D/(D+1)$, $n!/(1/2)_n\le2n+1$ and
$|\arctan t|\le\frac12\log(2D+1)$ prove
\begin{equation}\tag{SR14}
 K_D^\sigma(z,\bar z)
 \le\frac{e^2}{\sqrt{2\pi}}(D+1)^2\sqrt{2D+1}(2D+1)^{|z|}.
\end{equation}
Indeed $r^{-2n}\le e^2$ and
$|1+t^2|^{-1/2}\le(D+1)/\sqrt{2D+1}$; summing $2n+1$ gives
$(D+1)^2$. For $P\in y^r\C[y]$, the maximum principle for $P(z)/z^r$
on $|z|=2T$, followed by integration over $[-T,T]$, gives
\begin{equation}\tag{SR15}
 \int_{-T}^T|P|^2\sigma\dd y
 \le e^2(D+1)^2\sqrt{2D+1}(2D+1)^{2T}4^{-r}\|P\|_\sigma^2.
\end{equation}
Take $D=2q$ and
\begin{equation}\tag{SR16}
 r=\left\lceil\frac{
 \log[e^2(2q+1)^2\sqrt{4q+1}]+2T\log(4q+1)+(4B+60)L}
 {\log4}\right\rceil.
\end{equation}
Then $r=O_h(q/L^5)$ and the multiplier in (SR15) is at most $q^{-4B-60}$.
The further finite guard $r+d<q$ holds eventually.

For each actual space $\mathcal S=\PP_N$ or $Q_k\PP_{N-q}$ in (SR3), set
\begin{equation}\tag{SR17}
 \mathcal S^\circ=\mathcal S\cap(y+i)^d y^r\C[y],\qquad
 0\longrightarrow\mathcal S^\circ\longrightarrow\mathcal S
 \longrightarrow\mathcal S/\mathcal S^\circ\longrightarrow0.
\end{equation}
The codimension is at most $r+d$: it is the rank of the remainder map modulo
$(y+i)^dy^r$. This construction works also for the rank-one or empty relation.
For $P\in\mathcal S^\circ$, (SR13)--(SR16) bound changes of density inside
$[-T,T]$ relative to its full norm by $2q^{-2B-60}$.
For the other integration region, repeated multiplication gives
\begin{equation}\tag{SR18}
 \int_{|y|>128q}|P|^2\sigma\dd y
 \le\left(\frac{2(3q+1)}{128q}\right)^{2q}\|P\|_\sigma^2
 \le256^{-q}\|P\|_\sigma^2.
\end{equation}
The last inequality holds for $q\ge1$. Both exterior measures are bounded
by $q^B\sigma$, so their relative discrepancy is at most $2q^{2B}256^{-q}$.
Let $\zeta_q=2q^{-2B-60}+2q^{2B}256^{-q}$ and impose $\zeta_q<1/2$.
We obtain, for every $\mathcal S^\circ$,
\begin{equation}\tag{SR19}
 \left|\log\frac{\det\Gram_{\mathcal S^\circ}[\omega_k]}
 {\det\Gram_{\mathcal S^\circ}[\nu_1]}\right|
 \le2\dim(\mathcal S^\circ)\zeta_q.
\end{equation}
To restore the complementary directions, choose any fixed splitting of
(SR17). The determinant factors exactly into the restricted determinant and
the determinant of the Schur minimum on $\mathcal S/\mathcal S^\circ$.
Taking the infimum in (SR13) preserves both form inequalities, whence the
logarithmic determinant comparison between any two quotient metrics has
absolute value at most $2B(r+d)L$. Thus no discarded central or tail direction
is hidden in (SR19).

Division $P=(y+i)^dR\mapsto R$ is a linear isomorphism of $\mathcal S^\circ$
onto a fixed subspace of $\PP_{2q-d}\cap y^r\C[y]$. Under this isomorphism
the reference density $\nu_0$ is
$A_q\bar s_q(1+y^2)^\epsilon\sigma$. Its scalar $A_q$ cancels in the
projection kernel. It is at least $e^{-\eta}\sigma$ globally, and at most
$C\sigma$ on $\mathfrak A$, since
$\epsilon\log(1+(128q)^2)\le3$ after a finite guard.
The variational formula $K(y,y)=\sup_{P\ne0}|P(y)|^2/\|P\|^2$ and the
already proved whole-line bound
\begin{equation}\tag{SR20}
 K_D^\sigma(y,y)\sigma(y)\le20\log(D+2)
\end{equation}
give $K_{\mathcal S^\circ}^{\nu_0}(y,y)\nu_0(y)\le C L$ on $\mathfrak A$.
The proof of (SR20) uses Koelink--Van der Jeugt's exact Poisson kernel
\cite{a02sr:KVJ}, with $x=y/2$ and its density factor $1/2$, followed by Euler's
hypergeometric integral; the full proof is OR18--20 in the accepted predecessor.

Dyadic use of (SR6), including evenness for $y<0$, proves
\begin{equation}\tag{SR21}
 E_q:=\int_{\mathfrak A}|a_k-1|\dd y\le C_hq^{199/200},\qquad
 \int_{\mathfrak A}|\log a_k|\dd y\le(2+2BL)E_q.
\end{equation}
For the second bound, the pointwise bounds (SR7) give
$q^{-B}\le a_k\le q^B$ on the annulus after increasing $B$.
On $|a_k-1|\le1/2$ use $|\log a_k|\le2|a_k-1|$; its complement has
measure at most $2E_q$. The dyadic sum is geometric, with no extra $L$ factor.

For a fixed finite space $V$ and positive reference $\nu$, in a
$\nu$-orthonormal frame, expansion of the full determinant gives
\[
 \det\Gram_V[a\nu]=\frac1{t!}\int
 |\det(P_i(y_j))|^2\prod_{j=1}^ta(y_j)\dd\nu(y_j),\quad t=\dim V.
\]
The same expression with $a=1$ is one. Jensen for the logarithm and
$\log\det A\le\Tr(A-I)$ therefore prove
\begin{equation}\tag{SR22}
 \int\log a\,K_V^\nu\dd\nu
 \le\log\frac{\det\Gram_V[a\nu]}{\det\Gram_V[\nu]}
 \le\int(a-1)K_V^\nu\dd\nu.
\end{equation}
Applying this with (SR20)--(SR21) costs $C_hq^{199/200}L^2$ on
$\mathcal S^\circ$. Its quotient in (SR17) costs at most $2B(r+d)L$.
Finally (SR10)--(SR11) compare $\nu_0$ with $A_q\tau_a$ on the whole space
with determinant cost at most $\dim\mathcal S\,\eta$. This is a uniform
form bound and introduces no kernel logarithm.
Summing the four full source and four full relation contributions in (SR3),
and using $\BB_k[A_q\tau_a]=\BB_k[\tau_a]$, proves
\begin{equation}\tag{SR23}
 |\BB_k[\omega_k]-\BB_k[\tau_a]|
 \le C_h\left[q\eta+q^{199/200}L^2+q/L^4+q\zeta_q\right].
\end{equation}
All original attained minima and both complementary integration regions have
now been explicitly transported.

\section{The moving exponent and the next coefficient}
For a fixed integer $j\ge1$, the monic frame
$1,\ldots,y^{j-1},(y+i)^j,y(y+i)^j,\ldots,y^{N-j}(y+i)^j$
has determinant one. The last block with density $\tau_j$ is exactly the
Gamma Gram through degree $N-j$, and its complementary attained jet quotient
has fixed dimension $j$ and logarithmic determinant $O_j(\log(N+2))$.
This last bound is the proved two-sided finite-jet estimate in OR28, derived
from the Gamma recurrence and evaluation kernel. Therefore, with the original
monic norms $\gamma_n=\sqrt{2\pi}(2n)!/4^n$,
\begin{equation}\tag{SR24}
 \log\frac{\det H_N[\tau_j]}{\det H_N[\sigma]}
 =-j\log\gamma_N+O_j(\log(N+2)).
\end{equation}
The full-root comparison RC1--14 and positive-ensemble free-energy proof OR31--40
give, at $n=q,q+1$,
\begin{equation}\tag{SR25}
 \log\frac{\det\Gram_{Q_k\PP_{n-1}}[\tau_j]}
 {\det\Gram_{Q_k\PP_{n-1}}[\sigma]}
 =-2jnL-jq\mathcal L(2,\pi)+O_{h,j}(\sqrt{qL}).
\end{equation}
This formula uses the actual complete polynomial $Q_k$, not the monomial
$y^q$. Its uniform fixed-neighbourhood free-energy estimate covers all three
integers $j=d-1,d,d+1$.

For every fixed finite polynomial space $V$, define
$f_V(a)=\log\det\Gram_V[(1+y^2)^{-a}\sigma]$.
The full positive determinant integral shows that $f_V''$ is the variance of
$\sum_i\log(1+y_i^2)$ and is nonnegative. Thus, for $0\le\epsilon\le1$,
\begin{equation}\tag{SR26}
 -\epsilon[f_V(d+1)-f_V(d)]
 \le f_V(d-\epsilon)-f_V(d)
 \le-\epsilon[f_V(d)-f_V(d-1)].
\end{equation}
Substitution of (SR24)--(SR25) proves their linear formulas at
$a=d-\epsilon$, uniformly for $0\le\epsilon\le1/2$, with the same error
orders. This uses convexity of the individual positive Grams before taking
any signed difference; no error is differentiated.
The low rank-one relation has logarithmic comparison $O_h(\log q)$ by
(SR12) applied to $P=Q_k$; the empty relation has determinant one.
Stirling's formula \cite{a02sr:DLMFGamma} applied to the four source signs gives
\[
 4aqL+(8a\log2-4a)q+O_h(L).
\]
The two high relations, including their subtraction in (SR3), give
$-4aqL-2aq\mathcal L(2,\pi)+O_h(\sqrt{qL})$.
Their sum is the exact cancellation
\begin{equation}\tag{SR27}
 \BB_k[\tau_{d-\epsilon}]-\BB_k[\sigma]
 =-(d-\epsilon)C_\Gamma q+O_h(\sqrt{qL}).
\end{equation}
Taking $\epsilon=1/(2L_c)$, combining (SR23), (SR27) and (SR5), and using
$q\ge(k+1)^2$, proves
\begin{theorem}
On the original fixed-packet domain, as $k=4\ell+1\to\infty$,
\begin{equation}\tag{SR28}
 \delta_{k,1}^{\rm ar}:=\BB_k[m_k]-\BB_k[\sigma]
 =-dC_\Gamma q+\frac{\log2}{\pi}I_hk^2
 +\frac{C_\Gamma q}{2(\log q+c_\zeta)}+R_{h,k},
\end{equation}
where, for all the stated finite guards,
\begin{equation}\tag{SR29}
 |R_{h,k}|\le E_{\rm cen}
 +C_h\left[q\eta+q^{199/200}L^2+q/L^4+q\zeta_q+\sqrt{qL}\right].
\end{equation}
In particular the denominator may be replaced by $\log q$ at additional
cost $|C_\Gamma c_\zeta|q/(2LL_c)$, and the total error is at most
\begin{equation}\tag{SR30}
 C_hq\left[\frac{(1+\log\log q)^2}{(\log q)^2}
                  +q^{-1/200}(\log q)^2\right]
 =o_h(q/\log q).
\end{equation}
\end{theorem}
For $m_0=1$, set $d_0=-2C_\Gamma+(\log2/\pi)I_h$; for $m_0\ge2$ set
$d_0=-dC_\Gamma$. Since $q-k^2=2k+1$ in the first case and
$k^2=o(q/\log q)$ in the second,
\begin{equation}\tag{SR31}
 \lim\frac{\log q}{q}(\delta_{k,1}^{\rm ar}-qd_0)=C_\Gamma/2.
\end{equation}
The limits with $\log k$ are $C_\Gamma/4$ and $C_\Gamma/6$, respectively.
The theta-source proof delivered with this continuation proves $d_0>1$ on
every actual simple quartet. Thus the incoming slow-tail note's formal zero
threshold is not attained on that original simple-source domain. Formula
(SR31) is valid there nonetheless; it refines its positive leading value.

\section{Growing Gamma bands with finite recurrence guards}
For $s\in\{1,k\}$, retain $\lambda=s/4$, $M_s=(2\pi)^{s/2}$,
\begin{equation}\tag{SR32}
 w_s(y)=\frac{M_s2^{2\lambda}}{4\pi\Gamma(2\lambda)}
 |\Gamma(\lambda+iy/2)|^2,\quad
 p_{n,s}(y)=n!P_n^{(\lambda)}(y/2;\pi/2),\quad
 h_{n,s}=M_s n!(2\lambda)_n.
\end{equation}
These retain the physical coordinate and source mass. At $s=1$,
$M_1 2^{1/2}/(4\pi\Gamma(1/2))=1/(2\pi)$, so $w_1=\sigma$ exactly.
The classical generating function and recurrence \cite{a02sr:DLMF} give, for
$f_n=\sqrt{w_s}p_{n,s}/\sqrt{h_{n,s}}$,
\begin{equation}\tag{SR33}
 a_{n+1}f_{n+1}+a_nf_{n-1}=yf_n,\qquad
 a_n=\sqrt{n(n+2\lambda-1)},\qquad
 \sum\frac{p_{n,s}}{n!}t^n=(1+t^2)^{-\lambda}e^{y\arctan t}.
\end{equation}
Hereafter $y,\lambda$ are fixed while constructing the exact solution at
infinity; the resulting estimates have explicit uniform dependence on them.
Set $P=1+\lambda+|y|$, $x_n=n+\lambda$, $v=-1/2-iy/2$, and
$u_n=i^nx_n^v$. Define
\[
 F(z)=\sqrt{(1+z/2)^2-(\lambda-1/2)^2z^2}\,(1+z)^v.
\]
On $|z|\le1/(2P)$ the square root is the analytic branch equal to one at
zero. Factoring its argument into
$(1+(1-\lambda)z)(1+\lambda z)$ verifies that neither factor vanishes and
that its modulus is at most two. Also
$|(1+z)^v|\le\exp(|v||z|/(1-|z|))\le e$.
Thus $|F(z)|\le6$ there. Its odd coefficient starts with
$F(z)-F(-z)=(2v+1)z+\cdots=-iyz+\cdots$.
Cauchy's coefficient bound, for $x_n\ge4P$, gives
\begin{equation}\tag{SR34}
 |F(x_n^{-1})-F(-x_n^{-1})+iy/x_n|
 \le 192P^3/x_n^3,
\end{equation}
where the geometric odd tail is bounded by a geometric tail over all degrees.
Since $a_{n+1}=\sqrt{(x_n+1/2)^2-(\lambda-1/2)^2}$ and the corresponding
formula with $-1/2$ holds for $a_n$, this proves
\[
 |a_{n+1}u_{n+1}+a_nu_{n-1}-yu_n|
 \le192P^3x_n^{-2}|u_n|.
\]
For the two-column matrix
$U_n=\left(\begin{smallmatrix}u_n&\bar u_n\\u_{n-1}&\bar u_{n-1}\end{smallmatrix}\right)$,
the phase difference is
$\pi/2-(y/2)\log(x_n/(x_n-1))$. Under $x_n\ge4P$ it differs from
$\pi/2$ by less than $1/4$. Direct inversion therefore yields
$\|U_n\|\le3x_n^{-1/2}$ and $\|U_n^{-1}\|\le3x_n^{1/2}$.
If $T_n$ is the exact recurrence transfer matrix, the coefficient error
\begin{equation}\tag{SR35}
 B_n=U_{n+1}^{-1}T_nU_n-I
 \quad\hbox{satisfies}\quad \|B_n\|\le10^4P^3/x_n^3.
\end{equation}
Indeed the first row residual is divided by $a_{n+1}\ge x_n/2$, and the
second row is exact; the displayed matrix bounds absorb the remaining factors.
Let
\begin{equation}\tag{SR36}
 H_n=10^5P^3/(x_n-1)^2,\qquad x_n\ge4P,\quad H_n\le1/4.
\end{equation}
Then each $\|B_n\|<1/2$ and
$2\sum_{j\ge n}\|B_j\|\le H_n$ by the integral bound for $\sum x_j^{-3}$.
Products of $(I+B_j)^{-1}$ converge, and their deviation from identity is
at most $e^{H_n}-1$. This constructs and bounds the exact solution with each
prescribed limiting coefficient vector.

For completeness, the limiting vector of the original polynomials follows
from the two singular factors of (SR33):
$(1-it)^{-\lambda+iy/2}(1+it)^{-\lambda-iy/2}$.
At $t=i$ the analytic factor has value $2^{-\lambda+iy/2}$ and the singular
coefficient is
$i^{-n}n^{\lambda+iy/2-1}/\Gamma(\lambda+iy/2)$; the other endpoint is its
conjugate. This coefficient follows directly from
$[t^n](1-t)^{-z}=\Gamma(n+z)/(\Gamma(z)\Gamma(n+1))$ and the fixed-parameter
Gamma ratio limit \cite{a02sr:DLMFGamma}.
Subtracting finitely many Taylor terms of each analytic factor makes the
remaining coefficient $o(n^{\lambda-1})$: choose the subtraction order larger
than $2\lambda+3$, and integrate the residual Fourier coefficient twice on
the circle away from the endpoints. The subtracted endpoint terms of order
at least one are $O(n^{\lambda-2})$. This proves the limiting vector for each
fixed pair $(\lambda,y)$. Combining its absolute factor with (SR32) gives
exactly $1/(2\sqrt\pi)$ in each column, with phase
$\arg\Gamma(\lambda+iy/2)-(y/2)\log2$.
Uniformity comes from (SR35)--(SR36), not from a fixed-parameter error term.
In particular, with
\[
 g_n(y)=\frac{1}{\sqrt{\pi x_n}}
 \cos\left(\frac{\pi n}{2}-\frac y2\log(2x_n)
                  +\arg\Gamma(\lambda+iy/2)\right),\quad
 E_n(y)=\frac{4(e^{H_n}-1)}{\sqrt{\pi x_n}},
\]
we have the finite uniform estimate
\begin{equation}\tag{SR37}
 |f_n-g_n|\le E_n.
\end{equation}

The original even integer $q$ has the endpoint weights
$c_0=c_q=1$, $c_r=2$ for $0<r<q$. Define
$\mathsf P_{q,s}(y)=\sum_{r=0}^qc_rf_{q+r}(y)^2$.
The nonoscillatory half of $g_n^2$ equals
$\frac1{2\pi}\sum c_r/(q+r+\lambda)$.
The elementary trapezoid error is bounded by
$q/[6\pi(q+\lambda)^3]$, and its integral is
$\pi^{-1}\log((2q+\lambda)/(q+\lambda))$.
For the oscillatory half put
$a(x)=e^{-iy\log(2(x+\lambda))}/(x+\lambda)$.
Exactly,
\begin{equation}\tag{SR38}
 \sum_{r=0}^qc_r(-1)^ra(q+r)
 =\sum_{j=0}^{q/2-1}
 [a(q+2j)-2a(q+2j+1)+a(q+2j+2)].
\end{equation}
Twice integrating $a''$ over unit intervals bounds its absolute value by
$q(1+|y|)(2+|y|)/[2(q+\lambda)^3]$.
Combining these finite bounds and (SR37) proves
\begin{equation}\tag{SR39}\begin{aligned}
 |\mathsf P_{q,s}(y)-\log2/\pi|
 \le{}&\frac{\lambda}{\pi q}+
 \frac{q}{6\pi(q+\lambda)^3}
 +\frac{q(1+|y|)(2+|y|)}{4\pi(q+\lambda)^3}\\
 &+\sum_{r=0}^qc_r\left[
 \frac{2E_{q+r}}{\sqrt{\pi x_{q+r}}}+E_{q+r}^2\right].
\end{aligned}\end{equation}
For $|y|\le kb$ and $s\in\{1,k\}$ all guards (SR36) hold eventually,
uniformly, since $P\le C_bk$ and $q\ge(k+1)^2$. The right side of (SR39) is
\begin{equation}\tag{SR40}
 O_b(k/q+k^3/q^2).
\end{equation}
This proves the entire growing-interval band assertion in the incoming note,
including its original masses, Gamma order, coordinate and endpoint weights.

\section{The combined mixed and nonlinear term}
Use the actual logarithm
\[
 H_k(y)=\log\frac{m_k(y)}{\mathfrak M^k\sigma(y)},\quad
 H_C=1_{[-kb,kb]}H_k,\quad H_O=H_k-H_C,
 \quad F(t,r)=\BB_k[\sigma e^{tH_C+rH_O}].
\]
The central local-limit calculation ACW/L in the accepted provider gives
\begin{equation}\tag{SR41}
 H_C=-W_k+R_k,\quad \int|R_k|\dd y\le C_hk,\quad
 \|H_C\|_\infty\le C_h(k+\log k).
\end{equation}
It retains the origin logarithm as $\log(|y|/k+k^{-1})$, whose integral after
$y=kt$ is finite. Its actual complete-relation central fraction satisfies
$\int_{-kb}^{kb}|P|^2\sigma\le\eta_{k,1}\|P\|_\sigma^2$
for every $P\in Q_k\PP_q$, with the exponentially small original
$\eta_{k,1}$ from the complete-root localization proof L22--35.
Thus the trace of the entire relation band on that interval is at most
$2(q+1)\eta_{k,1}$; no orthogonality or minimizer is changed.

Differentiating the original source-minus-relation factorization (SR3) yields
\begin{equation}\tag{SR42}
 F_t(0,0)=\int H_C(y)
 [\mathsf R_{q,1}(y)-\mathsf P_{q,1}(y)]\dd y
 =\frac{\log2}{\pi}I_hk^2+e_{\rm lin}.
\end{equation}
To bound this derivative at finite $k$, let $\beta_{q,k}$ be the explicit
right side of (SR39) evaluated at $s=1$, $y=kb$. Every term of that bound is
nondecreasing in $|y|$, so this also bounds it throughout $[-kb,kb]$; there
is no remaining optimization in this definition. Then
\begin{equation}\tag{SR43}
 |e_{\rm lin}|\le
 k^2I_h\beta_{q,k}
 +(\log2/\pi+\beta_{q,k})C_hk
 +2(q+1)\eta_{k,1}\|H_C\|_\infty=:E_{\rm lin}.
\end{equation}
By (SR40), $E_{\rm lin}=O_h(k^3/q+k^5/q^2+k)+o(1)=O_h(k)$.
Every term is a bound for an original restriction or its full preceding-relation
Gram; this is not an absolute allocation of the native subquotients.

For any fixed space $V$, if $A_V(t,r)=\Gram_V[\sigma e^{tH_C+rH_O}]$,
the compact support of $H_C$ and the source envelopes justify two derivatives
under each finite integral. Since $H_CH_O=0$ pointwise,
\begin{equation}\tag{SR44}
 \partial_t\partial_r\log\det A_V
 =-\Tr(A_V^{-1}A_{V,C}A_V^{-1}A_{V,O}).
\end{equation}
The same formula holds in the fixed full relation space. It is generally
nonzero: pointwise disjoint support does not commute through $A_V^{-1}$.
The second central derivative is the exact projection variance
\begin{equation}\tag{SR45}
 \partial_t^2\log\det A_V
 =\frac12\iint(H_C(x)-H_C(y))^2|K_V(x,y)|^2
 \dd\nu_{t,r}(x)\dd\nu_{t,r}(y),\quad
 \nu_{t,r}=\sigma e^{tH_C+rH_O}.
\end{equation}
To verify it, differentiate $\log\det A_V$ twice, obtaining
$\Tr(A_V^{-1}A_{V,CC})-\Tr((A_V^{-1}A_{V,C})^2)$;
write both traces in an orthonormal basis and use
$\int|K_V(x,y)|^2\dd\nu(y)=K_V(x,x)$. These full matrix identities retain
all cross terms. The original eight-term combination is signed.

Define the actual receiver terms
\[
 \mathcal C_k^{\rm mixed}=\int_0^1F_{tr}(0,r)\dd r,\qquad
 \mathcal N_k^{\rm nonlinear}=\int_0^1(1-t)F_{tt}(t,1)\dd t.
\]
The fundamental theorem and Taylor's integral identity give exactly
\begin{equation}\tag{SR46}
 \mathcal C_k^{\rm mixed}+\mathcal N_k^{\rm nonlinear}
 =F(1,1)-F(0,1)-F_t(0,0).
\end{equation}
Now $F(1,1)-F(0,1)=\BB_k[m_k]-\BB_k[\omega_k]$;
this equality uses the scalar cancellation in (SR2), not a change of endpoint.
Substitute the already proved finite-amplitude theorem (SR5) and the newly
proved derivative (SR42)--(SR43). The equal coefficients cancel once, giving
\begin{equation}\tag{SR47}
 |\mathcal C_k^{\rm mixed}+\mathcal N_k^{\rm nonlinear}|
 \le E_{\rm cen}+E_{\rm lin}=o_h(k^2).
\end{equation}
In fact this is $o_h(q/\log q)$ on each fixed multiplicity stratum:
divide its terms by $q/\log q$ and use $k^2\le q$,
$k\log^3q/q\to0$, and $k^4\log q/q^3\to0$.
Neither summand is separately assigned a sign or an $o(q)$ bound by this
calculation. Equation (SR46) is the exact morphism from the path's two named
remainders to the independently evaluated full central change.

\section{Receiving formulas and scope}
In the complete original action substitute (SR28) exactly once:
\begin{equation}\tag{SR48}
 J_k=\BB_k[\sigma]-dC_\Gamma q+\frac{\log2}{\pi}I_hk^2
 +\frac{C_\Gamma q}{2(\log q+c_\zeta)}
 +W_k^{\rm mono}-P_{k,-}-P_{k,+}+R_{h,k}.
\end{equation}
The full Gamma baseline and both penalties remain. This does not evaluate
their own lower-order coefficients or the seven absolute native allocations.
The order-$k$ comparison is exactly
$\delta_{k,k}^{\rm ar}=\delta_{k,1}^{\rm ar}+\BB_k[\sigma]-\BB_k[w_k]$;
no order conversion is omitted. For every invertible marked frame $\mathcal J$,
$\det(\mathcal J^*G_N[\nu]\mathcal J)=|\det\mathcal J|^2\det G_N[\nu]$,
so the same arithmetic--Gamma relative determinant receives (SR28) after
its original frame factor cancels. At a singular frame parameter one first
uses the proved regularized frame; the cancellation is not asserted for
a zero determinant without that map.

The source sign proved in the companion theta calculation sharpens the
simple-stratum sign to $d_0>7370383/7260000>1$. For each fixed $m_0\ge2$,
$d_0=-(4m_0-2)C_\Gamma<0$ remains. The new $q/\log q$ term is positive in
both cases. The original projected ratios $U,V$, all complementary CPQ
returns, and the seven absolute native quantities retain their definitions;
no numerical value for them follows solely from this scalar determinant.



% END PRESERVED INPUT 01_SLOW_TAIL_AND_NONLINEAR_RECEIVER.tex
\endgroup

\clearpage
\begingroup
\def\dd{\,d}

\setcounter{equation}{0}
\renewcommand{\theequation}{A02TS.\arabic{equation}}
\section*{The actual theta-source sign and physical equilibrium constant}
\addcontentsline{toc}{section}{The actual theta-source sign and physical equilibrium constant}
% BEGIN PRESERVED INPUT 02_THETA_SIGN_INDEPENDENT_DERIVATION.tex SHA256 e1fa0df19832c6cee27dd463a2fdd354de02d2b92d1e26babb6b4f59794102f8



This is a successor proof to the accepted signed arithmetic calculation,
not an alteration of its sealed reader. The incoming sign argument is
rederived below. In particular the equilibrium logarithmic moment is
derived directly from the original EIQ measure, so its closed form is
not an additional assumption. All finite inequalities are rational;
\texttt{check\_theta\_sign\_exact.py} implements precisely their finite sums.
The classical theta and special-function inputs are cited at use.

\section{Original source, original coefficient, and result}
Fix an actual off-line quartet of zeros of $g=2\xi$, with complete common
order $m_0\geq1$ and representatives
\begin{equation}\tag{TS1}
 \rho=\tfrac12+\delta+i\gamma,\qquad 0<\delta<\tfrac12,\quad\gamma>2,
 \qquad
 h(s)=\prod_{z\in\{\rho,\bar\rho,1-\rho,1-\bar\rho\}}(s-z)^{m_0}.
\end{equation}
The word actual means that $g(\rho)=0$ and that these are the stated
zeros of the original numerator. No existence of such a quartet is
asserted here. Define, with its original mass,
\begin{equation}\tag{TS2}
\begin{gathered}
 w_h(y)=\frac{|g(1/2+iy)/h(1/2+iy)|^2}{2\pi},\quad
 M_h(\theta)=\int_{\mathbb R}e^{\theta y}w_h(y)\dd y,\\
 \alpha_*=\frac\pi2,\quad
 I_h=\int_0^{\alpha_*}\left(\frac{M_h'(\theta)}{M_h(\theta)}\right)^2\dd\theta.
\end{gathered}
\end{equation}
The already proved original-source envelope is
$e^{\alpha_*|y|}w_h(y)\leq C_h(1+|y|)^{-p_h}$ with
$p_h=8m_0-9/2>3$; its complete provider is recorded in
\cite{a02ts:SBV}. Thus dominated convergence gives $M_h\in C^2$
on the closed tilt interval, with positive values there. The mass is
$M_h(0)$ and is not replaced by one.

Let $\rho_{a,b}$ denote the original equilibrium measure for
$V_{a,b}(x)=b\sqrt x-a\log x$ on $(0,\infty)$, and put
\begin{equation}\tag{TS3}
 \mathcal L(a,b)=\int\log x\dd\rho_{a,b}(x),\qquad
 \mathcal C_\Gamma=2\mathcal L(2,\pi)-8\log2+4.
\end{equation}
Section~\ref{a02ts:eqphysical} retains its physical elliptic integrals,
constructs this measure, and computes the moment. The symbol
$\mathcal C_\Gamma$ refers to this return coefficient, not a constant
in a pointwise upper bound for the Gamma density.

\begin{theorem}\label{a02ts:th:sign}
For every actual quartet in (TS1), at every complete multiplicity,
\begin{equation}\tag{TS4}
 \frac{M_h(\pi/2)}{M_h(0)}>61>e^{41/10},\qquad
 I_h>\frac{1681}{50\pi},\qquad 0<\mathcal C_\Gamma<\frac23.
\end{equation}
For $m_0=1$ the already evaluated coefficient therefore satisfies
\begin{equation}\tag{TS5}
 d_0(h)=-2\mathcal C_\Gamma+\frac{\log2}{\pi}I_h
 >\kappa_*:=\frac{7370383}{7260000}>1.
\end{equation}
The pressure bound (TS4) applies also when $m_0\geq2$.
At those fixed higher multiplicities the asymptotic coefficient is
$-(4m_0-2)\mathcal C_\Gamma<0$, because the original $k^2/q$ tends to zero.
\end{theorem}

\section{Positive theta kernel with the factor two retained}
For $x>0$ put
\[
 \vartheta(x)=\sum_{n\in\mathbb Z}e^{-\pi n^2x^2},\quad
 D=-x\partial_x,\quad (Jf)(x)=x^{-1}f(1/x).
\]
The theta transformation $\vartheta=J\vartheta$ is the specialization
of the classical modular identity, DLMF 20.7.32 at $z=0$,
$\tau=ix^2$ \cite{a02ts:theta}. Differentiating a term with $a=\pi n^2$
gives $D e^{-ax^2}=2ax^2e^{-ax^2}$ and
$D^2e^{-ax^2}=(-4ax^2+4a^2x^4)e^{-ax^2}$. Therefore
\begin{equation}\tag{TS6}
 T(x)=D(D-1)\vartheta(x)
 =2\sum_{n\geq1}(4\pi^2n^4x^4-6\pi n^2x^2)e^{-\pi n^2x^2}.
\end{equation}
The series and all its derivatives converge uniformly on compact
subsets of $x>0$. Direct differentiation of $Jf$ gives
$DJ=J(1-D)$, hence $D(D-1)J=JD(D-1)$ and $T=JT$.
Consequently $K(t)=e^{t/2}T(e^t)$ is even. For $t\geq0$ every term
of (TS6) is strictly positive, since $2\pi n^2e^{2t}-3>0$.
It follows that $K(t)>0$ for every real $t$. For $x\geq1$, the
Gaussian series in (TS6) is bounded by $C(1+x^4)e^{-\pi x^2}$;
the same type of bound holds after any fixed number of logarithmic
derivatives. Reflection gives decay faster than every exponential
at the other end. This proves all Mellin and Fourier integrability
statements used next.

For $\Re s>1$ each Mellin integral is absolutely integrable and its
sum is absolutely convergent. Substitution $u=\pi n^2x^2$ gives
\[
 \int_0^\infty(4\pi^2n^4x^4-6\pi n^2x^2)e^{-\pi n^2x^2}x^s\frac{dx}{x}
 =\frac12\pi^{-s/2}n^{-s}s(s-1)\Gamma(s/2).
\]
Here $4\Gamma(s/2+2)-6\Gamma(s/2+1)=s(s-1)\Gamma(s/2)$.
The literal factor two in (TS6) yields
\begin{equation}\tag{TS7}
 \int_0^\infty T(x)x^s\frac{dx}{x}
 =s(s-1)\pi^{-s/2}\Gamma(s/2)\zeta(s)=2\xi(s)=g(s).
\end{equation}
The convention for $\xi$ is DLMF 25.4.4 \cite{a02ts:zeta}.
Rapid decrease at both ends extends the left side to an entire
function; the identity continues to every $s$.

Set $G(z)=g(1/2+z)$. Substitution $x=e^t$ in (TS7), then evenness
of $K$, gives the absolutely convergent representation
\begin{equation}\tag{TS8}
 G(z)=2\int_0^\infty K(t)\cosh(zt)\dd t
 =\sum_{n=0}^\infty a_nz^{2n},\qquad
 a_n=\frac2{(2n)!}\int_0^\infty t^{2n}K(t)\dd t>0.
\end{equation}
Convergence is uniform on compact $z$-sets, by the exponential
majorants already proved. Thus these are the Taylor coefficients
of the original numerator.

\section{A fixed interval determined by actual zeros}
The residue-one pole of $\zeta$ at one, the Gamma recurrence, and
$\zeta(2)=\pi^2/6$, $\zeta(6)=\pi^6/945$ give
\begin{equation}\tag{TS9}
 G(1/2)=1,\qquad G(3/2)=\pi/3,\qquad G(11/2)=4\pi^3/63.
\end{equation}
These classical values and conventions are recorded in DLMF
25.2(i), 25.4.4 and 25.6.1 \cite{a02ts:zeta,a02ts:zetavalues}.
Since $(3/2)^{2n}\geq9(1/2)^{2n}$ for $n\geq1$, (TS8) gives
\[
 G(3/2)-a_0\geq9(G(1/2)-a_0).
\]
Subtracting the series at $3/2$ and $1/2$ also gives
$G(3/2)-G(1/2)\geq2a_1$. Thus
\begin{equation}\tag{TS10}
 a_0\geq\frac{27-\pi}{24}>\frac{167}{168},\qquad
 \frac{a_1}{a_0}\leq\frac{4(\pi-3)}{27-\pi}<\frac4{167}.
\end{equation}
The rational estimates use $\pi<22/7$, proved by a finite series
below. Positivity of all coefficients gives, for $|z|\leq11/2$,
\begin{equation}\tag{TS11}
 |G(z)|\geq a_0-|G(z)-a_0|
 \geq2a_0-G(11/2)
 >\frac{167}{84}-\frac4{63}\left(\frac{22}{7}\right)^3
 =\frac{1475}{86436}>0.
\end{equation}
Apply (TS11) to the actual zero $z=\delta+i\gamma$. It implies
$\delta^2+\gamma^2>121/4$ and hence
\begin{equation}\tag{TS12}
 \gamma^2-\delta^2>121/4-2(1/2)^2=119/4.
\end{equation}
No numerical zero table or extra hypothesis on the height is used.

The elementary inequalities $|\cos u|\leq1$ and
$\cos u\geq1-u^2/2$ in (TS8) give
$|G(iy)|\leq a_0$ and $G(iy)\geq a_0-a_1y^2$.
Put $T_*=27/5$ and $c_*=4/167$. The exact identities
\[
 119/4-T_*^2=59/100>0,\qquad
 1-c_*T_*^2=1259/4175>0
\]
show that squaring is legitimate on $0\leq y\leq T_*$:
\begin{equation}\tag{TS13}
 a_0^2(1-c_*y^2)^2\leq G(iy)^2\leq a_0^2.
\end{equation}

\section{The complete denominator and the entire outside mass}
Multiplying the four original linear factors on the source line gives
\begin{equation}\tag{TS14}\begin{aligned}
 H(y)&=[(y-\gamma)^2+\delta^2][(y+\gamma)^2+\delta^2]\\
 &=(y^2-\gamma^2+\delta^2)^2+4\delta^2\gamma^2>0,\\
 h(1/2+iy)&=H(y)^{m_0},\qquad
 w_h(y)=\frac{G(iy)^2}{2\pi H(y)^{2m_0}}.
\end{aligned}\end{equation}
In particular the full power and all four roots remain present.
Since $H'(y)=4y[y^2-(\gamma^2-\delta^2)]<0$ on $(0,T_*]$,
$R_h(y)=H(y)^{-2m_0}$ increases there.

For a finite positive measure $\nu$ and real integrable functions
$f,R$, direct expansion gives
\begin{equation}\tag{TS15}
 \left(\int fR\dd\nu\right)\left(\int1\dd\nu\right)
 -\left(\int f\dd\nu\right)\left(\int R\dd\nu\right)
 =\frac12\iint(f(y)-f(z))(R(y)-R(z))\dd\nu(y)\dd\nu(z).
\end{equation}
On $[0,T_*]$, use $d\nu=G(iy)^2dy$, $f(y)=\cosh(\alpha_*y)$
and $R=R_h$. These functions are bounded and the measure is finite,
so Fubini applies. The right side is nonnegative because both
$f$ and $R$ increase. All denominators below are positive by (TS13).
Using (TS13) only after this exact comparison yields
\begin{equation}\tag{TS16}
 \frac{\int_0^{T_*}\cosh(\alpha_*y)w_h(y)\dd y}
      {\int_0^{T_*}w_h(y)\dd y}
 \geq\frac{\int_0^{T_*}\cosh(\alpha_*y)G(iy)^2\dd y}
              {\int_0^{T_*}G(iy)^2\dd y}
 \geq J_*:=\frac1{T_*}\int_0^{T_*}\cosh(\alpha_*y)(1-c_*y^2)^2\dd y.
\end{equation}
On this interval $0<(1-c_*y^2)^2\leq1$, hence
$J_*\leq\cosh(\alpha_*T_*)$. On the entire remaining half-line
$y\geq T_*$, $\cosh(\alpha_*y)\geq\cosh(\alpha_*T_*)\geq J_*$.
Therefore
$\int_0^\infty(\cosh(\alpha_*y)-J_*)w_h(y)dy\geq0$:
the inner integral is nonnegative by (TS16) and the outer one
is nonnegative pointwise. Both integrals are finite by the retained
envelope. The exact evenness in (TS14) then gives
\begin{equation}\tag{TS17}
 \frac{M_h(\alpha_*)}{M_h(0)}
 =\frac{\int_0^\infty\cosh(\alpha_*y)w_h(y)\dd y}
        {\int_0^\infty w_h(y)\dd y}\geq J_*.
\end{equation}
This argument places no monotonicity requirement on $R_h$ outside
$[0,T_*]$ and removes no outer mass.

\section{Finite rational pressure estimate}
Because $\alpha_*>157/100$, the positive Taylor series of $\cosh$
and exact polynomial integration give
\begin{equation}\tag{TS18}\begin{aligned}
 J_*&\geq S_*:=\sum_{j=0}^8\frac{(157/100)^{2j}}{(2j)!}
 \left[\frac{T_*^{2j}}{2j+1}-\frac{2c_*T_*^{2j+2}}{2j+3}
                +\frac{c_*^2T_*^{2j+4}}{2j+5}\right],\\
 \frac{611585}{10000}&<S_*<\frac{611586}{10000},\qquad S_*>61.
\end{aligned}\end{equation}
Every bracket is exactly
$T_*^{-1}\int_0^{T_*}y^{2j}(1-c_*y^2)^2dy>0$.
For $x=41/10$, all terms after $x^{21}/21!$ have successive
ratios at most $x/22<1$, so
\begin{equation}\tag{TS19}
 e^{41/10}\leq\sum_{j=0}^{20}\frac{(41/10)^j}{j!}
 +\frac{(41/10)^{21}}{21!}\frac1{1-(41/10)/22}<61.
\end{equation}
Both (TS18) and (TS19) are finite rational inequalities verified
by the delivered exact cross-products. They imply
$\log(M_h(\alpha_*)/M_h(0))>41/10$. The positive $C^1$ function
$M_h$ satisfies the fundamental theorem of calculus on the whole
closed interval. Cauchy--Schwarz thus proves
\begin{equation}\tag{TS20}
 I_h\geq\frac1{\alpha_*}
 \left(\int_0^{\alpha_*}\frac{M_h'(\theta)}{M_h(\theta)}\dd\theta\right)^2
 =\frac2\pi\log^2\frac{M_h(\pi/2)}{M_h(0)}
 >\frac{1681}{50\pi}.
\end{equation}

\section{Physical EIQ measure and an independent exact moment derivation}
\label{a02ts:eqphysical}
Use $a,b>0$ for the potential parameters in this section, to keep them
distinct from $\alpha_*=\pi/2$ and all original packet parameters.
Define the physical complete elliptic integrals
\[
 K(\kappa)=\int_0^{\pi/2}\frac{d\theta}{\sqrt{1-\kappa^2\sin^2\theta}},
 \qquad E(\kappa)=\int_0^{\pi/2}\sqrt{1-\kappa^2\sin^2\theta}\dd\theta.
\]
The conventions agree with Carlson's DLMF Chapter 19 \cite{a02ts:elliptic}.
For $r\in(0,1)$, put $\kappa=\sqrt{1-r^2}$. Differentiating these
integrals gives $dE/dr=r(K-E)/(1-r^2)$ and
$dK/dr=(r^2K-E)/(r(1-r^2))$; the latter follows also by integrating
the derivative of $\sin\theta\cos\theta/
\sqrt{1-\kappa^2\sin^2\theta}$ between its zero boundary values.
It follows that
\begin{equation}\tag{TS21}
 \frac d{dr}\frac{rK}{E}
 =\frac{(K-E)(E-r^2K)}{(1-r^2)E^2}>0.
\end{equation}
The first numerator factor is the integral of
$(1-r^2)\sin^2\theta/\sqrt{1-(1-r^2)\sin^2\theta}$;
the second is the corresponding integral with $\cos^2\theta$.
Both are positive. Dominated convergence gives $rK\to0$, $E\to1$
as $r\downarrow0$, since the integrand of $rK$ is bounded by one.
At $r\uparrow1$, $K,E\to\pi/2$. Thus there is a unique solution of
\begin{equation}\tag{TS22}
 \frac{rK}{E}=\frac a{a+2},\qquad
 v=\frac{(a+2)\pi}{bE},\quad u=rv,
 \quad uK=\frac{a\pi}{b},\quad vE=\frac{(a+2)\pi}{b}.
\end{equation}

Put $A=u^2$, $B=v^2$, $m=(A+B)/2$, $R_s=\sqrt{(A+s)(B+s)}$,
$R_0=uv$, and $V(x)=b\sqrt x-a\log x$.
The substitutions $x=A\cos^2\theta+B\sin^2\theta$ give
\[
 \int_A^B\frac{V'(x)}{\sqrt{(B-x)(x-A)}}\dd x=0,\qquad
 \int_A^B\frac{xV'(x)}{\sqrt{(B-x)(x-A)}}\dd x=2\pi.
\]
Indeed the integrals with numerators $x^{-1/2},x^{-1},x^{1/2}$
are respectively $2K/v$, $\pi/(uv)$, and $2vE$.
Since $\int_0^\infty s^{-1/2}(x+s)^{-1}ds=\pi/\sqrt x$,
Tonelli in the first identity gives
\[a/R_0=(b/(2\pi))\int_0^\infty ds/(\sqrt sR_s)\].
Define the positive function and measure
\begin{equation}\tag{TS23}\begin{aligned}
 H_*(x)&=\frac{a}{xR_0}-\frac b{2\pi}\int_0^\infty
           \frac{ds}{\sqrt s(x+s)R_s}
 =\frac b{2\pi x}\int_0^\infty\frac{\sqrt s}{(x+s)R_s}\dd s>0,\\
 d\rho_{a,b}(x)&=\frac{\sqrt{(B-x)(x-A)}}{2\pi}H_*(x)
                 \mathbf1_{(A,B)}(x)\dd x.
\end{aligned}\end{equation}
All auxiliary integrals converge at zero and infinity. This is exactly
the EIQ density \cite{a02ts:EIQ}, with no coordinate rescaling.

For completeness its mass and Euler equality can be checked directly.
Let $R(z)=\sqrt{(z-A)(z-B)}$ with $R(z)/z\to1$ at infinity.
Substitution $x=m+(B-A)\cos\theta/2$ and then
$t=\tan(\theta/2)$ gives the arcsine transform $1/R(z)$.
Polynomial division of $(B-x)(x-A)$ then gives
\begin{equation}\tag{TS24}
 \frac1\pi\int_A^B\frac{\sqrt{(B-x)(x-A)}}{(x+s)(z-x)}\dd x
 =1-\frac{R_s+R(z)}{z+s}.
\end{equation}
The removable value at $z=-s$ is understood. The same division gives
\[\pi^{-1}\int_A^B\sqrt{(B-x)(x-A)}(x+s)^{-1}dx=m+s-R_s\].
Using (TS23) and $\int_0^\infty s^{-1/2}(1-s/R_s)ds=2vE$,
which follows by averaging the Stieltjes representation of $\sqrt x$
against the arcsine measure, yields
\[
 \int d\rho_{a,b}=-a/2+bvE/(2\pi)=1.
\]
Furthermore (TS24), the Stieltjes identity, and the exact cancellation
\[a/R_0=(b/(2\pi))\int_0^\infty ds/(\sqrt sR_s)\] give
\begin{equation}\tag{TS25}
 \int\frac{d\rho_{a,b}(x)}{z-x}
 =\frac{V'(z)-R(z)H_*(z)}2.
\end{equation}
These identities first hold off both $[A,B]$ and the negative real
axis, where all integrations are absolutely convergent. Taking real
boundary values on $(A,B)$ gives the principal-value derivative
of the logarithmic potential, because the density is locally $C^1$;
subtracting its value at the singular point proves this limit.
As $R(x+i0)$ is purely imaginary there, (TS25) implies
\begin{equation}\tag{TS26}
 V(x)-2\int\log|x-t|\dd\rho_{a,b}(t)=\ell\qquad(A\leq x\leq B).
\end{equation}
The endpoints follow by continuity; bounded compactly supported
density makes the logarithmic singularity uniformly integrable.
Outside $[A,B]$ its derivative is $R(x)H_*(x)$, negative on $(0,A)$
and positive on $(B,\infty)$. This proves the strict exterior Euler
inequality. It identifies the constructed measure with the unique
EIQ equilibrium: for the difference $\nu$ from any competing finite
energy probability measure, expansion of energy gives
$\int(V-2U_\rho-\ell)\dd\nu-\iint\log|x-y|\dd\nu(x)\dd\nu(y)$.
The first term equals its integral against the competing probability
and is nonnegative. The second is nonnegative, since for a zero-mass
signed measure the identity
\[
 -\iint\log|x-y|\dd\nu(x)\dd\nu(y)
 =\frac12\int_0^\infty\frac1t
       \iint e^{-t(x-y)^2}\dd\nu(x)\dd\nu(y)\dd t
\]
follows from the scalar identity for $-\log d$ by bounded truncation
and dominated convergence. The inner double integral is
$(4\pi t)^{-1/2}\int e^{-\zeta^2/(4t)}|\widehat\nu(\zeta)|^2d\zeta$.
It is positive unless $\nu=0$: a zero Fourier transform forces every
Gaussian convolution, and hence the measure, to vanish. The
logarithmic integrability needed here follows from finite energy and
the coercivity of $V$; cross terms with $\rho$ are integrable by its
bounded compact density. Infinite-energy competitors cannot improve
the finite energy of (TS23).

We now derive the closed logarithmic moment without importing ECL14.
Write
\[
 C=\frac{(u+v)^2}{4},\quad\eta=\frac{v-u}{v+u},\quad
 d\nu_0(x)=\frac{dx}{\pi\sqrt{(B-x)(x-A)}},\quad
 d\omega_0(x)=\frac{R_0}{x}\dd\nu_0(x).
\]
Both measures are probabilities. Under $x=m+(B-A)\cos\theta/2$
the first is $d\theta/\pi$ and
$x=C(1+2\eta\cos\theta+\eta^2)$,
$R_0=C(1-\eta^2)$. Therefore the second measure has the Poisson
kernel density $(1-\eta^2)/(1+2\eta\cos\theta+\eta^2)$ and
\begin{equation}\tag{TS27}
 \int\cos(j\theta)\dd\omega_0=(-\eta)^j\quad(j\geq0).
\end{equation}
This follows by expanding the two geometric series and integrating;
all series converge absolutely for $\eta<1$.
The absolutely convergent expansion of $\log x$ gives
\begin{equation}\tag{TS28}
 \int\log x\dd\nu_0=\log C,\qquad
 \int\log x\dd\omega_0=\log C+2\log(1-\eta^2)
 =2\log R_0-\log C.
\end{equation}
For $t=m+(B-A)\cos\phi/2\in[A,B]$, the logarithmic distance has
the Abel-convergent cosine expansion
\[
 \log|x-t|=\log\frac{B-A}{4}
       -2\sum_{j\geq1}\frac{\cos(j\theta)\cos(j\phi)}j.
\]
To verify it, factor $\cos\theta-\cos\phi$ into the two sine
factors and take real parts of $-\sum z^j/j=\log(1-z)$.
Inserting an Abel factor $s^j$, $0<s<1$, permits termwise
integration. As $s\uparrow1$, convergence holds in $L^1(d\theta)$
by the locally integrable logarithmic singularity; the Poisson
kernel in (TS27) is bounded, so the same is true against $d\omega_0$.
The endpoint values follow by this identical one-sided integrability.
Thus (TS27)--(TS28) give the exact potential identities
\begin{equation}\tag{TS29}
 \int\log|x-t|\dd\nu_0(x)=\log\frac{B-A}{4},\qquad
 \int\log|x-t|\dd\omega_0(x)=\log t+\log\eta.
\end{equation}
Finally the endpoint integrals in (TS22) give
\[
 \int\sqrt x\dd\nu_0=2vE/\pi=2(a+2)/b,\qquad
 \int\sqrt x\dd\omega_0=2uK/\pi=2a/b.
\]
Integrate (TS26) first against $\nu_0$, then against $\omega_0$.
Fubini applies because both auxiliary densities have integrable
square-root endpoints and the density of $\rho$ is bounded;
the logarithmic singularities are absolutely integrable. This yields
\begin{align*}
 \ell&=2(a+2)-a\log C-2\log((B-A)/4),\\
 2\mathcal L(a,b)+2\log\eta
 &=2a-a(2\log R_0-\log C)-\ell.
\end{align*}
Since $(B-A)/(4\eta)=C$, substitution proves the exact formula
\begin{equation}\tag{TS30}
 \boxed{\mathcal L(a,b)=(a+1)\log C-a\log R_0-2.}
\end{equation}
In particular, at $(a,b)=(2,\pi)$,
\begin{equation}\tag{TS31}
 uK=2,\quad vE=4,\quad 2rK=E,\qquad
 \mathcal L(2,\pi)=6\log((u+v)/2)-2\log(uv)-2.
\end{equation}
All these variables refer to the original physical potential.
Inserting $u=rv$ and $v=4/E$ in (TS3) gives
\begin{equation}\tag{TS32}
 \boxed{\mathcal C_\Gamma=4\log H_*,\qquad
 H_*:=\frac{(1+r)^3}{8rE(\sqrt{1-r^2})}.}
\end{equation}
In particular the denominator contains $E$, not $2E/\pi$.

\section{Exact finite bounds for the physical elliptic constant}
For rational $x\in(0,1)$ set $m_x=1-x^2$ and
\[
 t_n(x)=\left(\frac{\binom{2n}{n}}{4^n}\right)^2m_x^n,
 \quad K_N(x)=\sum_{n=0}^Nt_n(x),\quad
 E_N(x)=1-\sum_{n=1}^N\frac{t_n(x)}{2n-1}.
\]
The binomial series for $(1-z)^{-1/2}$, integrated against
$z=m_x\sin^2\theta$, gives $2K/\pi=\sum_{n\geq0}t_n$.
Here $\int_0^{\pi/2}\sin^{2n}\theta\,d\theta=
(\pi/2)\binom{2n}{n}/4^n$, obtained by integration by parts from
$n=0$. The binomial series for $(1-z)^{1/2}$ similarly gives
$2E/\pi=1-\sum_{n\geq1}t_n/(2n-1)$.
Since $t_{n+1}/t_n=m_x(2n+1)^2/(2n+2)^2<m_x$, geometric
majoration proves both finite enclosures
\begin{equation}\tag{TS33}
 K_N\leq\frac{2K}{\pi}\leq K_N+\frac{t_{N+1}}{1-m_x},\qquad
 E_N-\frac{t_{N+1}}{(2N+1)(1-m_x)}
 \leq\frac{2E}{\pi}\leq E_N.
\end{equation}
Take $N=512$, $r_-=1601/10000$, $r_+=801/5000$.
With each $K,E$ evaluated at $\sqrt{1-r_\pm^2}$, exact rational
substitution in (TS33) proves
\begin{equation}\tag{TS34}
 -\frac{443}{10^8}<2r_-\frac{2K_-}{\pi}-\frac{2E_-}{\pi}
 <-\frac{441}{10^8}<0,
 \quad
 0<\frac{2597}{10^7}<2r_+\frac{2K_+}{\pi}-\frac{2E_+}{\pi}
 <\frac{2598}{10^7}.
\end{equation}
The strict monotonicity (TS21) implies $r_-<r<r_+$.
The defining integrand shows that $E(\sqrt{1-r^2})$ increases
with $r$. Define $\underline E_-=E_{512}(r_-)-
t_{513}(r_-)/(1025(1-m_{r_-}))$. Using the physical $\pi/2$
factor and $157/50<\pi<22/7$, the following rational bounds hold:
\begin{equation}\tag{TS35}
 \frac{1176}{1000}
 <\frac{(1+r_-)^3}{8r_+(11/7)E_{512}(r_+)}
 <H_*<\frac{(1+r_+)^3}{8r_-(157/100)\underline E_-}
 <\frac{1179}{1000}<\frac{85}{72}<e^{1/6}.
\end{equation}
The penultimate inequality is rational. The last follows by retaining
the first three positive terms $1+1/6+1/72$ in the exponential series.
The left endpoint exceeds one. Equations (TS32) and (TS35) prove
$0<\mathcal C_\Gamma<2/3$.

For completeness all elementary constants used above have rational
certificates as well. The tangent addition formulas yield
$\tan(2\arctan(1/5))=5/12$,
$\tan(4\arctan(1/5))=120/119$, and
$\tan(4\arctan(1/5)-\arctan(1/239))=1$.
The final angle is positive, by
$\arctan x\geq x/(1+x^2)$ and $\arctan x\leq x$;
it is less than $4/5<1<\pi/2$, where
$\pi/2=2\arctan1>1$ follows by integrating $(1+t^2)^{-1}>1/2$
on $[0,1)$. Hence Machin's identity is
$\pi=16\arctan(1/5)-4\arctan(1/239)$.
The alternating series through eight terms, with its first omitted
term, encloses $\pi$ strictly between $157/50$ and $22/7$.
Also integration of the geometric series on $[0,1/3]$ gives
$\log2=2\sum_{j\geq0}(1/3)^{2j+1}/(2j+1)>56/81>69/100$.
All stated finite comparisons are made with exact integer
cross-multiplication by the delivered certificate; its 33 comparisons
pass identically with and without Python optimization.

\section{Coefficient and finite receivers on the original domain}
Combining (TS20), $\mathcal C_\Gamma<2/3$, $\log2>69/100$ and
$\pi<22/7$ gives
\begin{equation}\tag{TS36}
 -2\mathcal C_\Gamma+\frac{\log2}{\pi}I_h
 >-\frac43+\frac{1681}{50}\frac{69}{100}\frac{49}{484}
 =\frac{7370383}{7260000}>1.
\end{equation}
This completes the source and equilibrium part of Theorem~\ref{a02ts:th:sign}.

The scalar receiving map is the already accepted original one
\cite{a02ts:SBV}. To make its scope explicit, retain
\begin{equation}\tag{TS37}\begin{aligned}
 k&=4\ell+1\geq5,\quad e_k=1+k(m_0-1),\quad q=e_k(k+1)^2,\\
 \chi_k(S)&=\prod_{a,b=0}^k
 [S-k/2-(2a-k)\delta-i(2b-k)\gamma]^{e_k},\\
 G_N[\nu]&=(J_NH_N[\nu]^{-1}J_N^*)^{-1},\qquad J_N=\operatorname{rem}_{\chi_k},\\
 \mathcal B_k[\nu]&=\log\det G_{q-1}[\nu]+\log\det G_q[\nu]
 -\log\det G_{2q-1}[\nu]-\log\det G_{2q}[\nu],\\
 d\sigma(y)&=\frac{|\Gamma(1/4+iy/2)|^2}{2\pi}\dd y,\qquad
 \delta_{k,1}^{\rm ar}=\mathcal B_k[w_h^{*k}(y)dy]-\mathcal B_k[d\sigma].
\end{aligned}\end{equation}
Here $H_N$ is the original degree-$N$ polynomial Gram matrix in
the physical coordinate $S=k/2+iy$. The coordinate morphism is
$P(S)\mapsto P(k/2+iy)$, with inverse substitution
$y=(S-k/2)/i$; no roots or lower components are removed.
The Gamma mass is $\sqrt{2\pi}$ and the arithmetic mass is
$M_h(0)^k$. The accepted theorem states, on its fixed-packet
sufficiently-large domain $k\geq k_0(h)$,
\begin{equation}\tag{TS38}\begin{aligned}
 \delta_{k,1}^{\rm ar}
 &=-(4m_0-2)\mathcal C_\Gamma q+\frac{\log2}{\pi}I_h k^2+\epsilon_{h,k},\\
 |\epsilon_{h,k}|&\leq E_{h,k}=C_h^{\rm SBV}q\mathfrak r(q),\qquad
 \mathfrak r(q)=\frac{\log\log(q+16)}{\log(q+2)}
                       +q^{-1/200}\log^2(q+2)\longrightarrow0.
\end{aligned}\end{equation}
Its full accepted proof and providers are preserved separately.
The present argument establishes the missing sign by applying
that theorem, and does not reassign its proof as unfinished work.

For $m_0=1$, define $A_*=5683461/2420000$ and note
$A_*-4/3=\kappa_*$. Then $q=(k+1)^2$ and the exact finite lower
receiver is
\begin{equation}\tag{TS39}
 \delta_{k,1}^{\rm ar}
 >A_*k^2-\frac43q-E_{h,k}
 =\kappa_*q-A_*(2k+1)-E_{h,k}.
\end{equation}
Thus $\delta_{k,1}^{\rm ar}>q$ whenever $k\geq k_0(h)$ and
\begin{equation}\tag{TS40}
 C_h^{\rm SBV}\mathfrak r(q)+A_*\frac{2k+1}{q}<\kappa_*-1.
\end{equation}
This guard holds eventually for each fixed actual simple packet;
no packet-independent value for $k_0(h)$ or $C_h^{\rm SBV}$ is claimed.
Dividing (TS38) by $q$ gives the positive simple coefficient in (TS5).
At fixed $m_0\geq2$, $k^2/q\to0$, so it instead gives
$\delta_{k,1}^{\rm ar}/q\to-(4m_0-2)\mathcal C_\Gamma<0$.
The same pressure estimate remains true, but its $k^2$ term is
lower order on this original higher-multiplicity scale.

The exact four-endpoint receiving identity for any invertible marked
frame $\mathcal J$ is
$\det(\mathcal J^*G_N[\nu]\mathcal J)=|\det\mathcal J|^2\det G_N[\nu]$.
It follows by multiplicativity of determinants. If the same frame is
used at the four endpoints, its factor cancels with coefficients
$1+1-1-1=0$. The arithmetic-to-Gamma ratio of these complete
four-endpoint products consequently has logarithm
$\delta_{k,1}^{\rm ar}$ exactly; it exceeds $e^q$ under (TS40).
This proves the full-frame scalar transport without assigning a
sign to any proper projected subframe or individual pairing.

Finally the complete action receiver retains its existing identity
\begin{equation}\tag{TS41}
 J_k^{\rm action}
 =\mathcal B_k[d\sigma]+\delta_{k,1}^{\rm ar}
   +W_k^{\rm mono}-P_{k,-}-P_{k,+}.
\end{equation}
Substitution of (TS39) increases this receiver by the displayed
positive scalar correction on the simple stratum, while retaining
the full Gamma return, monic window, and both penalties. The
order-$k$ identity
$\delta_{k,k}^{\rm ar}=\delta_{k,1}^{\rm ar}
 -(\mathcal B_k^{(k)}-\mathcal B_k^{(1)})$ remains unchanged.
The seven absolute allocations and the projected responses $U,V$
are not evaluated by these scalar identities. Nothing here closes
the RH programme or infers a sign for the complete action after
omitting its other terms.



% END PRESERVED INPUT 02_THETA_SIGN_INDEPENDENT_DERIVATION.tex
\endgroup

\clearpage
\begingroup


\setcounter{equation}{0}
\renewcommand{\theequation}{A02CRP.\arabic{equation}}
\section*{Original conductor pivots and all signed directional flags}
\addcontentsline{toc}{section}{Original conductor pivots and all signed directional flags}
% BEGIN PRESERVED INPUT 03_CORRELATED_PIVOT_AND_DIRECTION_PROOF.tex SHA256 afd268ff1a72a90130df3d5605290081f3b1e11ecc086e9ed37ba5d32fb8085c
% Successor calculation. Earlier sealed editions remain unchanged.
\begingroup
\section{Original conductor pivots and correlated directional receivers}
\label{a02crp:sec:CRP}
\subsection{Objects, measures, and the attained fibre}
This calculation integrates the first corrected-row note, equations
(1)--(25), in the supplied transcript \texttt{Untitled 1775.md}, lines
1--851. The coefficient estimates are the fully proved LT1--22/LRC6--9
of \path{LRC_TRANSPORT_COMPLETE.tex}; the original affine fibre is
ATA1--8 of \path{ATTAINED_FIBRE_AND_PROFILE_PROPAGATION.tex}.
All programme and source attributions in those providers are retained.
Here every new step connecting those estimates to individual attained
pivots and their correlated directional sums is proved explicitly.

Fix the actual simple quartet, period, logarithm branch and conductor.
Write
\[
\begin{gathered}
a\equiv1\pmod4,\quad a\ge257,\quad 0<\delta<\tfrac12,\quad\gamma>2,\\
q=(a+1)^2,\quad Q=(a-7)^2,\quad\Delta=16a-48,\\
\beta_{b,u,w}=b/2+(2u-b)\delta+i(2w-b)\gamma,
\quad\chi_b(S)=\prod_{u,w=0}^{b}(S-\beta_{b,u,w}),\\
E_A(z)=\sum_{u,w=0}^{8}a_{uw}e^{\beta_{8,u,w}z},\quad
\mu_j=E_A^{(j)}(0),\quad v=\operatorname{ord}_0 E_A\le80,
\quad\mu_v\ne0,\\
s\in\{1,a\},\quad c=a/2,\quad c'=a/2-4,\quad
m=\Delta-v,\quad L=q-v,\\
0\le r\le M\le q+32a,\quad N=q+M,\quad D=N-v,\quad d=D+1.
\end{gathered}\tag{CRP1}
\]
The finite analytic root guards of LRC13 are imposed on both actual
grids, with the same original order $s$: for $b=a,a-8$,
\[
n_b=(b+1)^2+(s-1)/4\ge100,\qquad
\max\{b\sqrt{\delta^2+\gamma^2},(s-1)/2+1/2\}\le2^{-17}n_b.
\tag{CRP2}
\]
These are the domain of the supplied finite single-norm estimates;
the exact matrix arguments below require no additional root guard.
The original Hilbert space is
\[
\|P\|_{s,e}^2=\int_{\mathbb R}|P(e+iy)|^2dm_s(y),\qquad
dm_s(y)=\frac{M_s2^{s/2}}{4\pi\Gamma(s/2)}
 |\Gamma(s/4+iy/2)|^2dy,\quad M_s=(2\pi)^{s/2}.\tag{CRP3}
\]
No half-line, mass or centre has been changed. The original orthonormal
polynomials, including their phases, are
\[
\begin{gathered}
p_{0,s}=1,\ p_{1,s}=y,\quad
p_{n+1,s}=yp_{n,s}-n(n-1+s/2)p_{n-1,s},\\
h_{n,s}=M_sn!(s/2)_n,\quad
\phi_{n,s,e}(S)=p_{n,s}((S-e)/i)/\sqrt{h_{n,s}},\quad
\rho_n=\sqrt{n!/(s/2)_n},\\
\sum_{n\ge0}\frac{p_{n,s}(y)}{n!}t^n
 =(1+t^2)^{-s/4}e^{y\arctan t}.\tag{CRP4}
\end{gathered}
\]
The generating and orthogonality formulas are the original
Meixner--Pollaczek dictionary in NIST DLMF (18.23.7), (18.19.8)--(18.19.9)
\cite{human:dlmf18}, with $p_n(y)=n!P_n^{(s/4)}(y/2;\pi/2)$.
LT2--3 proves the factors in this dictionary, including $dy=2dx$ and
the original mass $M_s$.

The literal conductor is $(\mathcal T_AP)(S')=
\sum a_{uw}P(S'+\beta_{8,u,w})$. Since
$\mathcal T_AS^n=\sum_{j=v}^n\binom nj\mu_j S'^{n-j}$, its full
kernel on $\mathbb C[S]_{\le N}$ is $\mathbb C[S]_{<v}$ and it is
onto $\mathbb C[S']_{\le D}$. In the original quotient frame
$\phi_{v,s,c},\ldots,\phi_{N,s,c}$ and target frame
$\phi_{0,s,c'},\ldots,\phi_{D,s,c'}$, its square matrix is
\[
\begin{gathered}
A=A_{D,s},\qquad A_{ij}=\rho_i^{-1}g_{j-i}\rho_{j+v}\quad(i\le j),
\quad A_{ij}=0\quad(i>j),\\
w(t)=-i\arctan t,\quad g(t)=t^{-v}e^{-4w(t)}E_A(w(t)),
\quad g_0=(-i)^v\mu_v/v!.\tag{CRP5}
\end{gathered}
\]
This is LT4--6, obtained by applying the conductor to CRP4 and comparing
coefficients. The phase of $g_0$, all zeros of $g$, and all actual
coefficients remain in this matrix.

Let $P_r^+$ minimize $\|\chi_aP\|_{s,c}^2$ over monic degree-$r$
polynomials, and put $\nu_r^+=\|\chi_aP_r^+\|_{s,c}^2$.
Existence, uniqueness and orthogonality follow by completing the
positive coefficient square: the minimizer is orthogonal to
$\chi_a\mathbb C[S]_{<r}$. Consequently the original coordinate
columns
\[
\begin{gathered}
\Phi_r=\chi_aP_r^+/\sqrt{\nu_r^+},\quad
\Phi=(\Phi_0,\ldots,\Phi_M),\quad\Phi^*\Phi=I_{M+1},\\
\quad E=[0_{d\times v},I_d],\quad B=E\Phi,
\quad K=B^*B,\quad H=B^*A^*AB
\end{gathered}\tag{CRP6}
\]
are well defined in the upper native degree-$N$ frame. The matrix $E$
is precisely the orthogonal removal of the complete conductor kernel.
The map $B$ is injective: any element of $\chi_a\mathbb C[S]_{\le M}$
in $\mathbb C[S]_{<v}$ is zero, because $q>v$. Thus $0<K\preceq I$.

Write $u_r=\mathcal T_A(\chi_aS^r)/\ell_r$ with
$\ell_r=\mu_v\binom{q+r}{v}$, and retain the complete minima
\[
\eta_r=\min_{z\in\mathbb C^r}\|u_r+\sum_{t<r}z_tu_t\|_{s,c'}^2,
\quad\nu^-_{m+r}=\min_{P\ {\rm monic},\deg P=m+r}
 \|\chi_{a-8}P\|_{s,c'}^2,
\quad\gamma_r=\log(\eta_r/\nu^-_{m+r})\ge0.\tag{CRP7}
\]
The original relation polynomials lie in this lower ideal. For
$P=S^r+\sum_{t<r}p_tS^t$, the exact mutually inverse coefficient maps are
\[
\frac{\mathcal T_A(\chi_aP)}{\ell_r}
 =u_r+\sum_{t<r}p_t\frac{\ell_t}{\ell_r}u_t,
\qquad p_t=z_t\frac{\ell_r}{\ell_t}.\tag{CRP8}
\]
The earlier monic minimizers form a triangular basis with diagonal one.
Their conductor images therefore span exactly the preceding relation
space, without dropping any lift. If $\delta_r$ is the degree-order
Cholesky pivot of $H$, its variational definition and CRP8 give
\[
\boxed{\delta_r=\frac{|\ell_r|^2\eta_r}{\nu_r^+},\qquad
\gamma_r=\log\delta_r+\log\nu_r^+-\log\nu^-_{m+r}
                         -2\log|\ell_r|.}\tag{CRP9}
\]
These are full attained pivots. The positive-square minimum is the
classical Schur operation of Boyd and Vandenberghe
\cite[Appendix A.5.5, p.~650]{human:boyd-vandenberghe}; CRP8--9 supplies
its exact original-fibre realization.

\subsection{The killed-space angle from actual root evaluations}
For $v=0$ set $C_{\rm ang}=0$; then $K=I$. Suppose $v>0$ and choose
the following actual upper roots, with no rescaling:
\[
z_h=(2h-1)\gamma+i\delta,\qquad
\beta_h=c+iz_h=\beta_{a,(a-1)/2,(a-1)/2+h},\quad0\le h<v.
\tag{CRP10}
\]
All their indices are in the original grid, since $a\ge257$ and
$v\le80$. Let $\mathcal E=(\phi_{n,s,c}(\beta_h))_{h<v,n\le N}$,
$\mathcal K=\mathcal E\mathcal E^*$, and let $\mathcal L$ be the first
$v$ columns of $\mathcal E$. Evaluation annihilates $\operatorname{im}\Phi$.
The matrix $\mathcal L$ is invertible, because its columns are degree
$0,\ldots,v-1$ polynomials with nonzero leading coefficients at distinct
points. Let $Z$ insert the first $v$ native coordinates. Sylvester's
identity, followed by inclusion of the evaluation adjoint range, gives
\[
\begin{split}
\det K&=\det(I-\Phi^*ZZ^*\Phi)
 =\det(Z^*(I-\Phi\Phi^*)Z),\\
Z^*(I-\Phi\Phi^*)Z&\succeq
 Z^*\mathcal E^*\mathcal K^{-1}\mathcal E Z
 =\mathcal L^*\mathcal K^{-1}\mathcal L,\\
\det K&\ge |\det\mathcal L|^2/\det\mathcal K>0.
\end{split}\tag{CRP11}
\]
For the determinant monotonicity used here, write a positive pair as
$Y^{1/2}(I+C)Y^{1/2}$ with $C\succeq0$ and multiply its positive
eigenvalues. Thus the inequality is also valid when the evaluation
space is a proper part of the full orthogonal complement.

Put
\[
\begin{gathered}
Z_v=\max_{h<v}|z_h|,\quad
\mathcal A=e^2(N+1)^{1+s/2}(2N+1)^{1+Z_v},\\
\mathcal V=(2\gamma)^{v(v-1)/2}\prod_{n=1}^{v-1}n!,\\
C_{\rm ang}=\left[v\log\mathcal A+
 \sum_{n=0}^{v-1}\log(n!(s/2)_n)-2\log\mathcal V\right]_+.
\end{gathered}\tag{CRP12}
\]
Here $\mathcal V$ is the literal Vandermonde modulus of the selected
$z_h$'s. On $|t|=R=N/(N+1)$, the original generating series obeys
\[
|\arctan t|\le\operatorname{artanh}R=\tfrac12\log(2N+1),\qquad
|(1+t^2)^{-s/4}|^2\le(1-R^2)^{-s/2}\le(N+1)^{s/2}.
\]
Cauchy's coefficient formula, $R^{-2n}\le e^2$ for $n\le N$, and
$\rho_n^2\le n!/(1/2)_n\le2n+1$ imply
\[
|\phi_{n,s,c}(\beta_h)|^2\le
 \frac{e^2}{M_s}(N+1)^{s/2}(2N+1)^{1+Z_v}.
\tag{CRP13}
\]
The last bound on $\rho_n$ follows by induction: multiplying $2n+1$
by $(n+1)/(n+1/2)$ gives $2n+2\le2n+3$. Hadamard's inequality is
obtained by orthogonalizing the rows, whose residual lengths do not
exceed their lengths. It gives $\det\mathcal K\le(\mathcal A/M_s)^v$.
The leading coefficients in CRP4 give exactly
$|\det\mathcal L|^2=\mathcal V^2/\prod_{n<v}h_{n,s}$.
Combining them with CRP11 cancels exactly the $v$ factors $M_s$, proving
\[
\boxed{0\le-\log\det K\le C_{\rm ang}
 =O_{\rm actual}(a\log q)=o(q).}\tag{CRP14}
\]
The constant depends on the original fixed quartet and conductor,
including $v$ and $\gamma$. The estimate is uniform in the stated
$M,s$ range; no bound uniform in a varying hypothetical quartet is
claimed.

\subsection{Summed absolute pivots from all exterior ranks}
For clarity the exact finite LT constants used below are
\[
\begin{gathered}
F_0^{\flat}=4+4\sqrt{\delta^2+\gamma^2},\quad
B_0=2^{v+8}\bigl(\sum|a_{uw}|\bigr)e^{2\pi\gamma},\\
B_2=B_0\{F_0^{\flat}(F_0^{\flat}+1)+4vF_0^{\flat}+4v(v+1)\},
\quad H_*=B_0(1+\sqrt3)+2\sqrt{10}B_2,\\
H_s=\begin{cases}\max(1,H_*\sqrt{2v+1}),&s=1,\\
\max(1,H_*),&s=a,\end{cases}\qquad
\kappa_s=\begin{cases}1+\sqrt{8/3},&s=1,\\2,&s=a,\end{cases}\\
J=d\log^+(1/|g_0|)+\tfrac12(s/2-1)_+v\sum_{j=1}^d1/j,\\
F_r=r\log H_s+4(d-r)\log\kappa_s\ (1\le r\le d),
\quad F_0=0,\quad I_r=J+F_{d-r}\ (1\le r\le d),\quad I_0=0.
\end{gathered}\tag{CRP15}
\]
LT7--22/LRC6--9 proves
$\log\|\bigwedge^r A\|\le F_r$ and
$\log\|\bigwedge^r A^{-1}\|\le I_r$ for every original rank.
The same bounds hold for ordinary transposes. Put
$F^*=\max_{0\le r\le M+1}F_r$ and
$I^*=\max_{0\le r\le M+1}I_r$.

We prove the needed selected-pivot estimates, rather than replacing
them by a diagonal trial norm. For an index set $J_+$, the actual
degree-order pivot minimizes the column residual over all earlier
columns. Restricting the minimum to earlier columns in $J_+$ can only
increase it. Multiplying these selected Gram--Schmidt residuals yields
\[
\prod_{r\in J_+}\delta_r\le\det H[J_+,J_+]
 =\|\bigwedge^{|J_+|}(AB_{J_+})\|^2
 \le e^{2F_{|J_+|}}.\tag{CRP16}
\]
The last bound uses that $B_{J_+}$ is a contraction, since the
corresponding columns of $\Phi$ are orthonormal.

For another set $J_-$ first condition every column in $J_-$ on all
columns in its complement, and then on its preceding columns in
$J_-$. These spaces include all original preceding columns, so the
new residuals are no larger. Their product is the Schur determinant
$\det H/\det H[J_-^c,J_-^c]$. The block inverse formula therefore gives
\[
\prod_{r\in J_-}\delta_r^{-1}
 \le\det(H^{-1}[J_-,J_-]).\tag{CRP17}
\]
To estimate the right side, put $B=WK^{1/2}$ where $W^*W=I$, and
$C=W^*A^*AW$. Complete $W$ to an orthonormal basis. If the resulting
blocks of $A^*A$ are $\left(\begin{smallmatrix}C&T\\T^*&R\end{smallmatrix}\right)$,
the original full Schur formula gives
\[
C^{-1}\preceq(C-TR^{-1}T^*)^{-1}=W^*(A^*A)^{-1}W.
\tag{CRP18}
\]
For a zero-dimensional complementary block this is equality. For
any isometric rank-$t$ column matrix $Y$, determinant monotonicity
and exterior norms give
$\det(Y^*C^{-1}Y)\le\|\bigwedge^t A^{-1}\|^2\le e^{2I_t}$.
Taking $Y$ to be a top-$t$ eigenvector frame proves
$\|\bigwedge^t C^{-1}\|\le e^{2I_t}$. Since
$H^{-1}=K^{-1/2}C^{-1}K^{-1/2}$ and every eigenvalue of $K$ lies in
$(0,1]$, the product of any $t$ eigenvalues of $K^{-1}$ is at most
$\det K^{-1}$. Thus
\[
\det(H^{-1}[J_-,J_-])\le e^{2I_{|J_-|}}/\det K.\tag{CRP19}
\]
Choose $J_+=\{r:\log\delta_r>0\}$ and
$J_-=\{r:\log\delta_r<0\}$. Empty sets contribute zero. Taking
logarithms in CRP16--19 and adding proves the finite result
\[
\boxed{\sum_{r=0}^M|\log\delta_r|
 \le 2F^*+2I^*+C_{\rm ang}=O_{\rm actual}(q).}\tag{CRP20}
\]
Indeed $d=O(q)$, $\log H_s,\log\kappa_s=O_{\rm actual}(1)$,
and $J=O_{\rm actual}(q+a\log q)=O_{\rm actual}(q)$.
This controls all pivots together and uses the actual polynomial
ideal position through CRP10--14. It does not estimate the sum of
the trial-to-attained slacks $\log(\tau_r/\eta_r)$ by $O(q)$.

\subsection{Exact directional transfer after the scalar reconstruction}
The companion scalar proof supplies the corrected profile
\[
\begin{split}
c_r&=\Delta\mathfrak f(r/q)
-\tfrac12\{(r+\Delta)\log(1+\Delta/r)-\Delta\},\quad r>0,\\
c_0&=\tfrac\Delta2\{1+\log(q/\Delta)\},\qquad
\mathcal E_c\ge\sum_{r=0}^M|\gamma_r-c_r|
 =O_{\rm actual}(q),\qquad c_r^+=\max(c_r,0).
\end{split}\tag{CRP21}
\]
Here $\mathcal E_c$ is its explicit finite sum of single-norm errors,
CRP20 and scalar correction errors; it is not an assumed smallness
condition. The subsequent proof transfers this proved error without
altering any original directional matrix.

We specify the map into the analytic coefficient frame as well. Let
$e_b=\chi_{a-8}\pi_b^-/\sqrt{\nu_b^-}$ be the original orthonormal
lower-ideal basis. In that basis write the full relation columns as
$\mathcal R_r=[\mathcal L_r;H_r]$, where the first block has $m$ rows.
Degree and the monic coefficient give
$H_r[t,t]=\sqrt{\nu^-_{m+t}}$ and $H_r[t,u]=0$ for $t>u$.
Thus the mutually inverse coefficient maps are $z=H_rb$ and
$b=H_r^{-1}z$, and the relation graph is
$\mathcal R_rH_r^{-1}=[X_r;I]$, $X_r=[x_0,\ldots,x_r]$.
The leading columns persist with $r$ because the inverse of an upper
triangular leading block persists. Put $\Omega_r=I_m+X_rX_r^*$.
For the original analytic columns $F_b=N_b/E_A$ and original
complex-linear functionals $\Lambda_{F_b}(S'^j)=F_b^{(j)}(0)$, define
$A_{\rm an}[h,b]=\Lambda_{F_b}(e_h)$ for $0\le h<m$.
This $m\times m$ matrix is the analytic coefficient map NTG16; the
$d\times d$ matrix $A_{D,s}$ in CRP5 is the conductor on polynomials.
Their exact common intermediate map is the full relation graph just
displayed: CRP8 constructs every relation polynomial before the
graph coordinate change, and functional annihilation gives
$X_r^{\mathsf T}A_{\rm an}+F_{\rm hi}=0$ after it.
NTG18--20's full lower-evaluation minimum therefore gives
$D_r=A_{\rm an}^*\overline{\Omega_r}A_{\rm an}$ in the original
analytic frame, with $D_{-1}=A_{\rm an}^*A_{\rm an}$.
Here the row index notation means exactly
$D_r=D^{\rm nat}_{L+r}=D^{\rm target}_{q+r}$, including $r=-1$.
The analytic frame is onto by the original complete functional map,
so $A_{\rm an}$ is invertible. Moreover
$\det\Omega_r/\det\Omega_{r-1}=\eta_r/\nu^-_{m+r}$: taking the
Gram of $[X_r;I]H_r$ gives
$\det\mathcal R_r^*\mathcal R_r=(\prod_{t\le r}\nu^-_{m+t})\det\Omega_r$,
while completing the same relation square gives
$\det\mathcal R_r^*\mathcal R_r=\prod_{t\le r}\eta_t$.
This proves the morphism carrying the scalar pivots to the native
analytic row, before any directional estimate.

The row after the full original relation minimum is retained with
the ordinary-transpose convention of NTG/NFI:
\[
\mathbf u_r=\frac{\boldsymbol\rho_{L+r}}{\sqrt{1+\zeta_{L+r}}}
   =-i^{-(L+r)}x_r^{\mathsf T}A_{\rm an},\quad
D_r=D_{r-1}+\mathbf u_r^*\mathbf u_r,\quad
t_r=\mathbf u_rD_{r-1}^{-1}\mathbf u_r^*=e^{\gamma_r}-1.\tag{CRP22}
\]
The symbol $\boldsymbol\rho_{L+r}$ here is the original residual \emph{row};
the scalar $\rho_n$ in CRP4--5 has already served only to specify the
conductor matrix. Their original distinction is preserved.
The full relation Gram defining $\zeta_{L+r}$ and this row is the
one used in ATA8 and NFI, not its diagonal. For the four-copy row
$U_r=\operatorname{diag}(\mathbf u_r,4)$ let
$\mathcal D_0=\operatorname{diag}(D_{r-1},4)$.
For any actual flag column map $X$, the determinant update is
\[
\begin{split}
K_{X,r}&=U_rX(X^*\mathcal D_0X)^{-1}X^*U_r^*,\\
\Delta\log\det(X^*\mathcal D X)&=\log\det(I_4+K_{X,r}).
\end{split}\tag{CRP23}
\]
The determinant identity follows by factoring the old positive Gram
and applying $\det(I+AB)=\det(I+BA)$; the inverse rank-one update is
the classical formula recorded by Hager \cite[p.~221, (2)]{human:hager1989}.
For empty flags the matrices and increments are zero.

If $t_r>0$, define
$Z_r=t_r^{-1/2}U_r\mathcal D_0^{-1/2}$ and let $P_{X,r}$ be the
orthogonal projector onto $\mathcal D_0^{1/2}\operatorname{im}X$.
Then
\[
Z_rZ_r^*=I_4,\quad M_{X,r}=Z_rP_{X,r}Z_r^*,\quad
0\preceq M_{X,r}\preceq I_4,\quad K_{X,r}=t_rM_{X,r}.
\tag{CRP24}
\]
All four eigenvalues $\lambda_{X,r,b}$ of $M_{X,r}$, including zero,
are retained. If $t_r=0$, the original row is zero, and all its exact
and approximating directional contributions are defined to be zero;
no direction is assigned to a zero row.

For $g\ge0$ define $F(g,\lambda)=\log(1+(e^g-1)\lambda)$ and
$C(g,\lambda)=(g+\log\lambda)_+$, with $C(g,0)=0$.
Direct differentiation and the two cases $e^g\lambda\le1$ and
$e^g\lambda\ge1$ give
\[
0\le F-C\le\log2,\qquad
\partial_gF=\frac{e^g\lambda}{1+(e^g-1)\lambda}\in[0,1].
\tag{CRP25}
\]
Moreover $\partial_gF$ is increasing in $\lambda$. The difference
$d_g(\lambda)=F-C$ increases from zero up to
$\lambda=e^{-g}$ and decreases thereafter to zero at $\lambda=1$.
At $g=0$ it vanishes identically. Hence it can be written as
$d_g=u_g-v_g$, with $u_g,v_g$ increasing, both starting at zero and
both ending at the same maximum $\log(2-e^{-g})\le\log2$.

The original two chains are
\[
J\subset W_i\subset W\subset T,\qquad
K_i\subset I\subset K_c\subset S.\tag{CRP26}
\]
Projection order and the Rayleigh min--max principle imply order of
the four eigenvalues at each position along each chain. For four
ordered scalar arguments $z_1\le z_2\le z_3\le z_4$, the sum of
the two disjoint increments of any increasing $u$ is between zero
and $u(1)-u(0)$. Applying this first to $\partial_gF$ shows that the
derivative of each chain sum is in $[0,4]$. Their difference,
with the exact signs below, is therefore $4$-Lipschitz in $g$.
Applying it to $u_g$ and $v_g$ shows that the discrepancy of each
chain sum lies between $-4\log2$ and $4\log2$. Thus the difference
of the two chains has discrepancy at most $8\log2$.

Retain the eight \emph{labelled occurrences}, even if some underlying
subspaces coincide:
\[
(T,K_c,W_i,K_i,W,S,J,I),\qquad
(\sigma_X)=(+,+,+,+,-,-,-,-).\tag{CRP27}
\]
For any original row set $\mathscr I$ with its original weights
$0\le\vartheta_r\le2$, put
$\widetilde f_{X,r}=\sum_{b=1}^4C(c_r^+,\lambda_{X,r,b})$ on nonzero
rows. Since $\gamma_r\ge0$,
$|\gamma_r-c_r^+|\le|\gamma_r-c_r|$. Freeze the actual matrices
$M_{X,r}$ while applying the scalar Lipschitz estimate; this does
not assert that the flags of a different metric equal these flags.
Summation proves
\[
\boxed{\left|\mathcal L_e^{\mathscr I}-
 \sum_{r\in\mathscr I}\vartheta_r\sum_X\sigma_X\widetilde f_{X,r}\right|
 \le8\mathcal E_c+8\log2\sum_{r\in\mathscr I}\vartheta_r.}
\tag{CRP28}
\]
For a single original nested quotient pair $A\subset B$, the
derivative of its four-eigenvalue difference is in $[0,4]$ and its
clipped discrepancy is at most $4\log2$. Therefore
\[
\boxed{|V_{B/A}^{\mathscr I}-\widetilde V_{B/A}^{\mathscr I}|
 \le8\mathcal E_c+4\log2\sum_{r\in\mathscr I}\vartheta_r.}
\tag{CRP29}
\]
This applies to STF's unchanged complementary pairs $W/W_i$,
$U/T$, $S/K_c$, and $I/K_i$, where $U$ is the full original ambient
space and its directional compression is $I_4$. It includes the
original quotient minima through the determinant differences in
CRP23, not new quotient ranks or substituted forms.

\subsection{One-copy allocations, exact quotient maps, and backward use}
Let $B_a$ be the original onto observation map, $J_a$ an injective
matrix with image $\ker B_a$, and let the fixed onto-value frame be
unchanged. Its restricted and quotient forms are exactly
\[
G_{R,r}=J_a^*D_rJ_a,\qquad
G_{K,r}=(B_aD_r^{-1}B_a^*)^{-1}.\tag{CRP30}
\]
If $Y$ is any right inverse of $B_a$, the fixed square frame
$[J_a,Y]$ has determinant independent of $r$. Completing the full
$J_a$ coefficient square proves
\[
|\det[J_a,Y]|^2\det D_r=\det G_{R,r}\det G_{K,r}.
\tag{CRP31}
\]
Thus its original four-endpoint return loses no frame factor.
For a nonzero row let $P_{R,r}$ be the orthogonal projector onto
$D_{r-1}^{1/2}\ker B_a$ and retain the actual projection energy
$p_r=t_r^{-1}\mathbf u_rD_{r-1}^{-1/2}P_{R,r}
D_{r-1}^{-1/2}\mathbf u_r^*\in[0,1]$.
Then CRP23 in one copy and CRP31 give
\[
\Delta S_R=F(\gamma_r,p_r),\qquad
\Delta S_K=\gamma_r-F(\gamma_r,p_r).\tag{CRP32}
\]
Set $\kappa_r=-\log p_r$, allowing $\kappa_r=+\infty$ when $p_r=0$.
Both functions in CRP32 are $1$-Lipschitz in the scalar argument.
CRP25 and $g-C(g,p)=\min(g,-\log p)$ now prove
\[
\begin{gathered}
\widetilde S_R=\sum_r\vartheta_r(c_r^+-\kappa_r)_+,\qquad
\widetilde S_K=\sum_r\vartheta_r\min(c_r^+,\kappa_r),\\
\boxed{|S_X-\widetilde S_X|\le2\mathcal E_c+
 \log2\sum_r\vartheta_r,\qquad X=K,R.}\tag{CRP33}
\end{gathered}
\]
Zero rows contribute zero in both expressions, as above. The exact
identities remain $S_K+S_R=\mathcal T$ and
$4S_K=\mathcal L_e+V_b+V_{\rm rem}+V_s+V_A$; every term uses the same
original row and quotient maps. CRP28--33 do not evaluate the
$p_r$ or four-dimensional spectra from dimensions or from the total.

The precise replacement in older PCB/LRS/FTR row receivers is the
following: replace the radial function $\Delta\mathfrak f(r/q)$
by CRP21's $c_r$ (and by $c_r^+$ only inside clipped directional
expressions), replace its weighted $O(q\log q)$ error by the
explicit $2\mathcal E_c$, and use CRP28--33 with the original
weights, flags and quotient pairs. Keeping the old uncorrected
profile and merely changing its allowance to $O(q)$ is invalid:
the subtracted correction has full-window weighted sum
$128q\log a+O(q)$, proved in the companion scalar calculation.
For prefix/tail receivers also replace the integrated centres by
the companion $\mathcal A_\Delta$ formulas, including their endpoint
terms. The old trial-to-attained and metric-chain comparisons remain
their own proved comparisons. Every old sealed edition is preserved;
this is a successor theorem with an explicit map back to each receiver.
\endgroup

% END PRESERVED INPUT 03_CORRELATED_PIVOT_AND_DIRECTION_PROOF.tex
\endgroup

\clearpage
\begingroup


\setcounter{equation}{0}
\renewcommand{\theequation}{A02CRS.\arabic{equation}}
\section*{Corrected radial profile and prefix, tail and proper-source returns}
\addcontentsline{toc}{section}{Corrected radial profile and prefix, tail and proper-source returns}
% BEGIN PRESERVED INPUT 04_CORRELATED_PROFILE_SCALAR_PROOF.tex SHA256 4c9069c298624a6dc42b253880002190d8af3c9d26d83ac28e3ad21afa261ae7
% Additive reconstruction of Untitled 1775.md, equations (9)--(20),(25).
% Inclusion file: requires amsmath,amssymb,hyperref.
\begingroup
\section{The corrected native radial profile and its endpoint defect}
\label{a02crs:sec:CRS}

This calculation receives the literal conductor-pivot identity from the
accompanying correlated-row proof and proves the scalar transformations,
finite error transfers and coefficients. Its source is the first
continuation in \texttt{Untitled 1775.md}, equations (9)--(20) and (25).
Earlier providers are LRS1--18 and LRS24--25, PCB1--6, the direct elliptic
derivative proof and ATA35--40. Their exact files are pinned in the
adjacent \texttt{ACCEPTANCE.json}. No sealed predecessor is edited.

\subsection{Original parameters and finite input identities}

Keep the actual simple quartet, period, branch and conductor. Its moment
is $\mu_v=E_A^{(v)}(0)\ne0$, $0\le v\le80$. Write
\[
\begin{gathered}
 a\equiv1\pmod4,\quad a\ge257,\quad s\in\{1,a\},\quad m=\Delta-v,\\
 q=(a+1)^2,\quad Q=(a-7)^2,\quad \Delta=q-Q=16a-48,\\
 \ell=(s-1)/4,\quad n_+=q+\ell,\quad n_-=Q+\ell,\quad
 \alpha_A=v!/\mu_v,\quad \ell_r=\mu_v\binom{q+r}{v}.
\end{gathered}\tag{CRS1}
\]
The degree-$a-8$ lower packet keeps the same order $s$. Retain the two
original LRC13 guards:
\[
 n_b=(b+1)^2+\ell\ge100,\qquad
 \max\{b\sqrt{\delta^2+\gamma^2},2\ell+1/2\}
 \le2^{-17}n_b,\qquad b=a,a-8.                         \tag{CRS2}
\]
Any GPA provider used elsewhere keeps its own guard. The present
calculation uses the retained Gamma LRC endpoints with their original
weight. All sums allow $0\le M\le q+32a$.
The original degree ranges hold: $M\le2n_+$ and $m+M\le2n_-$.
For the latter,
$q-3\Delta-32a+v+2\ell=a^2-78a+145+v+2\ell>0$ for $a\ge257$.
All scalar arguments below lie in $[0,3]$, since
$M+\Delta\le q+48a-48<3Q$.

Let $\nu^+_{r,s},\nu^-_{m+r,s}$ be the original full monic minima and
$\eta_{r,s}$ the attained preceding-relation minimum. Put
\[
 h_{d,s}=(2\pi)^{s/2}d!(s/2)_d,\quad
 \gamma_{r,s}=\log(\eta_{r,s}/\nu^-_{m+r,s})\ge0.
\]
The squared Cholesky pivot $\delta_r$ of the complete conductor-image
Gram in its original degree order satisfies
\[
 \delta_r=|\ell_r|^2\eta_{r,s}/\nu^+_{r,s}.             \tag{CRS3}
\]
This is the affine fibre identity ATA6 followed by the full Schur
minimum, as reconstructed in the accompanying conductor proof.

For root count $p=q,Q$, let $e_{p,j,s}$ be the larger absolute distance
of the two finite LRS4 endpoints from $2(p+\ell)\psi(j/(p+\ell))$.
Use the scalar-moment endpoints at $j=0,1$. LRS11 proves
\[
\begin{split}
 \log\nu^{(p)}_{j,s}-\log h_{p+j,s}
 &=2n\psi(j/n)+\epsilon_{p,j,s},\quad
 |\epsilon_{p,j,s}|\le e_{p,j,s},\quad n=p+\ell,\\
 e_{p,j,s}&\le C_{\rm actual}
 \left(1+\log^+\frac{n}{j+1}+\frac{(\ell+1)^2}{n+j}\right),
 \qquad 0\le j\le2n .
\end{split}                                                    \tag{CRS4}
\]
Its retained Gamma parity parameter is $(n+\varepsilon-3/4)/h$.
The positive-real Stirling remainder used in LRS7 has human provenance
in NIST DLMF 5.11.1 and 5.11(ii) \cite[\S5.11]{a02crs:crs:dlmf5}.
The norm masses and parity parameter are unchanged here.
The errors are summable: integrating $\log(n/x)$ on $(0,n)$ gives
$\sum_{j=0}^{2n}\log^+(n/(j+1))\le n$, and
$\sum_{j=0}^{2n}(\ell+1)^2/(n+j)\le3(\ell+1)^2=O(q)$.
There are $O(q)$ constant terms. Thus both original row-error sums
have size $O_{\rm actual}(q)$.

\subsection{The exact elliptic dictionary}

Use the real Legendre modulus convention in NIST DLMF 19.2.4,
19.2.5 and 19.2.8 \cite[\S19.2]{a02crs:crs:dlmf19}:
\[
 K(k)=\int_0^{\pi/2}(1-k^2\sin^2\theta)^{-1/2}\,d\theta,\quad
 E(k)=\int_0^{\pi/2}(1-k^2\sin^2\theta)^{1/2}\,d\theta.
\]
The original parameter and functions are
\[
\begin{gathered}
 k(t)=\sqrt{1-\rho(t)^2},\qquad
 \rho(t)K(k(t))/E(k(t))=(1+t)^{-1},\\
 f(t)=\log\frac{1+\rho(t)}{1-\rho(t)},\qquad
 \psi(t)=\log\frac{1+\rho(t)}{E(k(t))}
                         +t\log\frac{k(t)}{E(k(t))}.
\end{gathered}                                                   \tag{CRS5}
\]
Equivalently $\psi(t)=(1+t)(1-\log(1+t))-t\lambda(2/t,\pi/t)/4$
with the original Robin convention. The pinned direct integral proof
establishes uniqueness of $\rho$, the derivative identities, and
the following global and endpoint bounds:
\[
\begin{gathered}
 \psi'(t)=\log(k(t)/E(k(t)))<0,\quad
 2\psi(t)-2(1+t)\psi'(t)=f(t),\quad \psi(0)=a_0=\log(4/\pi),\\
 0<-tf'(t)\le1,\quad
 |f(t)+\tfrac12\log t|\le(22/5)\sqrt t,\quad
 |tf'(t)+\tfrac12|\le9\sqrt t\quad(0<t\le1/20).
\end{gathered}                                                   \tag{CRS6}
\]
For the dictionary identity, substitute CRS5: its left side is
$2\log((1+\rho)/k)=\log((1+\rho)/(1-\rho))$ exactly.
The derivative bounds in the provider are proved by differentiating
the displayed defining integrals, without differentiating an
asymptotic remainder.
Put
\[
 \eta(t)=f(t)+\tfrac12\log t,\quad \eta(0)=0,\qquad
 H_\eta=\int_0^3|\eta'(t)|\,dt
          \le18/\sqrt{20}+\tfrac12\log60.              \tag{CRS7}
\]
Indeed $|\eta'|\le9/\sqrt t$ up to $1/20$, and thereafter
$-tf'\in[0,1]$ gives $|\eta'|\le1/(2t)$. This proves absolute
continuity of $\eta$ on $[0,3]$ and the finite bound.
Also $\int_0^3|\psi'|=\psi(0)-\psi(3)<\infty$.

\subsection{A finite summed scalar error}

Define
\[
 h_r^{\rm prof}=2n_+\psi(r/n_+)-2n_-\psi((m+r)/n_-)+\Gamma_r.
                                                               \tag{CRS8}
\]
The factorial correction is exactly
\[
\begin{split}
 \Gamma_r&=\log\frac{h_{q+r,s}}{h_{q+r-v,s}}-2\log|\ell_r|\\
 &=2\log|\alpha_A|+
   \sum_{j=0}^{v-1}\log\left(1+\frac{s/2-1}{q+r-j}\right).
\end{split}                                                    \tag{CRS9}
\]
The norm ratio contributes
$\prod_{j<v}(q+r-j)(q+r-j-1+s/2)$ and
$|\ell_r|^2=|\mu_v/v!|^2\prod_{j<v}(q+r-j)^2$.
Their quotient proves CRS9 with the original mass cancelled only
inside this matched ratio. Empty products apply at $v=0$.
Combining CRS3--4 gives the exact equality
\[
 \gamma_{r,s}-h_r^{\rm prof}
   =\log\delta_r+\epsilon_{q,r,s}-\epsilon_{Q,m+r,s}.  \tag{CRS10}
\]
Let $B_M=2F_*+2I_*+C_{\rm ang}$ be the finite bound for
$\sum_{r=0}^M|\log\delta_r|$ proved in the accompanying conductor
calculation. Its proved size is $O_{\rm actual}(q)$.
Consequently
\[
 E_1=B_M+\sum_{r=0}^M(e_{q,r,s}+e_{Q,m+r,s}),\qquad
 \sum_{r=0}^M|\gamma_{r,s}-h_r^{\rm prof}|
 \le E_1=O_{\rm actual}(q).                           \tag{CRS11}
\]
No operator-norm error is charged independently to each pivot.

Set $F(n,r)=2n\psi(r/n)$. Direct differentiation and CRS6 give
$F_n-F_r=f(r/n)$. Integrate on $(n,r)=(n_+-u,r+u)$ and retain
the actual $m=\Delta-v$:
\[
\begin{split}
 F(n_+,r)-F(n_-,r+m)
 =\int_0^\Delta f\left(\frac{r+u}{n_+-u}\right)\,du
       +2\int_{r+m}^{r+\Delta}\psi'(b/n_-)\,db .
\end{split}                                                    \tag{CRS12}
\]
At $r=0$ these are convergent improper integrals by CRS6.
For $x>0$ put
\[
 c(x)=\Delta f(x/q)
 -\tfrac12\{(x+\Delta)\log(1+\Delta/x)-\Delta\}.        \tag{CRS13}
\]
Its exact integral presentation, valid also at $x=0$, is
\[
 c(x)=-\tfrac12\int_0^\Delta\log((x+u)/q)\,du
                     +\Delta\eta(x/q),\qquad
 c_0=c(0)=\tfrac{\Delta}{2}(1+\log(q/\Delta)),\quad c_r=c(r).
                                                               \tag{CRS14}
\]
This proves the initial-row constant with no unspecified offset.

All scalar terms between CRS12 and CRS13 have the following
finite summed bound:
\[
\begin{split}
 B_\Gamma&=(M+1)
 \left(2|\log|\alpha_A||+\frac{2v|s/2-1|}{q-v+1}\right),\\
 B_v&=2(v+1)n_-[\psi(0)-\psi(3)],\\
 B_{\rm den}&=(M+1)\Delta(\ell+\Delta)/Q,\qquad
 B_\eta=(\Delta^2/2+\Delta)H_\eta,\\
 B_{\rm sc}&=B_\Gamma+B_v+B_{\rm den}+B_\eta
                                      =O_{\rm actual}(q).
\end{split}                                                    \tag{CRS15}
\]
For $B_\Gamma$, each argument $(s/2-1)/(q+r-j)$ has modulus
less than $1/2$, and $|\log(1+x)|\le2|x|$ on that interval.
For $B_v$, the intervals $[r+m,r+\Delta]$ overlap at most
$v+1$ times almost everywhere; integrate $2|\psi'(b/n_-)|$,
then substitute $t=b/n_-$. At $v=0$ the actual error is zero.
For $B_{\rm den}$, CRS6 yields, almost everywhere,
\[
 \left|f\left(\frac{r+u}{q+\ell-u}\right)
                    -f\left(\frac{r+u}{q}\right)\right|
 \le|\log(q/(q+\ell-u))|\le(\ell+\Delta)/Q .
\]
This includes the convergent limiting calculation at $r=0$.
The remaining regular difference is
$\int_0^\Delta[\eta((r+u)/q)-\eta(r/q)]\,du$.
For fixed $u$, the intervals $[r/q,(r+u)/q]$ overlap at most
$u+1$ times almost everywhere. Sum their absolute derivative
integrals and integrate in $u$ to obtain $B_\eta$.
These four calculations prove
\[
 \sum_{r=0}^M|h_r^{\rm prof}-c_r|\le B_{\rm sc}.        \tag{CRS16}
\]
Thus the original, possibly smaller finite radius satisfies
\[
\begin{gathered}
 E_c=E_1+\sum_{r=0}^M|h_r^{\rm prof}-c_r|
                 \le E_1+B_{\rm sc}=O_{\rm actual}(q),\\
 \boxed{\displaystyle\sum_{r=0}^M|\gamma_{r,s}-c_r|\le E_c.}
\end{gathered}                                                   \tag{CRS17}
\]
No independently known total or saturation argument enters this proof.

\subsection{Primitive and original endpoint weights}

Let $\mathfrak F(0)=0$, $\mathfrak F'=\psi$ be the original LRS15
primitive. Integration of CRS6 gives
\[
 \mathcal J(t):=\int_0^t f(u)\,du
 =4\mathfrak F(t)-2(1+t)\psi(t)+2a_0,\qquad
 2\mathcal J(1)=C_\partial.                            \tag{CRS18}
\]
Its original energy convention is retained:
\[
 \mathfrak F(t)=\frac{t^2}{4}F_{\rm energy}(2/t,\pi/t)
 +\frac{(1+t)^2}{2}\left(\frac32-\log(1+t)\right)-\frac34,
 \qquad F_{\rm energy}=-E_{\min}.
\]
For $R>0$, define
\[
\begin{split}
 A_\Delta(R)=\tfrac12[&(R+\Delta)^2\log(R+\Delta)
       -R(R+2\Delta)\log R\\
       &-\Delta^2\log\Delta-R\Delta],\qquad A_\Delta(0)=0 .
\end{split}                                                    \tag{CRS19}
\]
Direct differentiation gives
$A_\Delta'(R)=(R+\Delta)\log(1+\Delta/R)-\Delta$,
twice the correction subtracted in CRS13. Its limit at zero is zero.
Consequently
\[
 2\int_0^R c(x)\,dx=2\Delta q\mathcal J(R/q)-A_\Delta(R).
                                                               \tag{CRS20}
\]
The logarithmic part of CRS14 has derivative
$-\tfrac12\log(1+\Delta/x)$ and hence total variation
$\tfrac12[(q+\Delta)\log(q+\Delta)-q\log q-\Delta\log\Delta]$
on $[0,q]$. Its other term has variation at most $\Delta H_\eta$.
Therefore
\[
\begin{split}
 {\rm Var}_{[0,q]}c&\le V_c
 :=\tfrac12[(q+\Delta)\log(q+\Delta)-q\log q-\Delta\log\Delta]
                                    +\Delta H_\eta,\\
 V_c+|c_0|&=O(a\log q)=O(q).
\end{split}                                                    \tag{CRS21}
\]
For the size, the bracket equals
$q\log(1+\Delta/q)+\Delta\log(1+q/\Delta)$; use
$\log(1+x)\le x$. The profile is absolutely continuous including zero.

The original weights are
\[
 \mathcal T^{\le R}=\gamma_0+2\sum_{r=1}^R\gamma_r,\qquad
 \mathcal T^{>R}=2\sum_{r=R+1}^{q-1}\gamma_r+\gamma_q,\quad 0\le R<q .
                                                               \tag{CRS22}
\]
On an integer interval $[u,w]$, the doubled trapezoid
$c_u+2\sum_{u<r<w}c_r+c_w$ differs from $2\int_u^w c$
by at most $2{\rm Var}_{[u,w]}c$. On each unit interval the
absolute error is bounded by the integral of
$|c_j-c(x)|+|c_{j+1}-c(x)|$, at most twice its variation;
sum the unit intervals. The prefix in CRS22 has one extra
$c_R$ relative to this trapezoid and the tail has its negative.
Since $|c_R|\le|c_0|+V_c$, CRS17 proves simultaneously
\[
\begin{split}
 C^{\rm pre}(R)&=2\Delta q\mathcal J(R/q)-A_\Delta(R),\\
 C^{\rm tail}(R)&=2\Delta q[\mathcal J(1)-\mathcal J(R/q)]
                                    -A_\Delta(q)+A_\Delta(R),\\
 |\mathcal T^{\le R}-C^{\rm pre}(R)|,\
 |\mathcal T^{>R}-C^{\rm tail}(R)|
 &\le2E_c+|c_0|+3V_c=O_{\rm actual}(q).
\end{split}                                                    \tag{CRS23}
\]
The row weights cost at most $2E_c$. At $R=0$ the prefix is
$\gamma_0$ and its centre is zero, so the same bound holds directly.

\subsection{Moving-cutoff logarithms with a finite remainder}

Taylor's integral formula for $g(y)=(1+y)^2\log(1+y)$,
with $g'''(y)=2/(1+y)$, gives for every $R>0$
\[
\begin{split}
 A_\Delta(R)
 &=\tfrac{\Delta^2}{2}\log(R/\Delta)+\tfrac{3\Delta^2}{4}
  +\frac{R^2}{2}\int_0^{\Delta/R}
                         \frac{(\Delta/R-u)^2}{1+u}\,du,\\
 0&\le A_\Delta(R)-\tfrac{\Delta^2}{2}\log(R/\Delta)
                -\tfrac{3\Delta^2}{4}\le\frac{\Delta^3}{6R}.
\end{split}                                                    \tag{CRS24}
\]
This strengthens the received remainder and does not require $R\ge\Delta$.
For the prescribed $R_*=\lfloor a^{3/2}/\log a\rfloor$,
\[
\begin{gathered}
 \Delta^2/2=128q-1024a+1024,\\
 \log(R_*/\Delta)=\tfrac12\log a-\log\log a-\log16+O(1/a),\\
 \log(q/\Delta)=\log a-\log16+O(1/a),\qquad
 \Delta^3/R_*=O(a^{3/2}\log a)=O(q).
\end{gathered}
\]
The floor changes $\log R_*$ by $O(\log a/a^{3/2})$, absorbed
above. Inserting these expressions in CRS24 gives
\[
\begin{split}
 A_\Delta(R_*)&=64q\log a-128q\log\log a+O(q),\\
 A_\Delta(q)&=128q\log a+O(q),\\
 A_\Delta(q)-A_\Delta(R_*)&=64q\log a+128q\log\log a+O(q).
\end{split}                                                    \tag{CRS25}
\]
Thus
\[
 \boxed{\displaystyle
 \mathcal T^{>R_*}=2\Delta q\int_{R_*/q}^1 f(t)\,dt
   -64q\log a-128q\log\log a+O_{\rm actual}(q).}        \tag{CRS26}
\]
The entire equilibrium integral remains in this formula.
For fixed $1<\beta\le2$ and $R=\lfloor a^\beta\rfloor$,
the same computation gives
$A_\Delta(R)=128(\beta-1)q\log a+O_\beta(q)$.
The total is
$\Delta qC_\partial-128q\log a+O(q)
=16C_\partial aq-128q\log a+O(q)$, consistent with LRS24.

\subsection{Independent proper-source displacement}

Apply the proved profile at the independent original degree $k$,
its own centre $k/2-4$, order $s_0=1,k$, conductor and guards.
Put $q_k=(k+1)^2$, $b=8k+24$, $\Delta_k=16k-48$.
Here $c_r,E_c$ mean its degree-$k$ quantities through row $q_k+2b$.
ATA35--37 prove for the actual attained quotient $g(N)=\log\det Q_N$
that $0\le g(q_k+r)-g(q_k+r-1)\le\gamma_{r,s_0}$.
The exact four-cutoff displacement expands into negative intervals
$[0,b-1]$, $[1,b]$ and positive intervals
$[q_k,q_k+2b-1]$, $[q_k+1,q_k+2b]$. Keep each occurrence:
\[
\begin{split}
 L_k&=\sum_{r=0}^{b-1}c_r+\sum_{r=1}^{b}c_r,\qquad
 H_k=\sum_{r=q_k}^{q_k+2b-1}c_r+
                               \sum_{r=q_k+1}^{q_k+2b}c_r,\\
 -L_k-2E_c&\le\delta V_Q\le H_k+2E_c .
\end{split}                                                    \tag{CRS27}
\]
For the lower bound discard the nonnegative high increments and
bound the low increments by the native rows; CRS17 costs at most
$2E_c$. For the upper bound discard the low increments and apply
the same argument to the two high intervals. There is no additional
four-copy coefficient.

The low sum is the doubled trapezoid on $[0,b]$. CRS20--21 yield
$L_k=2\Delta_kq_k\mathcal J(b/q_k)-A_{\Delta_k}(b)+O(k\log k)$.
CRS6--7 give
$\mathcal J(t)=\tfrac t2(1-\log t)+O(t^{3/2})$ at zero.
Also $A_{\Delta_k}(b)=O(k^2)$: put $b=\Delta_kx$ in CRS19;
the $\log\Delta_k$ coefficients cancel and $x$ stays in a compact
positive interval. Hence
\[
\begin{split}
 L_k&=\Delta_kb[1+\log(q_k/b)]
        +O(\Delta_kb^{3/2}/\sqrt{q_k}+k^2)\\
    &=128q_k\log k+O(q_k).
\end{split}                                                    \tag{CRS28}
\]
Here $\Delta_kb=128k^2-1152=128q_k+O(k)$ and
$\log(q_k/b)=\log k-\log8+O(1/k)$.

For $q_k\le r\le q_k+2b$, CRS6 and CRS13 give
\[
 0\le\Delta_k f(1)-c_r
       \le2\Delta_k b/q_k+\Delta_k^2/(4q_k).
\]
Indeed $0\le f(1)-f(r/q_k)\le\log(r/q_k)\le2b/q_k$.
The subtracted correction equals
$\tfrac12\int_0^{\Delta_k}\log(1+u/r)\,du
\le\Delta_k^2/(4r)$. There are $4b$ high occurrences, so
\[
 H_k=4b\Delta_k f(1)+O(k)=512q_k f(1)+O(k).            \tag{CRS29}
\]
Combining CRS27--29 and $E_c=O_{\rm actual}(q_k)$ gives
\[
 \boxed{-128q_k\log k-O_{\rm actual}(q_k)
       \le\delta V_Q\le512q_k f(1)+O_{\rm actual}(q_k).}\tag{CRS30}
\]
This calculates the displacement allowances on the stated original
intervals. The independent proper-source standard-window determinant
remains in its original metric and is not evaluated by that displacement.

CRS25 distributes the existing native $-128q\log a$ between prefix
and tail. It still cancels the literal lower-root $+128q\log a$
in LRS22. Thus the existing complete-action identity remains
$J_a-\delta^{\rm ar}_{a,1}
=C_Bq^2+8q\log a+O_{\rm actual,h}(q)$.
The signed arithmetic correction, directional allocation exponents
and projected arithmetic responses acquire no assigned values from
this scalar calculation.
\endgroup



% END PRESERVED INPUT 04_CORRELATED_PROFILE_SCALAR_PROOF.tex
\endgroup

\clearpage
\begingroup
\def\CC{\mathbb C}
\def\rank{\operatorname{rank}}
\def\coker{\operatorname{coker}}
\def\im{\operatorname{im}}
\def\tr{\operatorname{tr}}
\def\Sym{\operatorname{Sym}}
\def\diag{\operatorname{diag}}
\def\SF{\operatorname{SF}}

\setcounter{equation}{0}
\renewcommand{\theequation}{A02MAR.\arabic{equation}}
\section*{Complete marked arithmetic recovery and symmetric-power frame}
\addcontentsline{toc}{section}{Complete marked arithmetic recovery and symmetric-power frame}
% BEGIN PRESERVED INPUT 05_MARKED_ARITHMETIC_RECOVERY.tex SHA256 8c966281fbeb5784e53b651160157ca591f8e310d1c31eb312991fd3e2aebb0b



This receiving proof integrates the two new marked-recovery transcripts
and the marked arithmetic package \cite{a02mar:arrival}. It preserves the sealed
MII1--34 proof \cite{a02mar:mii} and adds the full scalar equation, its complete
initial-value recovery, and the exact second fundamental map of the
derivative-generated module. A polynomial right inverse for the enlarged
frame is constructed below. The new simultaneous-mark comparison retains
the original full-unit inclusion and computes its analytic correction.
The arithmetic scalar is received through exact metric identities; its
analytic value and remainder belong to the separately supplied signed-value
proof, and are not reproved by a determinant change of coordinates.

\section{Full primary polynomial and literal coefficient maps}
Fix the original $u\ne0$, the quartet parameters
$\rho=1/2+\delta+i\gamma$, $0<\delta<1/2$, $\gamma>2$, the
complete multiplicity $m_0\ge1$, and the original admitted odd $k$.
Set
\begin{equation}\tag{MAR1}\label{a02mar:mar:objects}
 \begin{gathered}
 e=1+k(m_0-1),\qquad q=e(k+1)^2,\qquad c=k/2,\\
 \beta_{ab}=c+(2a-k)\delta+i(2b-k)\gamma,\qquad
 \chi(S)=\prod_{a,b=0}^{k}(S-\beta_{ab})^e,\qquad\lambda=k\rho,\\
 \chi_\lambda=\chi/(S-\lambda)^e,\quad
 b_\lambda=\chi_\lambda(\lambda)^{-1},\quad
 p_\lambda=\chi/(S-\lambda),\quad v=b_\lambda p_\lambda.
 \end{gathered}
\end{equation}
Every centre is retained with its full power $e$. In $E=\CC[S]/\chi$,
the polynomial $v$ is precisely the terminal class
$e_\lambda(S-\lambda)^{e-1}$: modulo $(S-\lambda)^e$ it is
$(S-\lambda)^{e-1}$, because
$\chi_\lambda(S)/\chi_\lambda(\lambda)=1+O(S-\lambda)$;
at every other centre it is zero to its entire multiplicity.
Thus $A v=\lambda v$ for the arithmetic multiplication matrix $A$.
When $e=1$, and only then in this simple-root formula,
$b_\lambda=1/\chi'(\lambda)$. The coefficient $b_\lambda$ does not
replace the original arithmetic unit $v_h=2\xi/h$ in any inclusion or metric.

Put $R=\CC[t]$ and
\begin{equation}\tag{MAR2}\label{a02mar:mar:module}
 \mathscr H=R[S]/(u\partial_S+\chi-t)(R[S]),\qquad
 \nabla=\partial_t-A(t)/u,
 \quad A(t)=A+tB,\quad B=\mathbf e_0\mathbf e_{q-1}^{\mathsf T}.
\end{equation}
The superscript $\mathsf T$ is the ordinary coefficient transpose.
Subtracting $(u\partial_S+\chi-t)$ applied to a leading monomial
reduces every polynomial uniquely to degree less than $q$. Uniqueness
holds because a nonzero polynomial $f$ has
$\deg((u\partial_S+\chi-t)f)=q+\deg f$. Hence
$\mathscr H=R\{1,S,\ldots,S^{q-1}\}$ as an $R$-module.
The cochain commutator
$[\partial_t-S/u,u\partial_S+\chi-t]=0$ proves the connection.
On the chosen representatives, multiplication by $S$ requires only
the top relation $\chi-t$, giving exactly the companion matrix in
\eqref{a02mar:mar:module}.

Let $R_j$ be the representative of $S^jp_\lambda$. Multiplying
$(S-\lambda)p_\lambda=\chi$ by $S^{j-1}$ and using the differential
relation proves
\begin{equation}\tag{MAR3}\label{a02mar:mar:rec}
 R_j-\lambda R_{j-1}=tS^{j-1}-u(j-1)S^{j-2}\quad(1\le j\le q),
 \qquad v_j(t):=\nabla^jv=b_\lambda(-u)^{-j}R_j.
\end{equation}
The final term is zero at $j=1$. The right side has degree below $q$.
The jet formula follows by applying the cochain connection to the fixed
polynomial before reduction; the commutator proves compatibility with
reduction. No step in this argument requires $e=1$.

\section{The complete scalar equation and all its initial data}
Let $\Phi'=\chi$, $\Phi(0)=0$, and retain one of the original decay
cycles. Write
\[
 Z(t)=\int_\Gamma e^{(\Phi(S)-tS)/u}\,dS,\quad
 Y(t)=\int_\Gamma v(S)e^{(\Phi(S)-tS)/u}\,dS,
 \qquad D=-u\frac{d}{dt}.
\]
All coefficient differentiations on compact $t$ sets are dominated by
the original leading decay. Differentiation inserts $S$ and integration
by parts gives
\begin{equation}\tag{MAR4}\label{a02mar:mar:scalar-maps}
 [\chi(D)-t]Z=0,\qquad
 Y=b_\lambda p_\lambda(D)Z,\qquad
 (D-\lambda)Y=b_\lambda t Z.
\end{equation}
Consequently the maps between solution spaces on the punctured mark line
are the explicit inverse pair
\begin{equation}\tag{MAR5}\label{a02mar:mar:inverse-solutions}
 Z\longmapsto b_\lambda p_\lambda(D)Z,
 \qquad Y\longmapsto (b_\lambda t)^{-1}(D-\lambda)Y,
\end{equation}
and the scalar equation for $Y$ is
\begin{equation}\tag{MAR6}\label{a02mar:mar:Yeq}
 \bigl[p_\lambda(D)t^{-1}(D-\lambda)-1\bigr]Y=0.
\end{equation}
For the converse, define $Z$ by the second map for a solution of
\eqref{a02mar:mar:Yeq}. That equation gives $b_\lambda p_\lambda(D)Z=Y$;
applying $D-\lambda$ gives $(\chi(D)-t)Z=0$. This proves both
compositions, not merely a map in one direction.

Write $p_\lambda(S)=\sum_{j=0}^{q-1}p_jS^j$, $p_{q-1}=1$.
The Leibniz rule and $[D,t]=-u$ give
$D^jt^{-1}=\sum_{a=0}^j j!u^a t^{-a-1}D^{j-a}/(j-a)!$.
Thus all coefficients of the polynomial scalar operator are specified by
\begin{equation}\tag{MAR7}\label{a02mar:mar:P}
 \mathscr P_\lambda=
 \sum_{j=0}^{q-1}p_j\sum_{a=0}^j
 \frac{j!}{(j-a)!}u^a t^{q-a-1}D^{j-a}(D-\lambda)-t^q.
\end{equation}
Its leading derivative coefficient is $(-u)^q t^{q-1}$.
After dividing by it, the coefficient of derivative order $m$ has
a pole of order at most $q-m$: in the first term $m=j-a+1$,
so $a\le q-m$; the factor $-\lambda$ gives the stronger bound
$a\le q-1-m$. The final order-zero term has no additional pole.
This proves regular singularity of this scalar presentation at zero.

\begin{theorem}[Indicial polynomial and complete recovery]\label{a02mar:mar:scalar-theorem}
The indicial polynomial and exponent set are
\begin{equation}\tag{MAR8}\label{a02mar:mar:indicial}
 I_\lambda(r)=(-u)^q r(r-2)(r-3)\cdots(r-q),
 \qquad\{0,2,3,\ldots,q\}.
\end{equation}
Every local solution of \eqref{a02mar:mar:Yeq} on a punctured disc extends
holomorphically to zero. Its local monodromy is the identity, and no
logarithmic solution is omitted. The data $Y(0),Y''(0),\ldots,Y^{(q)}(0)$
are unrestricted initial coordinates for its $q$-dimensional solution space.
\end{theorem}
\begin{proof}
Applying \eqref{a02mar:mar:P} to $t^r$, the terms with smallest exponent
$r-1$ come exactly from the leading $j=q-1$ term and the $D$ part
of $D-\lambda$. Their coefficient is
\[
 u^q\sum_{m=1}^{q}(-1)^m\frac{(q-1)!}{(m-1)!}(r)_{\underline m}.
\]
It vanishes at $r=0$. For each integer $2\le r\le q$, terms
$m>r$ vanish and
$(r)_{\underline m}/(m-1)!=m\binom rm$; the sum of these
alternating terms is the derivative at $1$ of $(1-x)^r$, hence
zero. The coefficient of $r^q$ is $(-u)^q$. These $q$ roots and
the leading coefficient prove \eqref{a02mar:mar:indicial}.

The equation for $Z$ has polynomial coefficients and constant nonzero
leading coefficient. Its companion system has a unique entire
fundamental matrix: Picard iteration converges uniformly on each compact
disc by the bound for a bounded polynomial coefficient matrix, and the
solutions continue across overlapping discs. Its determinant solves a
scalar first-order equation and is nonzero. Thus there are $q$
independent entire solutions. The first map in
\eqref{a02mar:mar:inverse-solutions} is injective, since $Y=0$ implies
$tZ=0$. It therefore supplies $q$ independent entire solutions of
\eqref{a02mar:mar:Yeq}. An order-$q$ nonsingular equation on the punctured
disc has solution dimension $q$; these entire solutions exhaust it.
This proves the assertion about monodromy and logarithms.

For complete initial recovery set $y_n=Y^{(n)}(0)$ and
$z_n=Z^{(n)}(0)$. The last equation in \eqref{a02mar:mar:scalar-maps}
and the first one give, respectively,
\begin{equation}\tag{MAR9}\label{a02mar:mar:initial}
 \begin{gathered}
 y_1=-\lambda y_0/u,\qquad
 z_n=\frac{-u y_{n+2}-\lambda y_{n+1}}{b_\lambda(n+1)}
             \quad(0\le n\le q-2),\\
 z_{q-1}=(-u)^{-(q-1)}
 \left[b_\lambda^{-1}y_0-
             \sum_{j=0}^{q-2}p_j(-u)^jz_j\right],\\
 z_{n+q}=(-u)^{-q}
 \left[nz_{n-1}-\sum_{j=0}^{q-1}\chi_j(-u)^jz_{n+j}\right]
          \quad(n\ge0),\qquad z_{-1}=0,
 \end{gathered}
\end{equation}
where $\chi(S)=\sum_{j=0}^q\chi_jS^j$, $\chi_q=1$.
The first two rows recover all initial data of $Z$ from the stated
$q$ data of $Y$. Conversely applying the first map of
\eqref{a02mar:mar:inverse-solutions} to that $Z$ gives the same $y_0$,
then the same $y_1$, and successively the same $y_2,\ldots,y_q$.
The final row recursively specifies every subsequent derivative.
\end{proof}

\section{The entire derivative-lattice defect and a polynomial recovery map}
Let
\[
 O=[v_0(t),v_1(t),\ldots,v_{q-1}(t)],\qquad
 \mathscr M=\sum_{j=0}^{q-1}Rv_j(t)\subset\mathscr H.
\]
The recurrence \eqref{a02mar:mar:rec} gives, exactly as a column computation,
\begin{equation}\tag{MAR10}\label{a02mar:mar:lattice}
 \det O=c_\lambda t^{q-1},\quad
 c_\lambda=(-1)^{(q-1)(q+2)/2}b_\lambda^q u^{-q(q-1)/2},
 \qquad\mathscr H/\mathscr M\simeq R/(t^{q-1}).
\end{equation}
For clarity, remove the column scalars and replace $R_j$ by
$R_j-\lambda R_{j-1}$ from right to left. The columns become
$p_\lambda,t,tS-u,\ldots,tS^{q-2}-u(q-2)S^{q-3}$.
The monic first column eliminates the highest degree and contributes
$(-1)^{q-1}$ in the determinant. The remaining diagonal is $t$.
Restoring the scalars proves its sign and power of $u$.
The residual presentation is
$te_0=0$, $te_j=uj e_{j-1}$, $1\le j\le q-2$.
Its explicit isomorphism is
\begin{equation}\tag{MAR11}\label{a02mar:mar:phi}
 \varphi(e_j)=\frac{j!}{(q-2)!}u^{j-q+2}t^{q-2-j}\pmod{t^{q-1}},
 \qquad\varphi^{-1}(1)=e_{q-2}.
\end{equation}
All displayed relations vanish under this map. Conversely they express
every generator in terms of $e_{q-2}$ by exactly this formula, leaving
only its relation $t^{q-1}e_{q-2}=0$.

\begin{theorem}[The exact connection-preservation map]\label{a02mar:mar:SFthm}
The second fundamental map of $\mathscr M$ is $R$-linear and surjective:
\begin{equation}\tag{MAR12}\label{a02mar:mar:SF}
 \SF:\mathscr M\longrightarrow\mathscr H/\mathscr M,
 \qquad m\longmapsto[\nabla m].
\end{equation}
It vanishes on $v_j$, $j<q-1$, and on its last basis vector its value,
under \eqref{a02mar:mar:phi}, is
\begin{equation}\tag{MAR13}\label{a02mar:mar:g}
 g(t)=b_\lambda(-u)^{1-q}
 \sum_{j=1}^{q-1}\frac{j!p_j}{(q-2)!}
                 u^{j-q+1}t^{q-1-j}\pmod{t^{q-1}}.
\end{equation}
Its constant term is $g_0=b_\lambda(-u)^{1-q}(q-1)\ne0$.
In particular,
\begin{equation}\tag{MAR14}\label{a02mar:mar:SFkernel}
 \begin{aligned}
 \ker\SF&=Rv_0+\cdots+Rv_{q-2}+t^{q-1}Rv_{q-1},\\
 \mathscr H&=\mathscr M+Rv_q=\mathscr M+\nabla\mathscr M.
 \end{aligned}
\end{equation}
\end{theorem}
\begin{proof}
The term $f'm$ in $\nabla(fm)=f'm+f\nabla m$ lies in
$\mathscr M$, proving $R$-linearity. The first $q-1$ images
vanish by definition. In the cokernel, $R_{q-1}=0$ and
$p_\lambda=0$. Equation \eqref{a02mar:mar:rec} gives
$[R_q]=[tS^{q-1}-u(q-1)S^{q-2}]=[-u p_\lambda']$:
replace $S^{q-1}$ by $-\sum_{j=0}^{q-2}p_jS^j$, then use
$tS^j=ujS^{j-1}$ in the cokernel. Multiplying by
$b_\lambda(-u)^{-q}$ and applying \eqref{a02mar:mar:phi} proves
\eqref{a02mar:mar:g}. It is a unit in $R/(t^{q-1})$ because its
constant term is nonzero. The columns $v_0,\ldots,v_{q-1}$
are an $R$-basis of $\mathscr M$ by their nonzero determinant.
The image of $\sum f_jv_j$ is therefore $f_{q-1}g(t)$;
it vanishes exactly when $t^{q-1}$ divides $f_{q-1}$. This
proves the kernel and surjectivity. The image of $v_q$ generates
the entire cokernel, proving both recovery equalities.
\end{proof}

The following strengthens this to an explicit global polynomial
right inverse of the enlarged $q$ by $(q+1)$ frame, with every
coefficient depending on the actual $p_\lambda$. Define a row
$L(t)$ on the original coefficient vector $f=\sum_{j=0}^{q-1}f_jS^j$ by
\begin{equation}\tag{MAR15}\label{a02mar:mar:liftrow}
 L(t)f=\sum_{j=0}^{q-2}(f_j-f_{q-1}p_j)
              \frac{j!}{(q-2)!}u^{j-q+2}t^{q-2-j}.
\end{equation}
This is a polynomial lift of $\varphi$ and is $R$-linear.
Take the degree-at-most-$q-2$ representative of $g(t)$ in
\eqref{a02mar:mar:g} and put
\begin{equation}\tag{MAR16}\label{a02mar:mar:polynomialinverse}
 \begin{gathered}
 h(t)=g_0^{-1}\sum_{a=0}^{q-2}
               \left(-\frac{g(t)-g_0}{g_0}\right)^a,\\
 \mathcal R(t)=
 \begin{bmatrix}
 \displaystyle\frac{\operatorname{adj}O(t)}{c_\lambda t^{q-1}}
                   \bigl(I-v_q(t)h(t)L(t)\bigr)\\[3pt]
 h(t)L(t)
 \end{bmatrix},\qquad
 [O(t)\mid v_q(t)]\mathcal R(t)=I_q.
 \end{gathered}
\end{equation}
Although a division is displayed in its formula, every entry of
$\mathcal R$ is polynomial. Indeed $h g=1\pmod{t^{q-1}}$ by
the finite geometric sum. For every original basis vector $f$,
$L(f-v_qhLf)=0\pmod{t^{q-1}}$, so $f-v_qhLf\in\mathscr M$.
It has unique polynomial $O$-coordinates. Cramer's formula for those
coordinates is precisely the upper block in \eqref{a02mar:mar:polynomialinverse};
its numerators are therefore divisible by $c_\lambda t^{q-1}$.
Multiplying the two blocks by $[O\mid v_q]$ proves the identity.
This is a specified polynomial recovery, rather than existence of a
single additional generating direction alone.

The cokernel in \eqref{a02mar:mar:lattice} has length $q-1$ and special-fibre
dimension one. It cannot receive the given connection from
$\mathscr H$, because \eqref{a02mar:mar:SF} is the explicit nonzero descent
map. Its exact first preservation kernel is
\eqref{a02mar:mar:SFkernel}; no lower primary factor has been removed.

\section{Algebraic simplicity and the entire analytic comparison}
On $\mathscr H$ put $D_{\mathscr H}=-u\nabla$. The map
$\CC[S]\to\mathscr H$, $F\mapsto[F]$, is a complex-vector-space
isomorphism, under which
\begin{equation}\tag{MAR17}\label{a02mar:mar:weyl}
 D_{\mathscr H}F=SF,\qquad
 tF=(\chi(S)+u\partial_S)F,
 \qquad[D_{\mathscr H},t]=-u.
\end{equation}
To prove bijectivity, observe that
$(\chi+u\partial_S)^nS^j$, $n\ge0,0\le j<q$, is monic of degree
$nq+j$. These polynomials form a triangular basis of $\CC[S]$
and map to the basis $t^nS^j$ of $\mathscr H$.
A nonzero submodule stable under $t,D_{\mathscr H}$ is an ideal
in $\CC[S]$ because it is stable under multiplication by $S$.
It is also stable under $u\partial_S=t-\chi(S)$.
Repeated differentiation of a nonzero polynomial yields a nonzero
constant. The ideal is therefore all of $\CC[S]$. This proves
algebraic differential simplicity. A nonzero rational differential
subspace intersects $\mathscr H$ nontrivially after clearing
denominators; the intersection is stable under the same operators.
Thus the localization over $\CC(t)$ is irreducible as well.

The entire fundamental matrix of \eqref{a02mar:mar:module} is specified by
\begin{equation}\tag{MAR18}\label{a02mar:mar:U}
 \begin{gathered}
 U'(t)=A(t)U(t)/u,\quad U(0)=I,\quad
 U(t)=\sum_{n\ge0}U_nt^n,\\
 (n+1)U_{n+1}=u^{-1}(AU_n+BU_{n-1}),\quad U_{-1}=0,
 \qquad\det U(t)=e^{qct/u}.
 \end{gathered}
\end{equation}
The series convergence and continuation follow from the same compact
Picard bound used above. Since $\tr A=qc$, $\tr B=0$, differentiating
its determinant gives the last equality and its nonvanishing.
For any specified fibre map $\Lambda_0:E\to F$, the exact analytic
horizontal extension to the trivial target connection is
\begin{equation}\tag{MAR19}\label{a02mar:mar:Lambda}
 \Lambda_h(t)=\Lambda_0U(t)^{-1},\qquad
 \frac{d}{dt}(\Lambda_h(t)f(t))=\Lambda_h(t)\nabla f(t).
\end{equation}
Its rank is the unchanged rank of $\Lambda_0$. Algebraic simplicity
forbids a nonzero algebraic horizontal map of intermediate rank,
while \eqref{a02mar:mar:Lambda} constructs its analytic realization with
all behaviour at infinity retained.

\section{The actual observed marked defect}
Use the original admitted observation $\Lambda_0$ and its original
period and output coordinates. Put
$y(t)=\Lambda_h(t)v$, $w(t)=\Lambda_h(t)1$ and
$g(t)=e^{\lambda t/u}y(t)$. Since
$A(t)v=\lambda v+b_\lambda t1$, one obtains
\begin{equation}\tag{MAR20}\label{a02mar:mar:observed}
 \begin{gathered}
 y'+\lambda y/u=-b_\lambda t w/u,\qquad
 g'=-b_\lambda t e^{\lambda t/u}w/u,\\
 g(t)=\Lambda_0v-\frac{b_\lambda}{2u}\Lambda_01\,t^2+O(t^3).
 \end{gathered}
\end{equation}
This Taylor statement is at the fixed actual degree and periods.
For the precise first order, define
$j_* =\min\{j\ge0:\Lambda_0 S^j\ne0\}$.
For a nonzero map it satisfies $j_*\le\dim\ker\Lambda_0$,
since $1,S,\ldots,S^{\dim\ker\Lambda_0}$ are independent.
For $j\le q-1$, the cochain connection applied to $1$ gives
$w^{(j)}(0)=(-u)^{-j}\Lambda_0S^j$; no polynomial reduction
is necessary in these degrees. Inserting its first nonzero
Taylor coefficient into \eqref{a02mar:mar:observed} proves
\begin{equation}\tag{MAR21}\label{a02mar:mar:firstdefect}
 \begin{gathered}
 g^{(i)}(0)=0\quad(1\le i\le j_*+1),\\
 g^{(j_*+2)}(0)
   =b_\lambda(j_*+1)(-u)^{-(j_*+1)}\Lambda_0S^{j_*}\ne0,
 \qquad j_*+2\le\dim\ker\Lambda_0+2.
 \end{gathered}
\end{equation}

Here is a direct original-unit test making that first order
exactly quadratic. In the original simple-quartet tensor
construction let $U_h=j_h(2\xi/h)$, let $P(\varpi)$ be the
original invertible one-packet period matrix, and put
$b_U=P(\varpi)U_h\ne0$. Before the specified invertible output
coordinate map, the constant-class observation is
\begin{equation}\tag{MAR22}\label{a02mar:mar:constant-survival}
 \Lambda_k1=\mathcal P_k b_U^{\otimes k},\qquad
 \mathcal P_k=\tfrac15\sum_{j=0}^4(D_5^{-j})^{\otimes k},
 \quad D_5=\diag(\zeta_5,\zeta_5^2,\zeta_5^3,\zeta_5^4).
\end{equation}
If $5\mid k$, choose any nonzero coordinate $(b_U)_a$.
The ordered all-$a$ tensor coordinate has value $(b_U)_a^k\ne0$
and character $ka=0\pmod5$, so survives the projector. Thus on
the original sequence $k=20n+5$ in its admission domain the
coefficient of $t^2$ in \eqref{a02mar:mar:observed} is nonzero.
If $b_U$ has two nonzero coordinates $a\ne b$, then for $k\ge4$
the five exponents $ka+r(b-a)$, $0\le r\le4$, cover all residues;
one ordered tensor with $r$ entries $b$ and $k-r$ entries $a$
therefore survives. If it has only one nonzero coordinate, the
exact criterion is $5\mid k$. Every unit factor is retained in
$b_U$. The arithmetic norm of the quadratic coefficient in the
actual output quotient metric $Q_N$ is exactly
$|b_\lambda|\|\Lambda_k1\|_{Q_N}/(2|u|)$; no value relative to
the final arithmetic benchmark is assigned to this norm here.

\section{The simultaneous-mark tensor inclusion and its correction}
This section uses the original simple one-packet algebra
$E_h=\CC[S]/h$, the complete nonzero unit $U_h$, and the actual
sum inclusion
$\eta_k f=U_h^{\otimes k}f(A_\Sigma)(1^{\otimes k})$.
At zero $A_\Sigma\eta_k=\eta_k A$ exactly. Let
$A_h(t)=A_h+tB_h$ be the original one-packet companion matrix,
$B_h=\mathbf e_0\mathbf e_3^{\mathsf T}$, and
$A_\Sigma(t)=\sum_{i=1}^k A_h(t)^{(i)}$.
Define the specified defect map on the original coefficient spaces by
\begin{equation}\tag{MAR23}\label{a02mar:mar:tensordefect}
 \mathfrak D_k=\left(\sum_iB_h^{(i)}\right)\eta_k-\eta_k B,
 \qquad A_\Sigma(t)\eta_k-\eta_k A(t)=t\mathfrak D_k.
\end{equation}
For the extremal terminal class, the entire original primary
coordinate of this map is
\begin{equation}\tag{MAR24}\label{a02mar:mar:Dv}
 \begin{aligned}
 \mathfrak D_kv={}&\frac{U_h(\rho)^k}{h'(\rho)}
 \sum_{j=1}^k e_\rho^{\otimes(j-1)}\otimes1
                         \otimes e_\rho^{\otimes(k-j)}
                 -b_\lambda U_h^{\otimes k},\\
 (\mathfrak D_kv)(\bar\rho,\ldots,\bar\rho)
       ={}&-b_\lambda U_h(\bar\rho)^k\ne0\quad(k\ge2).
 \end{aligned}
\end{equation}
Indeed $\eta_kv=U_h(\rho)^ke_\rho^{\otimes k}$,
$B_he_\rho=1/h'(\rho)$, and $Bv=b_\lambda1$ by their literal
highest coefficients. On the displayed reflected tuple each
term of the sum contains at least one unchanged factor
$e_\rho(\bar\rho)=0$. The other factor is nonzero because the
original $U_h$ is a unit. This proves the full nonzero defect,
including its original coefficients.

For the two marked connections,
$\nabla_\Sigma\eta_k-\eta_k\nabla=-t\mathfrak D_k/u$.
There is an exact analytic correction, not merely an unspecified
replacement. Let $U_h^{\rm mark}(t)$ solve the one-packet
fundamental equation with identity initial value and set
$U_\Sigma(t)=(U_h^{\rm mark}(t))^{\otimes k}$.
Then
\begin{equation}\tag{MAR25}\label{a02mar:mar:eta-horizontal}
 \begin{gathered}
 \eta_h(t)=U_\Sigma(t)\eta_k U(t)^{-1},\qquad
 \nabla_\Sigma\eta_h=\eta_h\nabla,\qquad\eta_h(0)=\eta_k,\\
 \eta_h'(0)=0,\qquad\eta_h''(0)=\mathfrak D_k/u,
 \qquad\eta_h(t)=\eta_k+\frac{t^2}{2u}\mathfrak D_k+O(t^3).
 \end{gathered}
\end{equation}
The differential equation follows by differentiating the three
factors. At zero its first derivative is the original arithmetic
intertwining difference divided by $u$, hence zero. Its second
derivative is \eqref{a02mar:mar:tensordefect} divided by $u$ because
the terms containing $\eta_h'(0)$ vanish. Thus the nonzero
coordinate in \eqref{a02mar:mar:Dv} also gives the exact quadratic
correction at that same coordinate. The maps $U_\Sigma,U$ are
invertible, so the corrected map has the original rank at every
finite mark. This construction concerns the specified mark line;
no algebraic extension over the whole coefficient plane is
asserted from it.

For any original Hermitian target form $H_\otimes$, use its actual
transport \[H_\otimes(t)=U_\Sigma(t)^{-*}H_\otimes U_\Sigma(t)^{-1}.\]
The exact accompanying source metric identity is
\begin{equation}\tag{MAR26}\label{a02mar:mar:tensormetric}
 \eta_h(t)^*H_\otimes(t)\eta_h(t)
    =U(t)^{-*}\eta_k^*H_\otimes\eta_k U(t)^{-1}.
\end{equation}
This follows by cancelling the identical fundamental matrices.
It retains the original pullback $\eta_k^*H_\otimes\eta_k$ and
does not assign it the value of a different attained quotient.

\section{Partial marked flags and exact arithmetic--period receivers}
Write $v_j=v_j(0)$. From \eqref{a02mar:mar:rec},
$v_1=-\lambda v_0/u$ and
\[
 v_j+(\lambda/u)v_{j-1}=b_\lambda(j-1)(-u)^{1-j}S^{j-2}
           \quad(2\le j\le q).
\]
For $j=2$ substitute the displayed $v_1$ in terms of $v_0$.
Thus for $0\le d\le q-2$ there is an explicit triangular
column map from $(v,1,S,\ldots,S^d)$ to
$(v,v_2,\ldots,v_{d+2})$. Its first diagonal is one, and its
remaining diagonals are $b_\lambda r(-u)^{-r}$,
$1\le r\le d+1$. For every original Hermitian form $G_N$,
\begin{equation}\tag{MAR27}\label{a02mar:mar:partialgram}
 \frac{\det\operatorname{Gram}_{G_N}(v,v_2,\ldots,v_{d+2})}
      {\det\operatorname{Gram}_{G_N}(v,1,S,\ldots,S^d)}
 =|b_\lambda|^{2(d+1)}((d+1)!)^2|u|^{-(d+1)(d+2)}.
\end{equation}
Both frames are independent, since $p_\lambda$ is monic of
degree $q-1$ and $d\le q-2$. The formula follows by taking
the squared modulus of that triangular determinant. It carries
any established original-source flag comparison to its exact
marked receiver without changing a period or terminal factor.

For the full frame $J=[v,v_2,\ldots,v_q]$ one has
\[
 \det J=\epsilon_q(q-1)!b_\lambda^q u^{-q(q-1)/2},\qquad
 \epsilon_q=(-1)^{(q-1)(q+2)/2}.
\]
Let $G_N[\nu]$ denote the unchanged original attained metric for
source $\nu$. At each individual original cutoff,
\begin{equation}\tag{MAR28}\label{a02mar:mar:relative}
 \frac{\det(J^*G_N[\nu_1]J)}{\det(J^*G_N[\nu_0]J)}
       =\frac{\det G_N[\nu_1]}{\det G_N[\nu_0]}.
\end{equation}
The actual signed relative marked return is consequently
\begin{equation}\tag{MAR29}\label{a02mar:mar:Z}
 \log\mathfrak Z_{k,\lambda}
 =\sum_{N=q-1,q,2q-1,2q}\epsilon_N
       \log\frac{\det G_N[\mu_k]}{\det G_N[\sigma]}
 =\delta_{k,1}^{\rm ar},\qquad(\epsilon_N)=(+,+,-,-).
\end{equation}
This exact equality identifies the scalar that receives the
separately established signed arithmetic value and its identical
error. It introduces no second occurrence of that correction.
It also proves that this relative scalar is independent of the
chosen marked period $u$ and branch, without a conditioning bound
for the individual frames.

The consecutive frame has, for $t\ne0$, the exact finite
regularization
\begin{equation}\tag{MAR30}\label{a02mar:mar:singulargram}
 |t|^{-2(q-1)}\det(O(t)^*G_N[\nu]O(t))
      =|b_\lambda|^{2q}|u|^{-q(q-1)}\det G_N[\nu].
\end{equation}
It follows directly from \eqref{a02mar:mar:lattice}. Its finite limit
at zero differs from the regularized $J$ determinant by
exactly $((q-1)!)^2$. Thus the same relative scalar is obtained
without taking logarithms of zero determinants at zero.

The complete period map must have the same coefficient convention.
Let $\Pi_w(u,0)$ be the original period matrix in the centred
increasing powers of $w=S-c$, and let
$T_{rj}=\binom jr c^{j-r}$ for $r\le j$, zero otherwise.
Then $P=\Pi_w T$ is the original $S$-coefficient period map,
$\det T=1$, and its known fully evaluated determinant is
\[
 |\det P|^2=(2\pi)^q|u|^q e^{2q\Re(C_\Phi/u)},\qquad
 C_\Phi=\Phi(c).
\]
The proof remains valid with the full multiplicities: $\chi(c+w)$
is still monic even of even degree $q$. A lower odd-phase
coefficient inserts an odd power. Ordinary division by the even
$\chi-t$ has zero multiplication trace by parity; the required
integration-by-parts correction is the derivative of the division
quotient and strictly lowers column degree, so has zero trace.
The period determinant is constant in those lower odd
coefficients. At the monomial phase the original ordered decay
rays give $F_{ja}=\zeta^{ja}-1$ and
$g_a=(q+1)^{a/(q+1)-1}e^{\pi ia/(q+1)}\Gamma(a/(q+1))$.
Here $F^*F=(q+1)(I+\mathbf1\mathbf1^*)$ and the Gauss product
\cite{a02mar:dlmf} is $(2\pi)^{q/2}/\sqrt{q+1}$. Their product gives
the displayed modulus with every phase, branch and factor
retained. No critical-root simplicity is used.

The full period-coordinate metric and marked period frame are
\begin{equation}\tag{MAR31}\label{a02mar:mar:periodmetric}
 \begin{gathered}
 Q_N^{\rm per}[\nu]=P^{-*}G_N[\nu]P^{-1},\qquad Y_\lambda=PJ,\\
 Y_\lambda^*Q_N^{\rm per}[\nu]Y_\lambda=J^*G_N[\nu]J,\\
 \det Q_N^{\rm per}[\nu]
  =(2\pi)^{-q}|u|^{-q}e^{-2q\Re(C_\Phi/u)}\det G_N[\nu].
 \end{gathered}
\end{equation}
These are direct matrix cancellations in the stated coordinates.
Equivalently, with $H_N^{\rm ar}=J^*G_NJ$ and
$H^{\rm per}=J^*P^*PJ$, the exact relative determinant and
balanced original return are
\begin{equation}\tag{MAR32}\label{a02mar:mar:periodreturn}
 \begin{gathered}
 \det\bigl((H_N^{\rm ar})^{-1/2}H^{\rm per}(H_N^{\rm ar})^{-1/2}\bigr)
 =\frac{(2\pi|u|)^q e^{2q\Re(C_\Phi/u)}}{\det G_N},\\
 \sum_N\epsilon_N\log\det\bigl((H_N^{\rm ar})^{-1/2}
                 H^{\rm per}(H_N^{\rm ar})^{-1/2}\bigr)
              =-\mathcal B_k[\nu].
 \end{gathered}
\end{equation}
The period factor cancels because it is independent of $N$ and
the original four signs sum to zero. The native matrices in
these formulas are unchanged; no absolute angular allocation or
same-class projected response is assigned by them.

The companion analytic proof \cite[SR28--SR30]{a02mar:signed} now evaluates
the exact scalar on the right of \eqref{a02mar:mar:Z} more precisely.
Receiving that proved value through \eqref{a02mar:mar:relative} gives
\begin{equation}\tag{MAR33}\label{a02mar:mar:refined-value}
 \log\mathfrak Z_{k,\lambda}
 =-(4m_0-2)C_\Gamma q+\frac{\log2}{\pi}I_hk^2
       +\frac{C_\Gamma q}{2(\log q+c_\zeta)}+R_{h,k}.
\end{equation}
Here $c_\zeta=2\gamma_{\rm E}-\log(2\pi)$ and $C_\Gamma,I_h$
are exactly the original source constants in SR4. The error
$R_{h,k}$ is literally the same error as in SR28--SR29, with all
its original finite guards. No new frame or period error occurs.
In particular its bound in SR30 is
\[
 |R_{h,k}|\le C_h q\left[
 \frac{(1+\log\log q)^2}{(\log q)^2}
             +q^{-1/200}(\log q)^2\right]=o_h(q/\log q),
\]
with the same permitted enlargement of $C_h$ as there.
Equation \eqref{a02mar:mar:refined-value} is an exact receiving
substitution into the original marked determinant. Its analytic
proof is the cited companion proof, not a consequence of
geometric irreducibility or of the frame determinant alone.

% Fragment: requires amsmath, amssymb, hyperref.
% Source received 2026-09-19, attachment 5a756bcd-ebf7-406e-ad09-902ea93a559e, section 6.
\section{The full symmetric-power coefficient connection and its finite arithmetic bridge}

\paragraph{Source and scope.}
The received continuation, section~6, supplies the odd-quintic coefficient
family, the weighted cyclic inclusion, and the proposed all-degree extension.
The Airy irreducibility input used below is Claude Sabbah and Jeng-Daw Yu,
\emph{Hodge properties of Airy moments}, arXiv:2112.13405v2 (2023),
section~6.a, Lemma~6.2, and section~6.b, Lemma~6.6:
\href{https://arxiv.org/html/2112.13405}{the primary paper}.
Its differential-Galois input is explicitly attributed there to Nicholas
Katz, \emph{On the calculation of some differential Galois groups},
Inventiones Mathematicae \textbf{87} (1987), 13--61, Theorem~4.2.7.
The coordinate maps, matrices, quotient basis, and nonzero connection
defect in this section are calculated below.  Their statements concern
the coefficient connection and its actual finite fibre maps.  No
dimension is substituted for a native arithmetic norm.

\subsection{The original coefficient phase and the two commuting connections}

Fix the original \(u\in\mathbb C^\times\), and put
\[
 \phi_{\beta,\eta}(x)=\frac{x^5}{5}+\frac{\beta x^3}{3}+\eta x,
 \qquad
 A(\beta,\eta)=\frac1{160}+\frac{\beta}{24}+\frac{\eta}{2}.
 \tag{SPB1}
\]
Over \(R=\mathbb C[\beta,\eta]\), use the two-term relative complex
\[
 R[x]\ \xrightarrow{\ D_x\ }\ R[x]\,dx,\qquad
 D_x f=(u f'+(x^4+\beta x^2+\eta)f)\,dx.
 \tag{SPB2}
\]
The quotient \(\mathcal U=\operatorname{coker}D_x\) is free on
\(e_j=[x^j\,dx]\), \(0\leq j\leq3\).  Indeed, \(D_x f\) has
degree \(\deg f+4\) and the same leading coefficient as \(f\);
successive leading-term subtraction gives a unique remainder of
degree below four.  This also proves injectivity of \(D_x\).

On both terms of the complex the coefficient operators are
\[
 \nabla_\eta=\partial_\eta+\frac{x}{u},
 \qquad
 \nabla_\beta=\partial_\beta+\frac{x^3}{3u}.
 \tag{SPB3}
\]
They commute with \(D_x\): for example,
\([\,\partial_\eta,D_x\,]=1\) and
\([\,x/u,u\partial_x\,]=-1\); for the other operator the
two contributions are \(x^2\) and \(-x^2\).
They commute with each other, so they descend to a flat connection.
In column coordinates relative to \(e_0,e_1,e_2,e_3\), they are
\(\partial_\eta+B_\eta\) and \(\partial_\beta+B_\beta\), where
\[
 B_\eta=\frac1u
 \begin{pmatrix}
 0&0&0&-\eta\\
 1&0&0&0\\
 0&1&0&-\beta\\
 0&0&1&0
 \end{pmatrix},
 \qquad
 B_\beta=\frac1{3u}
 \begin{pmatrix}
 0&-\eta&-u&\beta\eta\\
 0&0&-\eta&-2u\\
 0&-\beta&0&\beta^2-\eta\\
 1&0&-\beta&0
 \end{pmatrix}.
 \tag{SPB4}
\]
For completeness the reductions giving all their entries are
\[
 \begin{split}
 [x^4\,dx]&=-\beta e_2-\eta e_0,\\
 [x^5\,dx]&=-\beta e_3-\eta e_1-u e_0,\\
 [x^6\,dx]&=(\beta^2-\eta)e_2-2u e_1+\beta\eta e_0.
 \end{split}
 \tag{SPB5}
\]
These follow respectively from \(D_x1,D_xx,D_xx^2\).
In particular the two integration-by-parts entries \(-u\) and
\(-2u\) have been retained.  The cochain commutators above prove
\(\partial_\beta B_\eta-\partial_\eta B_\beta+
[B_\beta,B_\eta]=0\) with these signs.

The original family is
\(\mathcal V^{\mathrm{orig}}=\mathcal E^{A/u}\otimes\mathcal U\).
Consequently its actual two matrices are
\[
 B_\eta^{\mathrm{orig}}=B_\eta+\frac1{2u}I_4,\qquad
 B_\beta^{\mathrm{orig}}=B_\beta+\frac1{24u}I_4.
 \tag{SPB6}
\]
Multiplication of periods by \(e^{A/u}\), and its inverse
by \(e^{-A/u}\), specify the comparison; the additive phase
has not been discarded.

\subsection{An exact Airy restriction, including \(dx\) and the branch}

Choose one \(\nu\in\mathbb C^\times\) with \(\nu^5=u\).
On \(\beta=0\) set \(x=\nu z\), \(\eta=-\nu^4\tau\).
Then
\[
 \phi_{0,\eta}(\nu z)/u=z^5/5-\tau z,\qquad
 x^j\,dx\longmapsto \nu^{j+1}z^j\,dz.
 \tag{SPB7}
\]
This is also an exact map of the unscaled complexes (SPB2):
use \(f(x)\mapsto u f(\nu z)\) in degree zero and
\(f(x)\,dx\mapsto\nu f(\nu z)\,dz\) in degree one.
Indeed \(D_x=\nu^4(\partial_z+z^4-\tau)\) under substitution,
and \(\nu^5=u\).  Alternatively, dividing both differentials
in (SPB2) by the invertible scalar \(u\) gives the usual
twisted de Rham complex, for which ordinary substitution
is a chain isomorphism.

The pulled-back coefficient connection is
\(\nabla_\tau=\partial_\tau-z\).  Its cyclic class \([dz]\)
has successive derivatives \((-z)^j\,dz\), \(0\leq j\leq3\).
The relation \([z^4\,dz]=\tau[dz]\) therefore gives
\[
 (\nabla_\tau^4-\tau)[dz]=0.
 \tag{SPB8}
\]
Those four derivatives form a basis.  Thus the map from
\(\mathbb C[\tau]\langle\partial_\tau\rangle/
\mathbb C[\tau]\langle\partial_\tau\rangle
(\partial_\tau^4-\tau)\) sending \(1\) to \([dz]\)
is an isomorphism: reduction by the monic differential
operator gives four generators, and their images are independent.
This is precisely \(\mathrm{Ai}_4\) in the convention of
Sabbah--Yu, section~6.a.

The original phase on the same slice is
\[
 e^{A/u}=\exp\!\left(\frac1{160u}-\frac{\tau}{2\nu}\right).
 \tag{SPB9}
\]
For a period \(F=e^{1/(160u)-\tau/(2\nu)}Z\) of the original
family, the corresponding scalar equation is
\[
 \left[\left(\partial_\tau+\frac1{2\nu}\right)^4-\tau\right]F=0.
 \tag{SPB10}
\]
Equations (SPB7)--(SPB10) retain the branch, differential,
constant phase, and linear phase.

\subsection{Every symmetric power and the geometric permutation sign}

For \(k\geq0\), write
\[
 \mathcal U_k=\operatorname{Sym}^k\mathcal U,\qquad
 d_k=\binom{k+3}{3}.
 \tag{SPB11}
\]
Restriction, taking symmetric powers, and tensoring by a line
are compatible with the displayed coefficient maps.  Sabbah--Yu,
Lemma~6.6, proves irreducibility of
\(\operatorname{Sym}^k\mathrm{Ai}_4\).
Hence \(\mathcal U_k\) has no algebraic flat subbundle of rank
strictly between zero and \(d_k\): restriction of such a
subbundle and its locally free quotient to \(\beta=0\) would
be a flat subbundle of the same rank of the irreducible
Airy symmetric power.  The identical conclusion holds for
\[
 \operatorname{Sym}^k\mathcal V^{\mathrm{orig}}
 =\mathcal E^{kA/u}\otimes\mathcal U_k.
 \tag{SPB12}
\]
The inverse comparison is tensoring by \(\mathcal E^{-kA/u}\).
This argument concerns algebraic flat subbundles on the
coefficient plane; it does not replace the original arithmetic
coefficient values by the slice values.

There is a cochain realization retaining the form orientations.
Take the tensor product of \(k\) copies of (SPB2), in the ordered
variables \(x_1,\ldots,x_k\), with its Koszul differential.
Each factor splits, as an \(R\)-module complex, into its free
degree-one quotient and a contractible complex, by the unique
degree-reduction calculation following (SPB2).  Tensoring these
splittings proves that the total complex has cohomology only in
degree \(k\), with explicit identification
\[
 [f_1\,dx_1]\otimes\cdots\otimes[f_k\,dx_k]
 \longmapsto
 [f_1(x_1)\cdots f_k(x_k)\,
 dx_1\wedge\cdots\wedge dx_k].
 \tag{SPB13}
\]
It is the top relative twisted de Rham quotient for
\(\sum_{\ell=1}^k\phi_{\beta,\eta}(x_\ell)\).
A geometric permutation of the variables acts on (SPB13)
by its sign times the ordinary permutation of the tensor
factors, because it permutes the \(k\) displayed one-forms.
Thus the ordinary symmetric tensors correspond to the
sign-isotypic part of this geometric permutation action.
This gives the exact orientation convention underlying the
symmetric-power exponential realization; no factorial or
Koszul sign is absorbed into a basis convention.

\subsection{All-power connection matrices in literal orbit sums}

For \(n=(n_0,n_1,n_2,n_3)\in\mathbb Z_{\geq0}^4\), \(|n|=k\),
let \(E_n\) be the sum of all distinct ordered tensors having
\(n_j\) entries equal to \(e_j\).  These are literal orbit sums.
For a \(4\)-by-\(4\) matrix \(B\), the induced infinitesimal
tensor action is
\[
 \begin{split}
 B^{(k)}E_n={}&\left(\sum_{j=0}^3 n_jB_{jj}\right)E_n\\
 &+\sum_{\substack{0\leq j\leq3\\ n_j>0}}
 \sum_{\substack{0\leq i\leq3\\ i\ne j}}
 (n_i+1)B_{ij}E_{n-\mathbf e_j+\mathbf e_i}.
 \end{split}
 \tag{SPB14}
\]
For each target ordered tensor in the off-diagonal term
there are exactly \(n_i+1\) source positions that could have
been changed from \(j\) to \(i\).  This proves the coefficient.
The diagonal term counts the \(n_j\) original positions.
In particular the factor \(n_i+1\) in (SPB14) must not be
replaced by the divided-power or symmetric-monomial coefficient.
Equations (SPB4) and (SPB14) calculate every entry of both
connections, at every power.  The original scalar phase adds
\[
 \frac{k}{2u}I_{d_k}\quad\hbox{in the \(\eta\)-direction},\qquad
 \frac{k}{24u}I_{d_k}\quad\hbox{in the \(\beta\)-direction}.
 \tag{SPB15}
\]

\subsection{The weighted arithmetic image and its complete quotient basis}

This paragraph is on the original simple quartet.  Use its four
idempotents \(e_{00},e_{10},e_{01},e_{11}\), original nonzero units
\(U_{00},U_{10},U_{01},U_{11}\), and original roots
\(\lambda_{\epsilon\theta}\).  Their rectangular identity is
\[
 \lambda_{00}+\lambda_{11}=\lambda_{10}+\lambda_{01}.
 \tag{SPB16}
\]
For \(0\leq a,b\leq k\), the admissible integers are
\[
 n_0(a,b)=\max(0,a+b-k)\leq n\leq n_1(a,b)=\min(a,b).
 \tag{SPB17}
\]
Let \(E_{ab,n}\) denote the orbit sum with ordered multiplicities
\((k-a-b+n,a-n,b-n,n)\), and define
\[
 w_{ab,n}=U_{00}^{\,k-a-b+n}U_{10}^{\,a-n}
           U_{01}^{\,b-n}U_{11}^{\,n}.
 \tag{SPB18}
\]
Expansion over ordered tuples gives the actual inclusion
\[
 \eta_k(e_{ab})=\sum_{n=n_0(a,b)}^{n_1(a,b)}w_{ab,n}E_{ab,n}.
 \tag{SPB19}
\]
Each ordered tuple occurs once on both sides; hence no multinomial
factor multiplies \(w_{ab,n}\).
Define \(R_k\) by
\[
 R_k\left(\sum_nx_{ab,n}E_{ab,n}\right)
 =\frac{x_{ab,n_0(a,b)}}{w_{ab,n_0(a,b)}}e_{ab}
 \quad\hbox{on each block}.
 \tag{SPB20}
\]
Then \(R_k\eta_k=I_q\), \(q=(k+1)^2\).
Indeed the selected coordinate of (SPB19) is exactly the
denominator in (SPB20).

The sum action on this block has the exact eigenvalue
\[
 \Lambda_{ab}=(k-a-b)\lambda_{00}+a\lambda_{10}+b\lambda_{01}.
 \tag{SPB21}
\]
Its expression with multiplicities in (SPB17) differs from
this by \(n(\lambda_{00}-\lambda_{10}-\lambda_{01}+\lambda_{11})=0\).
Thus \(A_\Sigma\eta_k=\eta_k A\) and \(R_k A_\Sigma=A R_k\);
both maps are typed intertwiners for the original sum operator.

For every unselected \(n>n_0(a,b)\), put
\[
 (\mathcal Q_kx)_{ab,n}
 =x_{ab,n}-\frac{w_{ab,n}}{w_{ab,n_0(a,b)}}x_{ab,n_0(a,b)}.
 \tag{SPB22}
\]
Let \(\mathcal S_k\) insert these coordinates in the corresponding
unselected orbit columns and put zero in the selected columns.
Direct substitution proves all of the identities
\[
 \begin{gathered}
 \ker\mathcal Q_k=\operatorname{im}\eta_k,\qquad
 \mathcal Q_k\mathcal S_k=I,\qquad R_k\mathcal S_k=0,\\
 x=\eta_kR_kx+\mathcal S_k\mathcal Q_kx,\qquad
 [\,\eta_k\ \mathcal S_k\,]^{-1}
 =\begin{bmatrix}R_k\\ \mathcal Q_k\end{bmatrix}.
 \end{gathered}
 \tag{SPB23}
\]
Consequently the quotient basis is precisely
\([E_{ab,n}]\) for \(n>n_0(a,b)\), with the complete relation
\[
 [E_{ab,n_0}]
 =-\sum_{n>n_0}\frac{w_{ab,n}}{w_{ab,n_0}}[E_{ab,n}].
 \tag{SPB24}
\]
Counting four-part compositions and the nonempty \((a,b)\)
blocks now gives
\[
 \dim\operatorname{coker}\eta_k
 =\binom{k+3}{3}-(k+1)^2=\frac{k(k-1)(k+1)}6.
 \tag{SPB25}
\]
This proves the whole exact sequence and a splitting, rather
than only its dimensions.  Conjugating (SPB19)--(SPB24) by
the original invertible period map \(P(u)^{\otimes k}\)
gives the corresponding exact sequence in the period fibre.
All the original \(U_{\epsilon\theta}\) remain in those maps.

\subsection{The cyclic projector and its explicitly nonzero connection defect}

At the monomial coefficient point \(\beta=\eta=0\), the
geometric action \(x\mapsto\zeta_5x\) gives
\(e_j\mapsto\zeta_5^{j+1}e_j\), including the factor from \(dx\).
The cyclic projector on the symmetric fibre therefore acts by
\[
 \Pi_5 E_n=p_nE_n,\qquad
 p_n=\mathbf1_{\{\sum_{j=0}^3(j+1)n_j\equiv0\pmod5\}}.
 \tag{SPB26}
\]
For any nonidentity fifth root of unity the symmetric-character
generating function is
\[
 \prod_{a=1}^4(1-\zeta_5^a z)^{-1}=\frac{1-z}{1-z^5}.
 \tag{SPB27}
\]
Average the identity character and the four nonidentity
characters, and take the coefficient of \(z^k\).  This proves
\[
 \operatorname{rank}\Pi_5
 =\frac15\left[\binom{k+3}{3}+4\varepsilon_k\right],\qquad
 \varepsilon_k=
 \begin{cases}
 1&k\equiv0\pmod5,\\
 -1&k\equiv1\pmod5,\\
 0&k\equiv2,3,4\pmod5.
 \end{cases}
 \tag{SPB28}
\]
The exact connection defect of the constant projector in these
coefficient coordinates is
\[
 [B^{(k)},\Pi_5]_{mn}=(p_n-p_m)B^{(k)}_{mn}.
 \tag{SPB29}
\]
It can be evaluated at every \(k\geq2\), without an
irreducibility argument.  Use the occupation vector
\[
 n=
 \begin{cases}
 (k/2,0,0,k/2)& k\ \text{even},\\
 ((k+1)/2,0,1,(k-3)/2)& k\ \text{odd}.
 \end{cases}
 \tag{SPB30}
\]
Its weight is a multiple of five, so \(p_n=1\).
Set \(m=n-\mathbf e_0+\mathbf e_1\); then \(p_m=0\).
At \(\beta=\eta=0\), (SPB4) has
\((B_\eta)_{1,0}=1/u\).
Formula (SPB14), with \(n_1=0\), gives the exact nonzero entry
\[
 [B_\eta^{(k)},\Pi_5]_{mn}=\frac1u.
 \tag{SPB31}
\]
The scalar phase contribution in (SPB15) commutes with
\(\Pi_5\), so (SPB31) is unchanged for the original family.
The two vectors \(E_n,E_m\) also prove directly that the
projector has positive, proper rank for every \(k\geq2\).
For \(k=0,1\) its ranks are respectively \(1,0\).

Equation (SPB31) proves the failure of the literal constant
projector to be horizontal.  The stronger global statement
follows from the Airy restriction: an algebraic horizontal
idempotent extending this fibre projector would have an
image of its same positive proper rank, contradicting
(SPB11)--(SPB12).  The original cyclic image (SPB19) likewise
has no proper algebraic flat containing subbundle for \(k\geq2\).
For \(k=0,1\), its rank already equals \(d_k\).

The analytic comparison is explicit.  On a simply connected
analytic coefficient neighborhood, let \(F\) solve
\[
 dF=-\left(B_\beta^{(k)}\,d\beta+B_\eta^{(k)}\,d\eta
              +k\,dA/u\right)F,\qquad F(p_0)=I.
 \tag{SPB32}
\]
Flatness gives an invertible analytic solution.  For the
given fibre projector \(P_0\), the analytic projector
\(F P_0F^{-1}\) is horizontal, since
\(d(FP_0F^{-1})+[\Omega,FP_0F^{-1}]=0\).
Its inverse fibre comparison is conjugation by \(F^{-1}\).
Thus the obstruction just proved is specifically algebraic;
the analytic comparison map is retained in full.

\subsection{The metric carried by the explicit quotient}

The algebraic retraction (SPB20) specifies coordinates.
The actual period metric can also be retained exactly.
Let \(P(u)\) be the original period map on the four root
idempotents, let \(G=P(u)^*P(u)\), and take the tensor metric
on the ordered period tensors.  Its Gram matrix in the literal
orbit sums is the following explicit finite matrix:
\[
 H_{mn}=[X^mY^n]\left(\sum_{i,j=0}^3G_{ij}X_iY_j\right)^k.
 \tag{SPB33}
\]
Expanding the \(k\) factors enumerates every ordered left
tuple of multiplicity \(m\) and every ordered right tuple of
multiplicity \(n\), once.  Their products of \(G_{ij}\) are
exactly their tensor inner products.  This proves (SPB33)
with all tuple multiplicities.
In the explicit frame \(T=[\,\eta_k\ \mathcal S_k\,]\), put
\[
 T^*HT=\begin{pmatrix}H_{11}&H_{12}\\H_{21}&H_{22}\end{pmatrix}.
 \tag{SPB34}
\]
The restriction metric is
\(H_{11}=\eta_k^*H\eta_k\), and the attained quotient metric
on the actual coordinates (SPB22) is exactly
\[
 H_{\mathrm{quot}}
 =H_{22}-H_{21}H_{11}^{-1}H_{12}.
 \tag{SPB35}
\]
Indeed every lift of a quotient coordinate \(y\) is
\(\mathcal S_ky+\eta_kz\).  Completing the square in its
actual norm gives the unique minimizing coordinate
\(z=-H_{11}^{-1}H_{12}y\), and substituting it gives (SPB35).
This proves the metric receiving map as well as its
algebraic quotient, while retaining the off-diagonal blocks.
The source columns in (SPB19) are sum-centre idempotents,
whereas the original native matrix is in increasing
\(S\)-coefficients.  Retain the actual polynomial
\(\chi_k(S)=\prod_{a,b=0}^k(S-\Lambda_{ab})\), and put
\[
 \begin{gathered}
 e_{ab}(S)=
 \prod_{\substack{0\leq a',b'\leq k\\(a',b')\ne(a,b)}}
 \frac{S-\Lambda_{a'b'}}{\Lambda_{ab}-\Lambda_{a'b'}},
 \qquad
 C_k=[\,\operatorname{coeff}_{1,S,\ldots,S^{q-1}}(e_{ab})\,]_{(a,b)},\\
 (C_k^{-1})_{(a,b),j}=\Lambda_{ab}^{\,j},\quad 0\leq j<q,
 \qquad
 \eta_k^{(S)}=\eta_k^{(\mathrm{id})}C_k^{-1},\qquad
 G_N^{(\mathrm{id})}=C_k^*G_N^{(S)}C_k .
 \end{gathered}
 \tag{SPB35a}
\]
Here one fixed ordering of \((a,b)\) is used in every
idempotent-coordinate matrix.  On the original quartet,
\(\Lambda_{ab}=k/2+(2a-k)\delta+i(2b-k)\gamma\), with the
original real nonzero \(\delta,\gamma\); equality of two
centres forces equality of both indices.  Evaluating the displayed
Lagrange polynomials at these centres proves \(C_k^{-1}C_k=I\),
with every root-difference factor retained.
The matrix \(\eta_k\) in (SPB19)--(SPB35) is precisely
\(\eta_k^{(\mathrm{id})}\).
Thus the period-to-native relative form in those same source
coordinates is
\[
 \bigl(G_N^{(\mathrm{id})}\bigr)^{-1/2}
 \bigl((\eta_k^{(\mathrm{id})})^*H\eta_k^{(\mathrm{id})}\bigr)
 \bigl(G_N^{(\mathrm{id})}\bigr)^{-1/2}.
 \tag{SPB36}
\]
In the original increasing-\(S\) coordinates, the identical
pair of forms is represented by
\(G_N^{(S)}\) and
\((\eta_k^{(S)})^*H\eta_k^{(S)}\).
Both pairs are related by the displayed congruence \(C_k\);
their relative operators are similar, because
\[
 (G_N^{(\mathrm{id})})^{-1}
   (\eta_k^{(\mathrm{id})})^*H\eta_k^{(\mathrm{id})}
 =
 C_k^{-1}(G_N^{(S)})^{-1}
   (\eta_k^{(S)})^*H\eta_k^{(S)}C_k .
\]
Neither (SPB28) nor (SPB25) replaces any entry of (SPB35)
or (SPB36).  These formulas preserve the exact place where
the original arithmetic metric enters the same finite maps.


\section{Receiving scope and exact earlier updates}
The scalar equation MII30 is extended to the entire explicit
coefficient operator MAR7, the indicial polynomial MAR8, and the
complete initial recovery MAR9. The consecutive-frame cokernel
MII19--21 is extended to its surjective second fundamental map,
its exact preservation kernel, and a polynomial right inverse
in MAR12--16. The formal matrix in MII25 is identified with the
entire analytic fundamental solution MAR18, and the observed
and tensor comparisons are the exact maps MAR19--26.
The general full-frame receiver MII31--34 is extended to every
partial marked flag by MAR27 and to the two original source
metrics and period coordinates by MAR28--32. These updates are
additive; the sealed earlier proof remains an exact source.

The ZIP's analytic central-response estimates and the transcripts'
signed arithmetic asymptotics are evaluated in their designated
analytic lanes. The present proof supplies their exact marked
receivers, including all factors. It assigns no part of the
relative scalar to the seven individual native returns. The
explicit $t$-dependent maps also do not assert preservation of
the original coefficient-cone subspaces under a different
connection; their defects and their actual maps are displayed.



% END PRESERVED INPUT 05_MARKED_ARITHMETIC_RECOVERY.tex
\endgroup

\clearpage
\begingroup
\def\C{\mathbb C}
\def\R{\mathbb R}
\def\PP{\mathcal P}
\def\Fcal{\mathscr F}
\def\Bcal{\mathcal B}
\def\norm#1{\left\lVert#1\right\rVert}
\def\Tr{\operatorname{Tr}}
\def\im{\operatorname{im}}
\def\Gram{\operatorname{Gram}}
\def\cond{\operatorname{cond}}

\setcounter{equation}{0}
\renewcommand{\theequation}{A02GMF.\arabic{equation}}
\section*{Growing marked flags and fixed-degree cumulant determinant}
\addcontentsline{toc}{section}{Growing marked flags and fixed-degree cumulant determinant}
% BEGIN PRESERVED INPUT 06_GROWING_MARKED_FLAGS.tex SHA256 5c6a510f304d610e782893e2efa9ce03277c73eee0180392dc8e380b957029c0



\begin{abstract}
For the original arithmetic source and its specified centre-filled
comparison, a linearly growing marked flag has large positive
individual logarithmic volume changes at every original cutoff.
Its balanced four-cutoff return nevertheless has an explicit
exponentially small bound.  The proof estimates every admissible
relation and terminal representative before taking the original
minima; it retains the full terminal--polynomial cross block.
The same estimates give linearly many large generalized
eigenvalues and a complementary subspace forcing all but
\(O_h(q/\log(q/k))\) eigenvalues into an exponentially short
interval about one.  Exact quotient determinant factorization
then carries the complete signed central response into the
attained complement of the marked flag.  The stated complement
is retained as its actual metric quotient.
\end{abstract}

\section{Original source, complete roots and fixed minimum maps}

Keep the complete hypothetical actual quartet and its full order:
\begin{equation}\tag{GMF1}
\begin{gathered}
 \rho=\tfrac12+\delta+i\gamma,\quad
 0<\delta<\tfrac12,\quad \gamma>2,\quad m_0\ge1,\\
 h(s)=\prod_{\zeta\in\{\rho,\bar\rho,1-\rho,1-\bar\rho\}}
                       (s-\zeta)^{m_0},\quad
 v_h=2\xi/h,\quad
 w_h(y)=|v_h(1/2+iy)|^2/(2\pi),\quad m_k=w_h^{*k},\\
 k\equiv1\pmod4,\ k\ge5,\quad
 e=1+k(m_0-1),\quad q=e(k+1)^2,\quad c=k/2,\\
 \chi(S)=\prod_{a,b=0}^k
 [S-c-(2a-k)\delta-i(2b-k)\gamma]^e,\qquad
 E=\C[S]/(\chi).
\end{gathered}
\end{equation}
The arithmetic source is the unchanged \(k\)-fold convolution,
with no numerical off-critical zero substituted for this original
counterfactual.  All limits fix this complete packet.

Write
\begin{equation}\tag{GMF2}
\begin{gathered}
 \alpha=\pi/2,\qquad
 M(\theta)=\int_\R e^{\theta y}w_h(y)\,dy,\quad
 M_0=M(0),\quad M_\alpha=M(\alpha),\\
 b=(\log M)'(\alpha)>0,\quad T=kb,\qquad
 \sigma(y)=|\Gamma(1/4+iy/2)|^2/(2\pi),\quad
 M_\sigma=\int_\R\sigma(y)\,dy=\sqrt{2\pi},\\
 \widetilde m_k(y)=
 \begin{cases}M_\alpha^k\sigma(y),&|y|\le T,\\
 m_k(y),&|y|>T.\end{cases}
\end{gathered}
\end{equation}
Both complete source masses are retained in every comparison.
This is precisely the central comparison in CPT4:
\(m_k(y)dy=M_\alpha^k\,d\nu_{1,1}\) and
\(\widetilde m_k(y)dy=M_\alpha^k\,d\nu_{0,1}\)
\cite{a02gmf:cpt}.  These equalities identify the two source
presentations without changing their norms.

For a positive source density \(a(y)\), let \(H_N[a]\) be the
Gram of \(1,S,\ldots,S^N\) in the original coordinate \(S=c+iy\).
Its inner product is conjugate-linear in the first variable.
For \(N\ge q-1\), let
\begin{equation}\tag{GMF3}
\begin{gathered}
 J_N:\C[S]_{\le N}\twoheadrightarrow E,\quad
                  J_NP=[P]_\chi,\qquad
 \ker J_N=\chi\C[S]_{\le N-q},\\
 G_N[a]=(J_NH_N[a]^{-1}J_N^*)^{-1},\quad
 R_N[a]=H_N[a]^{-1}J_N^*G_N[a].
\end{gathered}
\end{equation}
A negative degree means the zero space.  Multiplication gives
\(J_NR_N=I_E\), and \(R_Nv\) is orthogonal to every vector in
\(\ker J_N\).  Every representative is \(R_Nv+z\), so its squared
norm is \(v^*G_Nv+\norm z_a^2\).
This proves attainment of the original minimum and the exact
formula (GMF3).

For all claims about cutoff return take
\[
 N\in\{q-1,q,2q-1,2q\},
\]
and the proofs below in fact hold for every integer
\(q-1\le N\le2q\).
For a function \(F_N\) of these cutoffs, retain the signs
\[
 \mathcal R_{\rm ar}F_N=F_{q-1}+F_q-F_{2q-1}-F_{2q},\qquad
 \Bcal_k[a]=\mathcal R_{\rm ar}\log\det G_N[a].
\]
No cutoff is identified with another.

Fix an original root \(\lambda\) and set
\begin{equation}\tag{GMF4}
 \chi_\lambda=\frac{\chi}{(S-\lambda)^e},\qquad
 b_\lambda=\chi_\lambda(\lambda)^{-1},\qquad
 v_\lambda=\left[b_\lambda\frac{\chi}{S-\lambda}\right]_\chi .
\end{equation}
The complete-primary coefficient is nonzero.
It equals \(1/\chi'(\lambda)\) only when \(e=1\).
For \(0\le d\le q-2\), define the fixed flag and its displayed frame
\begin{equation}\tag{GMF5}
 \Fcal_{\lambda,d}
 =\operatorname{span}(v_\lambda,[1]_\chi,\ldots,[S^d]_\chi),
 \qquad F_{\lambda,d}=(v_\lambda,1,S,\ldots,S^d).
\end{equation}
The dimension is \(d+2\): the unique representative of
\(v_\lambda\) has degree \(q-1>d\) with leading coefficient
\(b_\lambda\ne0\), and a degree-\(<q\) representative is zero
in \(E\) only if it is the zero polynomial.

\section{Exact low-polynomial evaluation with the original Gamma mass}

Use the fixed auxiliary constants
\begin{equation}\tag{GMF6}
 \theta_*=\pi/4,\qquad
 r_*=\tanh(\pi/8),\qquad
 \ell_*=-\log r_*=\log\coth(\pi/8),\qquad
 \omega_*=\alpha-\theta_*=\pi/4.
\end{equation}
They are bounds on the original source, with no new source order.
The monic Gamma polynomials and their full norms are
\[
 \sum_{n\ge0}\frac{p_n(y)}{n!}z^n
 =(1+z^2)^{-1/4}e^{y\arctan z},\qquad
 \gamma_n=M_\sigma n!(1/2)_n.
\]
The generating function is the original
Meixner--Pollaczek specialization
\[p_n(y)=n!P_n^{(1/4)}(y/2;\pi/2).\]
Substitution in the generating formula
\cite[(18.23.7)]{a02gmf:dlmfMPgen} gives the displayed expression:
\[
 (1-iz)^{-1/4+iy/2}(1+iz)^{-1/4-iy/2}
       =(1+z^2)^{-1/4}e^{y\arctan z}.
\]
The analytic logarithms are the branches equal to zero at \(z=0\).
The Gamma norm is the original GTM7/CPT8 norm
\cite{a02gmf:providers,a02gmf:cpt}, since
\((1/2)_n=(2n)!/(4^n n!)\).

For \(|z|=r_*<1\), the power series for \(\arctan z\) gives
\(|\arctan z|\le\operatorname{artanh}r_*=\pi/8\), and
\(|1+z^2|\ge1-r_*^2\).
Cauchy's coefficient formula therefore gives
\[
 \frac{|p_n(y)|}{n!}
 \le r_*^{-n}(1-r_*^2)^{-1/4}
                      e^{|y|\operatorname{artanh}r_*}.
\]
The elementary central-binomial estimate
\[
 \frac{n!}{(1/2)_n}=\frac{4^n}{\binom{2n}{n}}\le2n+1
\]
follows because the largest of the \(2n+1\) binomial coefficients
is at least their average.
Thus the complete degree-\(d\) evaluation kernel obeys
\begin{equation}\tag{GMF7}
 K_d^\sigma(y,y)=\sum_{n=0}^d\frac{|p_n(y)|^2}{\gamma_n}
 \le A_de^{\theta_*|y|},\qquad
 A_d=\frac{(d+1)^2 e^{2d\ell_*}}
               {M_\sigma\sqrt{1-r_*^2}}.
\end{equation}
The factor \((d+1)^2\) is the exact sum
\(\sum_{n=0}^d(2n+1)\); every coefficient is included.
Expanding any polynomial in the orthogonal basis and applying
Cauchy--Schwarz proves
\[
 |P(c+iy)|^2\le K_d^\sigma(y,y)\norm P_\sigma^2
 \quad(P\in\C[S]_{\le d}).
\]
The coefficient map \(P(S)\mapsto P(c+iy)\) is invertible,
with determinant \(i^{d(d+1)/2}\) of modulus one, so this
estimate and all source determinant identities concern the
original coefficient space.

Evenness and the exact convolution transform give
\[
 \int_\R e^{\theta|y|}m_k(y)dy
 \le \int(e^{\theta y}+e^{-\theta y})m_k(y)dy
 =2M(\theta)^k,\qquad 0\le\theta\le\alpha.
\]
Combine this identity with (GMF7), first on the entire line
and then on the exact complement of \([-Y,Y]\), to obtain
\begin{equation}\tag{GMF8}
\begin{gathered}
 \norm P_{m_k}^2
 \le2A_dM(\theta_*)^k\norm P_\sigma^2,\\
 \int_{|y|>Y}|P(c+iy)|^2m_k(y)dy
 \le2A_dM_\alpha^k e^{-\omega_*Y}\norm P_\sigma^2
 \quad(Y>0).
\end{gathered}
\end{equation}
For the second inequality use
\(e^{\theta_*|y|}\le e^{-\omega_*Y}e^{\alpha|y|}\)
on the retained exterior.

Let \(C_\Gamma=\sqrt{2\sqrt5}e^{2/3}\) be the original upper
Gamma constant in RTM2/CPT3 \cite{a02gmf:providers,a02gmf:cpt}, and set
\begin{equation}\tag{GMF9}
 \eta_{k,d}=
 \frac{2C_\Gamma A_d}{\omega_*\sqrt{1+T}}e^{-\omega_*T},
 \qquad
 \zeta_{k,d}=2A_de^{-\omega_*T}.
\end{equation}
Indeed
\[
 \int_{|y|>T}|P(c+iy)|^2d\sigma
 \le2C_\Gamma A_d\norm P_\sigma^2
           \int_T^\infty\frac{e^{-\omega_*y}}{\sqrt{1+y}}dy
 \le\eta_{k,d}\norm P_\sigma^2.
\]
For the centre-filled norm keep both of its real regions:
\begin{equation}\tag{GMF10}
 M_\alpha^k(1-\eta_{k,d})\norm P_\sigma^2
 \le\norm P_{\widetilde m_k}^2
 \le M_\alpha^k(1+\zeta_{k,d})\norm P_\sigma^2.
\end{equation}
The lower bound retains the full central Gamma integral.
The upper bound retains its full-line upper bound and the
arithmetic exterior estimate (GMF8) at \(Y=T\).

Put
\begin{equation}\tag{GMF11}
 a_h=\log\frac{M_\alpha}{M(\theta_*)}>0,\qquad
 d_k=\left\lfloor\frac{a_h k}{4\ell_*}\right\rfloor.
\end{equation}
Strict positivity and strict convexity of \(\log M\) follow
by differentiating the actual source integral: its second
derivative is the variance of the positive tilted density
divided by its full mass, and that density is not supported
at a point.  Evenness gives \((\log M)'(0)=0\).
Consequently
\[
 0<a_h=\int_{\theta_*}^{\alpha}(\log M)'(\theta)d\theta
       <\omega_*b.
\]
For \(0\le d\le d_k\), \(2d\ell_*\le a_hk/2\) and
\(\log A_d\le a_hk/2+O_h(\log k)\).
Equations (GMF9) therefore prove, uniformly in these degrees,
\begin{equation}\tag{GMF12}
 \eta_{k,d},\ \zeta_{k,d}
 \le\exp[-a_hk/2+O_h(\log k)].
\end{equation}
The finite inequalities (GMF7)--(GMF10) remain the bounds used
at each stated \(k\); eventual notation has not replaced a guard.

\section{Uniform localization of every high-fibre polynomial}

Write \(B_h=42+8m_0\) and
\(\vartheta_h=\int_{-1}^1w_h(y)dy>0\).
The original threefold convolution lower bound is
\[
 w_h^{*3}(y)\ge c_h e^{-\alpha|y|}(1+|y|)^{-B_h}.
\]
In its convolution with the remaining \(k-3\) factors,
restrict every additional variable to \([-1,1]\).
The sum of those variables has modulus at most \(k-3\);
their retained product mass is \(\vartheta_h^{k-3}\).
The triangle inequality in the exponential and the polynomial
factor therefore proves on the entire real line
\begin{equation}\tag{GMF13}
 m_k(y)\ge c_h\vartheta_h^{k-3}
 e^{-\alpha(|y|+k-3)}(k-2+|y|)^{-B_h}.
\end{equation}
This derivation retains the original numerator and its
convolution mass.

Set
\begin{equation}\tag{GMF14}
\begin{gathered}
 R_k=k\sqrt{\delta^2+\gamma^2},\quad Y=\sqrt{kq},\\
 T\le Y,\qquad Y\le q/2,\qquad Y+2R_k<q,\\
 \underline m_{k,q}
 =c_h\vartheta_h^{k-3}e^{-\alpha(2q+k-3)}
                        (k-2+2q)^{-B_h},\\
 \rho_Y=\frac{Y+R_k}{q-R_k}<1,\qquad
 \mathcal M_k=2M_0^k+M_\alpha^kM_\sigma.
\end{gathered}
\end{equation}
The three displayed inequalities are explicit finite guards.
They hold eventually at each fixed packet, since \(q\ge(k+1)^2\).
The number \(\underline m_{k,q}\) is a lower bound for the
actual density at every point of \([q,2q]\).
Also
\(\int(m_k+\widetilde m_k)\le\mathcal M_k\), by (GMF2) and
\(\int m_k=M_0^k\).

Let \(\Psi(y)\) be monic of degree \(q\) or \(q-1\), with
every root, counted with multiplicity, of modulus at most \(R_k\).
For every \(F=\Psi p\) of degree at most \(2q\), define
\begin{equation}\tag{GMF15}
 \beta_k=
 \frac{\mathcal M_k}{\underline m_{k,q}}
 \frac{(q+2)^2}{q}\,9^{2(q+1)}
                     \rho_Y^{2(q-1)}.
\end{equation}
Then
\begin{equation}\tag{GMF16}
 \int_{|y|\le Y}|F(y)|^2(m_k(y)+\widetilde m_k(y))dy
 \le\beta_k\norm F_{m_k}^2,\qquad
 (1-\beta_k)\norm F_{m_k}^2
 \le\norm F_{\widetilde m_k}^2
 \le(1+\beta_k)\norm F_{m_k}^2.
\end{equation}

To prove the full coefficient estimate, \(\deg p\le q+1\).
The shifted orthonormal Legendre polynomials on \([q,2q]\)
are
\[
 q^{-1/2}\sqrt{2j+1}\,P_j(2y/q-3),\qquad 0\le j\le q+1.
\]
Use the orthogonality and recurrence in
\cite[Tables 18.3.1 and 18.9.1]{a02gmf:dlmfLegendre}.
For \(|x|\le4\), their recurrence, \(P_0=1,P_1=x\), gives
\(|P_j(x)|\le9^j\), because
\[
 8\cdot9^j+9^{j-1}<9^{j+1}.
\]
Under \(Y\le q/2\), \(|2y/q-3|\le4\) on \(|y|\le Y\).
Expand the complete polynomial \(p\) in this basis and
apply Cauchy--Schwarz.  Summing every evaluation square yields
\[
 \sup_{|y|\le Y}|p(y)|^2
 \le\frac{(q+2)^2}{q}9^{2(q+1)}
                         \int_q^{2q}|p(x)|^2dx.
\]
All roots retained in \(\Psi\) give
\[
 |\Psi(y)|\le(Y+R_k)^{\deg\Psi}\quad(|y|\le Y),\qquad
 |\Psi(x)|\ge(q-R_k)^{\deg\Psi}\quad(q\le x\le2q).
\]
Use the full mass \(\mathcal M_k\) for the numerator and the
actual lower density \(\underline m_{k,q}\) for the indicated
part of the denominator.  Division for \(p\ne0\) gives a
factor \(\rho_Y^{2\deg\Psi}\), at most
\(\rho_Y^{2(q-1)}\).  The zero polynomial is immediate.
This proves the first inequality (GMF16).
The measures agree for \(|y|>T\), and \(T\le Y\); hence
\[
 |\norm F_{\widetilde m_k}^2-\norm F_{m_k}^2|
 \le\int_{|y|\le Y}|F|^2(m_k+\widetilde m_k).
\]
This proves its two relative inequalities without selecting
a preferred representative.

In the original coefficient coordinate, the complete relation
polynomials have the form \(\chi P\), and every representative
of \(v_\lambda\) has the form
\begin{equation}\tag{GMF17}
 b_\lambda\frac{\chi}{S-\lambda}+\chi P
 =\frac{\chi}{S-\lambda}
                  [b_\lambda+(S-\lambda)P].
\end{equation}
The real-coordinate monic factors are
\(i^{-q}\chi(c+iy)\) and
\(i^{-(q-1)}(\chi/(S-\lambda))(c+iy)\).
Their roots have modulus at most \(R_k\).
The factors \(i^q,i^{q-1}\), and \(b_\lambda\) in (GMF17)
are retained scalars; (GMF16) applies to all coefficients and
therefore applies before both original minima.
Writing
\(E_{N,\lambda}[a]=v_\lambda^*G_N[a]v_\lambda\), it gives
\begin{equation}\tag{GMF18}
 1-\beta_k\le
 \frac{E_{N,\lambda}[\widetilde m_k]}{E_{N,\lambda}[m_k]}
 \le1+\beta_k,\qquad q-1\le N\le2q.
\end{equation}
At \(N=q-1\) the fibre consists of the unique degree-\(<q\)
representative; the same inequalities still hold.

Taking logarithms of the explicitly positive bound (GMF15),
and using
\[
 \log\rho_Y
 =-\tfrac12\log(q/k)
   +\log(1+R_k/\sqrt{kq})-\log(1-R_k/q),
\]
gives the fixed-packet estimate
\begin{equation}\tag{GMF19}
 \log\beta_k
 \le-(q-1)\log(q/k)+O_h(q+k+\log q).
\end{equation}
In particular \(\beta_k\) tends to zero faster than
\(\exp(-c q)\) for every fixed \(c>0\).
This does not remove its finite value from the following guards.

\section{Original attained low metric and the entire cross block}

The original global envelope and polynomial norm comparison are
\begin{equation}\tag{GMF20}
\begin{gathered}
 m_k(y)\ge a_k(1+y^2)^{-\beta_0}\sigma(y),\qquad
 \beta_0=21+4m_0,\\
 a_k=\frac{c_h\vartheta_h^{k-3}e^{-\alpha(k-3)}(k-2)^{-B_h}}
                 {C_\Gamma 2^{B_h/2}},\qquad
 D_d=[1+4(d+\beta_0)^2]^{\beta_0},\\
 \norm P_{m_k}^2\ge a_kD_d^{-1}\norm P_\sigma^2
                    \qquad(\deg P\le d).
\end{gathered}
\end{equation}
These are the original GTM1 and ACW25--26 bounds
\cite{a02gmf:providers,a02gmf:centralProvider}.  They also follow from
the sharper CPT7 bound
\([1+4(d+1)^2]^{-\beta_0}\), since \(\beta_0\ge1\)
\cite{a02gmf:cpt}.  Thus the incoming notation \(a_k^{\rm AMT}\)
in this calculation is fixed as this actual \(a_k\); no
independent source mass is introduced.

Use (GMF8),(GMF10),(GMF20) to define the finite low exterior
fractions
\begin{equation}\tag{GMF21}
\begin{gathered}
 \ell_m=
 2A_dM_\alpha^k\frac{D_d}{a_k}e^{-\omega_*Y},\qquad
 \ell_f=\frac{2A_d}{1-\eta_{k,d}}e^{-\omega_*Y},\\
 e_m=(\sqrt{\beta_k}+\sqrt{\ell_m})^2,\qquad
 e_f=\left(\sqrt{\frac{\beta_k}{1-\beta_k}}
                              +\sqrt{\ell_f}\right)^2,\\
 \eta_{k,d}<1/2,\quad \beta_k<1/4,\quad
 e_m<1/4,\quad e_f<1/4,\qquad
 z_m=\frac{e_m}{1-e_m},\quad z_f=\frac{e_f}{1-e_f}.
\end{gathered}
\end{equation}
The subscripts \(m,f\) denote the actual arithmetic and
centre-filled sources respectively.
Each displayed strict inequality is a finite guard.

For a low polynomial \(P\) of degree at most \(d\), (GMF21)
bounds its fraction of squared norm outside \([-Y,Y]\).
For any high polynomial covered by (GMF16), that formula
bounds its central fraction in the original source by
\(\beta_k\), and in the filled source by
\(\beta_k/(1-\beta_k)\).
The latter uses
\(\norm F_{\widetilde m_k}^2\ge(1-\beta_k)\norm F_{m_k}^2\).
Split the pairing into the two \emph{complete} regions
\(|y|\le Y\) and \(|y|>Y\), apply Cauchy--Schwarz on each,
and add before squaring.  This proves
\begin{equation}\tag{GMF22}
 |\langle P,F\rangle_a|^2
 \le e_a\norm P_a^2\norm F_a^2,\qquad
 a=m_k,\widetilde m_k,
\end{equation}
where \(e_{m_k}=e_m,e_{\widetilde m_k}=e_f\).
This bound holds for every high polynomial, uniformly in all
its coefficients.

Let \(\Pi_{\mathcal L_N}^a\) be the orthogonal projection in
\(\C[S]_{\le N}\) onto the \emph{whole} relation space
\(\mathcal L_N=\chi\C[S]_{\le N-q}\).
The original attained representative of a low class is
\(P-\Pi_{\mathcal L_N}^aP\).
Every vector of \(\mathcal L_N\) is a high polynomial in
(GMF22).  Taking the supremum over its unit sphere proves
\(\norm{\Pi_{\mathcal L_N}^aP}_a^2\le e_a\norm P_a^2\).
Pythagoras then gives the full matrix statement in the
fixed low-monomial frame:
\begin{equation}\tag{GMF23}
 (1-e_a)H_d[a]\preceq G_N[a]\big|_{\PP_d}
                         \preceq H_d[a].
\end{equation}
For the empty relation space at \(N=q-1\), this is the
exact identity \(G_{q-1}|_{\PP_d}=H_d\) with an admissible
looser bound.  Thus the original minimum has been evaluated
up to a proved error for every relation direction, rather
than by its value on one trial representative.

Let \(v_N=R_N[a]v_\lambda\) be the terminal minimizer.
By (GMF17) it is a high polynomial.  Its orthogonality to
\(\mathcal L_N\) gives
\[
 \langle v_N,P-\Pi_{\mathcal L_N}^aP\rangle_a
                       =\langle v_N,P\rangle_a.
\]
Equation (GMF22) and the lower inequality of (GMF23)
therefore bound the squared angle of the terminal value
with the \emph{entire} low quotient subspace by \(z_a\).
Explicitly, in the flag frame the Gram is
\[
 \begin{pmatrix} E_{N,\lambda}[a]&B_{N,a}\\
                 B_{N,a}^*&C_{N,a}\end{pmatrix},\qquad
 C_{N,a}=G_N[a]|_{\PP_d}.
\]
The coefficient row \(B_{N,a}\) contains all terminal--low
pairings.  Completing the square gives
\begin{equation}\tag{GMF24}
\begin{gathered}
 j_{N,a}=
 \frac{B_{N,a}C_{N,a}^{-1}B_{N,a}^*}{E_{N,\lambda}[a]},
 \quad 0\le j_{N,a}\le z_a<1,\\
 \det\Gram_{G_N[a]}F_{\lambda,d}
 =E_{N,\lambda}[a]\det C_{N,a}\,(1-j_{N,a}).
\end{gathered}
\end{equation}
For the upper bound on \(j\), its numerator divided by \(E\)
is exactly the supremum of the squared normalized terminal
pairing over the entire low subspace, by Cauchy--Schwarz
in the \(C_{N,a}\)-metric, with equality in the direction
\(C_{N,a}^{-1}B_{N,a}^*\).
Thus no entry of the original cross block has been discarded.

\section{Growing marked determinant and its four-cutoff cancellation}

For \(0\le x<1\) set \(h_0(x)=-\log(1-x)\), and define
\begin{equation}\tag{GMF25}
 \Omega_{k,d}=
 h_0(\beta_k)
 +(d+1)\{h_0(e_m)+h_0(e_f)\}
 +h_0(z_m)+h_0(z_f).
\end{equation}
The preceding finite guards make every term well defined.
Equations (GMF18),(GMF23),(GMF24) imply, simultaneously at
every original cutoff,
\begin{equation}\tag{GMF26}
\begin{split}
 \left|
 \log\frac{\det\Gram_{G_N[\widetilde m_k]}F_{\lambda,d}}
              {\det\Gram_{G_N[m_k]}F_{\lambda,d}}
 -\log\frac{\det H_d[\widetilde m_k]}{\det H_d[m_k]}
 \right|\le\Omega_{k,d}.
\end{split}
\end{equation}
Indeed the terminal norm-ratio logarithm is at most
\(h_0(\beta_k)\) in absolute value.
For each source the logarithmic determinant difference
between \(C_{N,a}\) and \(H_d[a]\) belongs to
\([-(d+1)h_0(e_a),0]\).
The two cross-block logarithms belong respectively to
\([-h_0(z_a),0]\).
Adding their absolute error bounds proves (GMF26) with
the stated full rank \(d+1\).

The source determinant term in (GMF26) is the same at
each cutoff \(N\).  Its four signed coefficients sum to zero.
Consequently
\begin{equation}\tag{GMF27}
 \boxed{
 \left|\mathcal R_{\rm ar}
 \log\frac{\det\Gram_{G_N[\widetilde m_k]}F_{\lambda,d}}
              {\det\Gram_{G_N[m_k]}F_{\lambda,d}}
 \right|\le4\Omega_{k,d}.}
\end{equation}

Take \(0\le d\le\min(d_k,q-2)\).
The finite constants show
\[
 \log\ell_m,\ \log\ell_f
 \le-\omega_*\sqrt{kq}+O_h(k+\log q).
\]
For \(\ell_m\), this uses the explicit formula for \(a_k\)
in (GMF20), the bound \(\log A_d=O_h(k+\log k)\), and
\(\log D_d=O_h(\log q)\); no unrecorded mass is cancelled.
For \(\ell_f\) use \(\eta_{k,d}<1/2\).
The bound (GMF19) for \(\beta_k\) is smaller than these
bounds on the eventual fixed-packet range.
Since \((\sqrt x+\sqrt y)^2\le2x+2y\), each \(e_a\) obeys
the same leading exponential bound.  The functions
\(h_0(x)\) and \(x/(1-x)\) are at most fixed multiples
of \(x\) on the finite guards in (GMF21).
Multiplication by \(d+1=O_h(k)\) adds only \(O_h(\log k)\)
to the logarithmic error.  Hence
\begin{equation}\tag{GMF28}
 \Omega_{k,d}\le
 \exp[-(\pi/4)\sqrt{kq}+O_h(k+\log q)]
 =O_h(e^{-(\pi/8)\sqrt{kq}}).
\end{equation}
All the strict finite guards hold eventually, uniformly in
this growing range, as proved by these explicit bounds.
At fixed packet \(d_k\le q-2\) eventually, since \(d_k=O_h(k)\)
and \(q\ge(k+1)^2\).

The per-cutoff logarithmic volumes themselves have a very
different evaluated lower bound.  From (GMF8),(GMF10),
\begin{equation}\tag{GMF29}
 \log\frac{\det H_d[\widetilde m_k]}{\det H_d[m_k]}
 \ge(d+1)\bigl[
 ka_h-\log(2A_d)+\log(1-\eta_{k,d})\bigr].
\end{equation}
This is the determinant of a full Loewner comparison on
the unchanged degree-\(d\) coefficient space.
At \(d=d_k=a_hk/(4\ell_*)+O(1)\),
\(\log A_d=2d\ell_*+2\log(d+1)+O(1)\).
Inserting this into (GMF29), and then (GMF26), yields
\begin{equation}\tag{GMF30}
 \boxed{
 \log\frac{\det\Gram_{G_N[\widetilde m_k]}F_{\lambda,d_k}}
              {\det\Gram_{G_N[m_k]}F_{\lambda,d_k}}
 \ge\frac{a_h^2}{8\ell_*}k^2-O_h(k\log k)}
\end{equation}
uniformly over the full original cutoff window.
The coefficient uses the fixed actual mass ratio \(a_h\).
It has not been obtained from a rank allocation of the
full signed determinant.

\section{Relative eigenvalues: both the large part and the majority}

Let
\[
 \mathcal A_N=
 G_N[m_k]^{-1/2}G_N[\widetilde m_k]G_N[m_k]^{-1/2}
\]
be the positive Hermitian representative of the relative metric.
The eigenvalues are independent of the chosen fixed value frame,
being the generalized eigenvalues of the two original forms.
On the low class subspace of dimension \(d+1\),
(GMF23),(GMF8),(GMF10) give the exact Rayleigh bound
\begin{equation}\tag{GMF31}
 \frac{\norm{[P]_\chi}_{G_N[\widetilde m_k]}^2}
      {\norm{[P]_\chi}_{G_N[m_k]}^2}
 \ge
 \frac{(1-e_f)(1-\eta_{k,d})e^{ka_h}}{2A_d}
 \quad(0\ne P\in\PP_d).
\end{equation}
The map from this class subspace into the Euclidean space of
\(\mathcal A_N\) is multiplication by \(G_N[m_k]^{1/2}\);
it is an isomorphism onto a \((d+1)\)-dimensional subspace.
If fewer than \(d+1\) eigenvalues were at least the right
side of (GMF31), this subspace would intersect the span
of the remaining smaller-eigenvalue eigenvectors nontrivially,
contradicting (GMF31).
At \(d=d_k\) this proves
\begin{equation}\tag{GMF32}
 \#\{\lambda_j(\mathcal A_N):
 \log\lambda_j\ge a_hk/2-O_h(\log k)\}\ge d_k+1.
\end{equation}
The argument proves the needed finite-dimensional min--max
consequence directly; it does not choose trial eigenvectors
or a preferred quotient section.

For the complementary estimate, take any literal monic
divisor \(D\mid\chi\) of degree \(j\), including its selected
repetitions.  In real coordinates make the monic factor
\(i^{-j}D(c+iy)\); every root still has modulus at most \(R_k\).
For any polynomial of degree at most \(2q\) divisible by
this factor, repeat the full-coefficient Legendre proof of
(GMF16), this time allowing the remaining degree through
\(2q\).  It gives relative error at most
\begin{equation}\tag{GMF33}
 B_{k,j}=A_k^{\rm high}\rho_Y^{2j},\qquad
 A_k^{\rm high}=
 \frac{\mathcal M_k}{\underline m_{k,q}}
 \frac{(2q+1)^2}{q}\,9^{4q}.
\end{equation}
Specifically the Legendre sum is bounded by
\((2q+1)^2 9^{4q}/q\); the remaining steps use exactly
the numerator mass and denominator density in (GMF14).
Put
\begin{equation}\tag{GMF34}
 j_k=\left\lceil
 \frac{\log A_k^{\rm high}+q}{2\log(1/\rho_Y)}
 \right\rceil ,
 \qquad 1\le j_k<q/2.
\end{equation}
The displayed degree inequalities are finite guards.
They hold eventually: \(\log A_k^{\rm high}=O_h(q+k+\log q)\),
while
\(2\log(1/\rho_Y)=\log(q/k)+o_h(1)\to\infty\).
This also proves
\(j_k=O_h(q/\log(q/k))\).
The ceiling in (GMF34) gives \(B_{k,j_k}\le e^{-q}\) exactly.

Choose the divisor \(D\) as the product of the first \(j_k\)
factors in the original ordered factor list for \(\chi\),
with each repetition retained.  The multiplication map
\begin{equation}\tag{GMF35}
 \C[S]/(\chi/D)\longrightarrow E,\qquad
 [P]_{\chi/D}\longmapsto[DP]_\chi
\end{equation}
is well defined and injective: \(\chi\mid DP\) is
equivalent in the polynomial domain to \(\chi/D\mid P\).
Its image \(\mathcal V_D\) has dimension \(q-j_k\).
Every polynomial representative of a class in \(\mathcal V_D\)
is divisible by \(D\), since its difference from a displayed
representative is divisible by \(\chi\).
The bounds (GMF33) therefore apply to every representative
in its original minimum fibre.  Taking both minima proves
\begin{equation}\tag{GMF36}
 (1-e^{-q})\,G_N[m_k]|_{\mathcal V_D}
 \preceq G_N[\widetilde m_k]|_{\mathcal V_D}
 \preceq(1+e^{-q})\,G_N[m_k]|_{\mathcal V_D}.
\end{equation}
This is the same exact divisibility-to-minimum map as
CPT25--CPT27, here with \(g=\chi/D\) and the new enlarged
localization interval \(Y\).

In the relative Euclidean coordinates used above, the
subspace in (GMF36) has codimension \(j_k\).
If the spectral subspace below \(1-e^{-q}\) had dimension
greater than \(j_k\), it would intersect this subspace,
contradicting its lower Rayleigh bound.
The same dimension argument above \(1+e^{-q}\) proves
\begin{equation}\tag{GMF37}
 \boxed{
 \#\{\lambda_j(\mathcal A_N)\notin
          [1-e^{-q},1+e^{-q}]\}
 \le2j_k=O_h(q/\log(q/k)).}
\end{equation}
Consequently, for every bounded continuous function \(f\),
\[
 \left|\frac1q\sum_{j=1}^q f(\lambda_j)-f(1)\right|
 \le\sup_{|x-1|\le e^{-q}}|f(x)-f(1)|
                    +4\norm f_\infty\,j_k/q\longrightarrow0 .
\]
Both estimates are uniform over all original cutoffs.
The linearly many large eigenvalues in (GMF32) coexist
with this concentration because their proportion tends
to zero.  Their logarithmic volume contribution is
explicitly retained in (GMF30).

\section{Marked derivatives, period data and the attained complement}

The original marked coefficient module is
\[
 \mathscr H=\C[u,u^{-1},t][S]\big/
       (u\partial_S+\chi-t)\C[u,u^{-1},t][S].
\]
Its vectors have unique representatives of degree below
\(q\); this quotient is a module, with the original marked
connection, and is not used as a quotient ring.
Indeed subtracting the differential relation applied to
a leading monomial reduces any polynomial's degree.
A nonzero \(P\) has
\(\deg((u\partial_S+\chi-t)P)=q+\deg P\), proving
uniqueness of a representative of degree below \(q\).
The cochain operator \(\partial_t-S/u\) descends because
\[
 [\partial_t-S/u,u\partial_S+\chi-t]=-1+1=0.
\]
On the chosen representatives it is exactly the original
\(\nabla_t=\partial_t-A(t)/u\).
At \(t=0\), identify its coefficient representatives with
the displayed vector space \(E\), and fix \(u\ne0\).
The exact marked recurrence from the marked-map
calculation \cite{a02gmf:markedProvider} follows directly here.
Let \(R_j\) be the unique representative of
\(S^j\chi/(S-\lambda)\).  The cochain formula gives
\(\nabla_t^jv_\lambda=b_\lambda(-u)^{-j}R_j\).
Applying the original relation to \(S^{j-1}\) gives
\[
 R_j-\lambda R_{j-1}
   =tS^{j-1}-u(j-1)S^{j-2}\qquad(1\le j\le q),
\]
whose right side already has degree below \(q\).
The final term is zero at \(j=1\).
Thus at \(t=0\) one has
\[
 v_0=v_\lambda,\quad v_j=(\nabla_t^j v_\lambda)|_{t=0},
 \qquad v_1=-\lambda v_0/u,
\]
\begin{equation}\tag{GMF38}
 v_j+\frac{\lambda}{u}v_{j-1}
   =b_\lambda(j-1)(-u)^{1-j}S^{j-2},
                    \qquad 2\le j\le q.
\end{equation}
It retains the complete-primary \(b_\lambda\).
For \(j=2\), insert the displayed \(v_1\), giving
\(v_2=\lambda^2v_0/u^2-b_\lambda/u\).
For each successive \(j\), (GMF38) expresses \(v_j\)
as a nonzero scalar multiple of \(S^{j-2}\) plus preceding
columns.  Thus in increasing order
\[
 (v_0,v_2,\ldots,v_{d+2})
       =F_{\lambda,d}\,\mathcal T_{\lambda,d}(u),
\]
where the triangular matrix has diagonal
\[
 1,\ b_\lambda(-u)^{-1},\
 2b_\lambda(-u)^{-2},\ldots,
 (d+1)b_\lambda(-u)^{-(d+1)}.
\]
Its complete modulus-squared determinant is
\begin{equation}\tag{GMF39}
 |\det\mathcal T_{\lambda,d}(u)|^2
 =|b_\lambda|^{2(d+1)}((d+1)!)^2
                        |u|^{-(d+1)(d+2)}.
\end{equation}
This proves the partial marked frame map and the exact
Gram congruence.  The scalar cancels between the two source
metrics at each cutoff because the frame is fixed.
Its magnitude is never set to one.  Equations (GMF26)--(GMF30)
therefore hold for the literal marked derivative columns.
The original period isomorphism has its own fixed matrix
\(\Pi\) at the retained nonzero period and branch.  Its exact
transported metric is
\[
 \widehat G_N[a]=(\Pi^{-1})^*G_N[a]\Pi^{-1},\qquad
 (\Pi F_{\lambda,d})^*\widehat G_N[a](\Pi F_{\lambda,d})
                  =F_{\lambda,d}^*G_N[a]F_{\lambda,d}.
\]
Thus the same estimate holds on the literal period images
when equipped with the transported arithmetic metric.
The Euclidean period Gram is separately
\(F_{\lambda,d}^*\Pi^*\Pi F_{\lambda,d}\);
its exact map is the displayed \(\Pi\), and it is not
substituted for either attained arithmetic form here.

Let \(r=d+2\), and complete \(F_{\lambda,d}\) by any fixed
coefficient columns \(C_{\lambda,d}\) to an invertible frame
\(\mathcal J=(F_{\lambda,d},C_{\lambda,d})\) of \(E\).
Use the quotient frame consisting of the images of
\(C_{\lambda,d}\) in \(E/\Fcal_{\lambda,d}\).
Then the actual quotient map \(\pi_\Fcal\) satisfies
\(\pi_\Fcal\mathcal J=(0,I_{q-r})\).
In this completed frame write the original Gram
\[
 \mathcal J^*G_N[a]\mathcal J=
 \begin{pmatrix}A_{N,a}&B_{N,a}\\B_{N,a}^*&C_{N,a}\end{pmatrix}.
\]
Every value \(z\) in the quotient is represented by
\(\mathcal J(x,z)\).  Completing its squared norm gives
the unique minimum at \(x=-A_{N,a}^{-1}B_{N,a}z\).
Consequently its exact attained metric is
\begin{equation}\tag{GMF40}
\begin{gathered}
 Q_{\Fcal,N}[a]
 =(\pi_\Fcal G_N[a]^{-1}\pi_\Fcal^*)^{-1}
 =C_{N,a}-B_{N,a}^*A_{N,a}^{-1}B_{N,a},\\
 |\det\mathcal J|^2\det G_N[a]
   =\det A_{N,a}\det Q_{\Fcal,N}[a].
\end{gathered}
\end{equation}
All original cross entries remain in this Schur complement.
The fixed Jacobian \(|\det\mathcal J|^2\) cancels between
the two sources.

Define the exact central signed endpoint value
\(\delta_k^{\rm cen}=\Bcal_k[m_k]-\Bcal_k[\widetilde m_k]\).
Taking the actual four signs in (GMF40) and using (GMF27)
proves the complete finite receiving theorem
\begin{equation}\tag{GMF41}
\begin{split}
 \mathcal R_{\rm ar}
 \log\frac{\det Q_{\Fcal,N}[\widetilde m_k]}
              {\det Q_{\Fcal,N}[m_k]}
 &=-\delta_k^{\rm cen}+\mathcal E_{\Fcal,k},\\
 \mathcal E_{\Fcal,k}
 &:=-\mathcal R_{\rm ar}
 \log\frac{\det\Gram_{G_N[\widetilde m_k]}F_{\lambda,d}}
              {\det\Gram_{G_N[m_k]}F_{\lambda,d}},
 \qquad |\mathcal E_{\Fcal,k}|\le4\Omega_{k,d}.
\end{split}
\end{equation}
This theorem evaluates the error of transport to the
\emph{actual attained complement}.  Its endpoint value is
the unchanged central endpoint, with the displayed minus
sign.  The accepted finite-amplitude endpoint evaluation is
inserted below with its full source error; no first
derivative is substituted for this endpoint difference.

For clarity the exact map to an independently specified
original observation \(O:E\to Y\) has its own criterion.
It factors through \(\pi_\Fcal\) precisely when
\(O|_\Fcal=0\): the necessity follows by applying the
factorization to the kernel of \(\pi_\Fcal\), and the
sufficiency follows by defining \(\bar O([v])=Ov\).
Without that vanishing the original observation remains
the pair of coefficient maps
\((OF_{\lambda,d},OC_{\lambda,d})\) in the completed frame,
with their full blocks.  In particular (GMF40) has not
identified its quotient with an original observation
quotient or with an induced arithmetic action.

\section{The preceding constant-class receiver is sharpened}

At \(N=q-1\) the relation space is zero.  Define the original
one-packet variance and two retained tails by
\[
 V_0=\frac{\int y^2w_h(y)dy}{M_0},\quad
 D_0=\log(M_\alpha/M_0),\quad
 s_T=\int_{|y|>T}\sigma(y)dy,\quad
 p_{k,T}=M_0^{-k}\int_{|y|>T}m_k(y)dy.
\]
The auxiliary unit-mass law \(w_h/M_0\) has mean zero by
evenness and variance \(V_0>0\).
Its \(k\)-fold sum has variance \(kV_0\), by expansion of
the centred square and vanishing of every cross mean.
Integrating \(y^2/T^2\ge1\) on \(|y|>T=kb\) gives
\[
 0\le p_{k,T}\le V_0/(kb^2),\qquad
 0\le s_T\le
 \frac{2C_\Gamma}{\alpha\sqrt{1+T}}e^{-\alpha T}.
\]
Keeping both tails and both original masses gives the
exact constant-class ratio
\begin{equation}\tag{GMF42}
 \frac{\norm{1}_{G_{q-1}[\widetilde m_k]}^2}
      {\norm{1}_{G_{q-1}[m_k]}^2}
 =e^{kD_0}(M_\sigma-s_T)+p_{k,T}.
\end{equation}
This verifies B7--B8 of the marked arithmetic integration
source \cite{a02gmf:markedProvider}.
The terminal Rayleigh quotient in (GMF18) is at most
\(1+\beta_k\); the ratio of the largest to the smallest
generalized eigenvalue is therefore at least
\begin{equation}\tag{GMF43}
 \cond(\mathcal A_{q-1})
 \ge\frac{e^{kD_0}(M_\sigma-s_T)+p_{k,T}}{1+\beta_k}
 =\sqrt{2\pi}\,e^{kD_0}(1-o_h(1)).
\end{equation}
If the sharper CPT pair radius is used in the terminal
comparison, the same proof uses its exact upper factor
instead; both are original-data bounds.
At every other original cutoff, (GMF23) with \(d=0\)
gives the uniform extension
\[
 \cond(\mathcal A_N)\ge
 \frac{(1-e_f|_{d=0})
       [e^{kD_0}(M_\sigma-s_T)+p_{k,T}]}
      {1+\beta_k}.
\]
Here the original arithmetic constant-class denominator
is at most \(M_0^k\), and the filled numerator is at least
\((1-e_f|_{d=0})\int\widetilde m_k\); the terminal
Rayleigh upper bound remains (GMF18).

For \(d=0\), \(H_0[m_k]=M_0^k\).
Equation (GMF10) therefore gives
\[
 \left|\log\frac{H_0[\widetilde m_k]}{H_0[m_k]}
            -kD_0-\log M_\sigma\right|
 \le h_0(\eta_{k,0})+\log(1+\zeta_{k,0}).
\]
Return it through the complete cross-block formula
(GMF26) and then the marked congruence (GMF39).
This proves at \emph{every} original cutoff the finite bound
\begin{equation}\tag{GMF44}
\begin{split}
 \left|
 \log\frac{\det\Gram_{G_N[\widetilde m_k]}(v_\lambda,\nabla_t^2v_\lambda)}
              {\det\Gram_{G_N[m_k]}(v_\lambda,\nabla_t^2v_\lambda)}
       -kD_0-\tfrac12\log(2\pi)
 \right|
 \le\Omega_{k,0}+h_0(\eta_{k,0})+\log(1+\zeta_{k,0}).
\end{split}
\end{equation}
Its right side is exponentially small at fixed packet.
This strengthens the prior first-cutoff B11 remainder
\(O_h(k^2/q^2)\), using full-fibre localization instead
of only its parity estimate.  The former statement and
its original proof remain preserved.

\section{The evaluated central value in the original attained complement}

The accepted signed arithmetic source proof
\cite[C32 and SAI1--SAI11]{a02gmf:signedCentral} evaluates the
\emph{finite-amplitude endpoint} appearing in (GMF41):
\[
 I_h=\int_0^\alpha[(\log M)'(\theta)]^2\,d\theta>0,\qquad
 \delta_k^{\rm cen}
  =\frac{\log2}{\pi}I_h k^2+E_k^{\rm cen}.
\]
On its original fixed-packet range the retained error is
\[
 |E_k^{\rm cen}|
 \le \mathfrak C_{h,k}^{\rm cen},\qquad
 \mathfrak C_{h,k}^{\rm cen}
 =A_h\left[
 \frac{k^2}{\log^3(q+2)}+k\log^2(q+2)+1+\frac{k^4}{q^2}
 \right].
\]
Here \(A_h\) is the fixed-packet bound in that accepted
source calculation, with its original exact Robin and
Neumann guards.  Specifically its finite monic-band part
is \(\sum_{r=0}^q c_rE_H(q+r)+k^2I_h/(16\pi q^2)\)
from SAI11; its source-localization and smoothing errors
are added as in C3,C8 before the displayed fixed-packet
bound is taken.  These terms are inherited with their
whole-line integration and both original endpoint weights.
This use is a receiving substitution of that completed
calculation, not a new proof by differentiating at zero.

At \(d=d_k\), insert the value into the exact (GMF41).
The complete result is
\begin{equation}\tag{GMF45}
\begin{split}
 \boxed{
 \mathcal R_{\rm ar}
 \log\frac{\det Q_{\Fcal,N}[\widetilde m_k]}
              {\det Q_{\Fcal,N}[m_k]}
 =-\frac{\log2}{\pi}I_h k^2+\epsilon_{\Fcal,k}},\\
 \epsilon_{\Fcal,k}=-E_k^{\rm cen}+\mathcal E_{\Fcal,k},
 \qquad
 |\epsilon_{\Fcal,k}|
 \le\mathfrak C_{h,k}^{\rm cen}+4\Omega_{k,d_k}.
\end{split}
\end{equation}
All localization guards (GMF14),(GMF21), the degree
guard \(d_k\le q-2\), and the source theorem's original
finite-amplitude guards remain in force for this receiver.
The separate majority-eigenvalue theorem retains its
own additional divisor-degree guard (GMF34).
On the simple stratum \(q=(k+1)^2\), the coefficient at
order \(q\) in this attained complementary return is
\(-(\log2/\pi)I_h\), since \(k^2/q\to1\) and the displayed
error is \(o_h(q)\).
For fixed \(m_0>1\) the exact displayed scale remains
\(k^2\), with \(q\asymp k^3\); that scale has not been
relabeled as an order-\(q\) coefficient.

% BEGIN COMPLETE FCD SOURCE
% Independent reconstruction of the fixed-d calculation, incoming (36)--(39).
% Include in a document loading amsmath, amssymb, and hyperref.
% This fragment has no external input, graphics, or bibliography dependency.
\section{Exact fixed-degree cumulant determinants with the original source masses}
\label{a02gmf:sec:FCD}

The incoming calculation is the second part of the supplied text headed
\emph{A growing marked flag has large individual metric changes---but an
exponentially small signed return}, equations (36)--(39).
This section proves its finite moment calculation, all four compressed
Hermite traces, the original Gamma constants, and the finite cubic
remainder. The attained-source comparison is retained as an exact receiving
identity at the end. Its analytic error is evaluated by the accompanying
growing-flag localization argument.

\subsection{Original objects and exact coordinate maps}

Retain the actual fixed packet, including its multiplicity:
\[
 h(s)=\prod_{\zeta\in\{\rho,\bar\rho,1-\rho,1-\bar\rho\}}
       (s-\zeta)^{m_0},\qquad
 \rho=\tfrac12+\delta+i\gamma,\quad
 0<\delta<\tfrac12,\quad\gamma>2,\quad h\mid2\xi.
\]
Put \(v_h=2\xi/h\), with its unchanged unit, and
\[
\begin{gathered}
 w_h(y)=\frac{|v_h(1/2+iy)|^2}{2\pi},\qquad
 m_k=w_h^{*k},\qquad c=k/2,\\
 q=[1+k(m_0-1)](k+1)^2,\qquad k\equiv1\pmod4.
\end{gathered}
\]
The original packet gives an even nonnegative source with all finite moments.
Evenness follows from the real coefficients of \(h\) and the conjugation
identity of \(\xi\). Since \(v_h\) is a nonzero analytic function, its
zeros on this line are isolated, so \(w_h\) is positive almost everywhere.
In particular the following original constants are positive:
\begin{equation}
 \begin{split}
 M(\theta)&=\int_{\mathbb R}e^{\theta y}w_h(y)\,dy,\qquad
 M_0=M(0),\qquad M_\alpha=M(\alpha),\qquad \alpha=\pi/2,\\
 V&=\frac1{M_0}\int_{\mathbb R}y^2w_h(y)\,dy>0,\qquad
 D_0=\log(M_\alpha/M_0).
 \end{split}
 \tag{FCD1}\label{a02gmf:eq:FCD1}
\end{equation}
The endpoint moment \(M_\alpha\) is the finite endpoint moment of this
packet; it is not replaced by a probability mass. Let
\[
 \mu_j=M_0^{-1}\int_{\mathbb R}y^jw_h(y)\,dy,\qquad
 c_4=\frac{\mu_4-3V^2}{V^2},\qquad
 c_6=\frac{\mu_6-15\mu_4V+30V^3}{V^3}.
\]
More generally \(c_{2r}=\kappa_{2r}(w_h/M_0)/V^r\), with \(c_2=1\);
the cumulants are the coefficients of the formal logarithm of the moment
series. Only finitely many are used for any fixed degree.

For a positive density \(\nu\), let \(H_d[\nu]\) denote the
\((d+1)\)-by-\((d+1)\) Gram matrix of
\((1,S,\ldots,S^d)\) on the literal line \(S=c+iy\):
\[
 (H_d[\nu])_{ij}
   =\int_{\mathbb R}\overline{(c+iy)^i}(c+iy)^j\nu(y)\,dy,
       \qquad 0\le i,j\le d.
\]
Define the auxiliary probability measure
\[
 d\mu^{(k)}(x)
  =M_0^{-k}\sqrt{kV}\,m_k(\sqrt{kV}\,x)\,dx
\]
and its moment matrix \(G_d^{(k)}\) in \(1,x,\ldots,x^d\).
The exact pullback map on polynomials is
\[
 \mathcal U_k:\mathbb C[S]_{\le d}\longrightarrow\mathbb C[x]_{\le d},
 \qquad P\longmapsto P(c+i\sqrt{kV}\,x).
\]
Its matrix \(B\), with columns indexed by powers of \(S\), has entries
\[
 B_{\ell j}=
 \begin{cases}
 \binom j\ell c^{j-\ell}(i\sqrt{kV})^\ell,&\ell\le j,\\
 0,&\ell>j .
 \end{cases}
\]
The inverse substitutes \(x=(S-c)/(i\sqrt{kV})\).
Consequently, with \(D=d(d+1)/2\),
\begin{equation}
 H_d[m_k]=M_0^kB^*G_d^{(k)}B,\qquad
 \det B=(i\sqrt{kV})^D,\qquad
 \det H_d[m_k]=M_0^{k(d+1)}(kV)^D\det G_d^{(k)} .
 \tag{FCD2}\label{a02gmf:eq:FCD2}
\end{equation}
This equality retains the physical centre, all factors \(i\), the entire
original mass, and the coordinate Jacobian. The auxiliary measure supplies
finite coefficient identities; \(\mathcal U_k\) and \eqref{a02gmf:eq:FCD2}
return them to the unchanged source metric.

\subsection{An exact finite polynomial in \(1/k\)}

For each moment order at most \(2d\), independence in the convolution and
the formal logarithm give
\begin{equation}
 \sum_{n=0}^{2d}\frac{\int x^n\,d\mu^{(k)}(x)}{n!}z^n
 \equiv
 \exp\left(\frac{z^2}{2}
       +\sum_{r=2}^{d}\frac{c_{2r}}{(2r)!}\,
                    k^{1-r}z^{2r}\right)\pmod {z^{2d+1}}.
 \tag{FCD3}\label{a02gmf:eq:FCD3}
\end{equation}
Indeed the cumulant of order \(2r\) of the sum is \(k\) times the
one-packet cumulant, and the exact map \(\mathcal U_k\) contributes
\((kV)^{-r}\). Odd cumulants vanish by evenness. These assertions can
also be obtained by taking the formal logarithm of
\((M(z/\sqrt{kV})/M_0)^k\); no infinite-series approximation is needed.

Write \(\varepsilon=1/k\), and define polynomials \(\mathfrak m_n(\varepsilon)\)
by the finite coefficients in \eqref{a02gmf:eq:FCD3}. Explicitly,
\[
 \mathfrak m_{2a}(\varepsilon)
 =(2a)!\!
 \sum_{\substack{n_1,\ldots,n_a\ge0\\
                         \sum_{r=1}^a r n_r=a}}
 \prod_{r=1}^a
  \frac{(c_{2r}\varepsilon^{r-1})^{n_r}}
       {n_r!((2r)!)^{n_r}},
 \qquad \mathfrak m_{2a+1}=0,\qquad \mathfrak m_0=1.
\]
Here \(c_2=1\) and \(\varepsilon^0=1\), including the \(r=1\) factor.
Define
\begin{equation}
 P_d(\varepsilon)=
 \frac{\det[\mathfrak m_{i+j}(\varepsilon)]_{i,j=0}^{d}}
      {\prod_{j=0}^d j!}.
 \tag{FCD4}\label{a02gmf:eq:FCD4}
\end{equation}
This is a finite real polynomial. At \(\varepsilon=0\) these moments are
Gaussian moments, and the calculation below proves \(P_d(0)=1\).
For the actual positive integer \(k\), its moment matrix is positive
definite: the integral of the square of a nonzero polynomial against
\(\mu^{(k)}\) is strictly positive. In particular \(P_d(1/k)>0\).
Combining \eqref{a02gmf:eq:FCD2} and \eqref{a02gmf:eq:FCD4} proves the exact identity
\begin{equation}
 \boxed{\displaystyle
 \det H_d[m_k]=M_0^{k(d+1)}(kV)^{d(d+1)/2}
                    \left(\prod_{j=0}^d j!\right)P_d(1/k).}
 \tag{FCD5}\label{a02gmf:eq:FCD5}
\end{equation}

\subsection{Hermite identities and all compressed traces}

We use the probabilists' Hermite generating function recorded by
T.~H.~Koornwinder, R.~Wong, R.~Koekoek, R.~F.~Swarttouw and
W.~P.~Reinhardt in the NIST Digital Library of Mathematical Functions,
\href{https://dlmf.nist.gov/18.12.E16}{equation 18.12.16}.
The product formula below is its direct coefficient consequence and agrees
with the probabilists' rescaling of their
\href{https://dlmf.nist.gov/18.18.E23}{equation 18.18.23}.
We give the derivation so that the trace calculation has no unstated
orthogonal-polynomial convention.
Set
\[
 d\gamma(x)=(2\pi)^{-1/2}e^{-x^2/2}\,dx,\qquad
 e_n(x)=\operatorname{He}_n(x)/\sqrt{n!},\qquad
 e^{xz-z^2/2}=\sum_{n\ge0}\operatorname{He}_n(x)z^n/n!.
\]
Completing the square in the Gaussian integral gives
\[
 \int e^{xz-z^2/2}e^{xt-t^2/2}\,d\gamma(x)=e^{zt}.
\]
Comparing the \(z^mt^n\) coefficients proves
\(\int \operatorname{He}_m\operatorname{He}_n\,d\gamma
   =n!\,\mathbf1_{\{m=n\}}\).
The defining generating function also shows that
\(\operatorname{He}_n\) is monic. Its triangular change from monomials has
determinant one; hence the Gaussian monomial Gram determinant equals
\(\prod_{n=0}^d n!\), proving \(P_d(0)=1\).
The identity
\[
 e^{xz-z^2/2}e^{xt-t^2/2}
    =e^{zt}e^{x(z+t)-(z+t)^2/2}
\]
proves
\begin{equation}
 \operatorname{He}_j\operatorname{He}_n
   =\sum_{r=0}^{\min(j,n)}
       r!\binom jr\binom nr\operatorname{He}_{j+n-2r}.
 \tag{FCD6}\label{a02gmf:eq:FCD6}
\end{equation}

Let \(T_j\) be multiplication by \(\operatorname{He}_j\), compressed
orthogonally in \(L^2(\gamma)\) to
\(\operatorname{span}\{e_0,\ldots,e_d\}\). Formula \eqref{a02gmf:eq:FCD6}
proves its complete entries:
\begin{equation}
 (T_j)_{mn}=
 \begin{cases}
 \displaystyle
 \frac{j!\sqrt{m!n!}}{a!\,b!\,r!},
 &
\begin{gathered}
 a=(j+n-m)/2,\quad b=(j+m-n)/2,\\ r=(m+n-j)/2
       \ \text{are nonnegative integers},
\end{gathered}\\[4pt]
 0,&\text{otherwise}.
 \end{cases}
 \tag{FCD7}\label{a02gmf:eq:FCD7}
\end{equation}
This expression also proves symmetry. Write
\(x^{\underline r}=x(x-1)\cdots(x-r+1)\), with
\(x^{\underline0}=1\). Its value is zero for nonnegative integral
\(x<r\). On the diagonal,
\[
 (T_{2r})_{nn}=\binom{2r}{r}n^{\underline r}.
\]
The polynomial identity
\((n+1)^{\underline{r+1}}-n^{\underline{r+1}}
 =(r+1)n^{\underline r}\)
telescopes. Thus, for every integer \(d\ge0\),
\begin{equation}
\begin{aligned}
 \operatorname{Tr}T_{2r}
    &=\frac{\binom{2r}{r}}{r+1}(d+1)^{\underline{r+1}},\\
 \operatorname{Tr}T_4&=2(d+1)d(d-1),\\
 \operatorname{Tr}T_6&=5(d+1)d(d-1)(d-2),\\
 \operatorname{Tr}T_8&=14(d+1)d(d-1)(d-2)(d-3).
\end{aligned}
 \tag{FCD8}\label{a02gmf:eq:FCD8}
\end{equation}

All nonzero bands of \(T_4\), before removing indices outside \(0,\ldots,d\),
are
\[
 \begin{split}
 (T_4)_{nn}&=6n(n-1),\\
 (T_4)_{n+2,n}=(T_4)_{n,n+2}
       &=4n\sqrt{(n+1)(n+2)},\\
 (T_4)_{n+4,n}=(T_4)_{n,n+4}
       &=\sqrt{(n+1)(n+2)(n+3)(n+4)} .
 \end{split}
\]
Consequently the compression, including its two boundary bands, gives
\[
 \begin{split}
 \operatorname{Tr}(T_4^2)
 ={}&36\sum_{n=0}^d n^2(n-1)^2
     +32\sum_{n=0}^{d-2}n^2(n+1)(n+2)\\
    &+2\sum_{n=0}^{d-4}(n+1)(n+2)(n+3)(n+4).
 \end{split}
\]
An upper limit below zero denotes the empty sum. Put
\(F_r=(d+1)^{\underline r}\). The polynomial identities
\[
 n^2(n-1)^2=n^{\underline4}+4n^{\underline3}+2n^{\underline2},
 \qquad
 n^2(n+1)(n+2)=(n+2)^{\underline4}+(n+2)^{\underline3}
\]
and the same telescoping formula give, respectively, the three sums
\[
 F_5/5+F_4+2F_3/3,\qquad F_5/5+F_4/4,\qquad F_5/5.
\]
The missing lower-index falling factorials in the shifted sums vanish, so
these formulas include every empty-sum case. Hence
\begin{equation}
 \boxed{\displaystyle
 \operatorname{Tr}(T_4^2)=14F_5+44F_4+24F_3
             =2(d-1)d(d+1)(10-13d+7d^2).}
 \tag{FCD9}\label{a02gmf:eq:FCD9}
\end{equation}
In particular \(T_4=0\) for \(d=0,1\). No large-\(d\) qualification
has been used in \eqref{a02gmf:eq:FCD8} or \eqref{a02gmf:eq:FCD9}.

\subsection{The two exact logarithmic coefficients}

The elementary Gaussian identity
\[
 \int e^{zx}\operatorname{He}_j(x)\,d\gamma(x)
       =z^j e^{z^2/2}
\]
follows by integrating the Hermite generating function against \(e^{zx}\)
and comparing its coefficients. Apply it coefficientwise to
\eqref{a02gmf:eq:FCD3}. For every polynomial \(F\) of degree at most \(2d\),
the first three coefficients of its exact moment functional are
\[
 \begin{split}
 L_\varepsilon(F)
 ={}&\int F\,d\gamma
       +\frac{c_4\varepsilon}{24}
                       \int F\operatorname{He}_4\,d\gamma\\
   &+\varepsilon^2\left[
        \frac{c_6}{720}\int F\operatorname{He}_6\,d\gamma
       +\frac{c_4^2}{1152}\int F\operatorname{He}_8\,d\gamma
                         \right]+\varepsilon^3 R_F(\varepsilon),
 \end{split}
\]
where \(R_F\) is a finite polynomial, possibly zero. This is a
finite coefficient identity of moments. It makes no claim about a
density expansion or a distributional remainder.
In the basis \(e_0,\ldots,e_d\), the Gram matrix is exactly
\[
 I+\varepsilon A+\varepsilon^2 B+\varepsilon^3 R(\varepsilon),
 \qquad
 A=\frac{c_4}{24}T_4,\qquad
 B=\frac{c_6}{720}T_6+\frac{c_4^2}{1152}T_8,
\]
with a polynomial matrix \(R\).
The change from monomials to \(e_n\) has diagonal
\((n!)^{-1/2}\); its squared determinant is
\((\prod_{n=0}^d n!)^{-1}\). Thus this Gram determinant is exactly
\(P_d(\varepsilon)\), including all higher coefficients.

For completeness, the first two logarithmic coefficients are obtained
directly from the determinant permutation formula:
\[
 \det(I+\varepsilon A+\varepsilon^2B+O(\varepsilon^3))
 =1+\varepsilon\operatorname{Tr}A+
 \varepsilon^2\left(\operatorname{Tr}B+
    \tfrac12[(\operatorname{Tr}A)^2-\operatorname{Tr}(A^2)]\right)
 +O(\varepsilon^3).
\]
The diagonal choices give the trace terms; selecting two distinct
first-order entries gives
\(\sum_{i<j}(A_{ii}A_{jj}-A_{ij}A_{ji})\), which is the stated half
difference. Taking the scalar logarithm therefore gives
\[
 [\varepsilon]\log P_d=\operatorname{Tr}A,\qquad
 [\varepsilon^2]\log P_d=\operatorname{Tr}B-\tfrac12\operatorname{Tr}(A^2).
\]
Substitution of \eqref{a02gmf:eq:FCD8} and \eqref{a02gmf:eq:FCD9} yields
\begin{equation}
 \begin{split}
 A_d^{\rm cum}&=\frac{c_4d(d^2-1)}{12},\\
 B_d^{\rm cum}
   &=\frac{d(d^2-1)}{288}
                \left[(16-11d)c_4^2+(2d-4)c_6\right],\\
 \log P_d(1/k)&=\frac{A_d^{\rm cum}}k
                         +\frac{B_d^{\rm cum}}{k^2}
                         +r_{d,k}^{\rm cum}.
 \end{split}
 \tag{FCD10}\label{a02gmf:eq:FCD10}
\end{equation}
The coefficient of \(c_4^2\) contains the trace of the
square of the \emph{compressed} multiplication matrix. Replacing it
by the compression of the full square would discard boundary bands and
would change the answer.
The exact relation between these two operators, on the polynomial domain,
is
\[
 \Pi_d M_{\operatorname{He}_4}^2\Pi_d
       -(\Pi_dM_{\operatorname{He}_4}\Pi_d)^2
   =\Pi_dM_{\operatorname{He}_4}(I-\Pi_d)
                             M_{\operatorname{He}_4}\Pi_d
   =L_d^*L_d,\qquad
 L_d=(I-\Pi_d)M_{\operatorname{He}_4}\Pi_d .
\]
Here \(L_d\) maps the degree-at-most-\(d\) Hermite space into
\(\operatorname{span}\{e_{d+1},\ldots,e_{d+4}\}\).
Inserting \(I=\Pi_d+(I-\Pi_d)\) proves the equality;
self-adjointness of real-polynomial multiplication gives the last
factorization. Formula \eqref{a02gmf:eq:FCD7} evaluates every entry of
this finite boundary map as well.

\subsection{An explicit finite \(k^{-3}\) remainder}

Use the whole polynomial from \eqref{a02gmf:eq:FCD4}, including every
cumulant required through degree \(2d\):
\[
 P_d(\varepsilon)=1+\sum_{j\ge1}a_j\varepsilon^j,\qquad
 C_0=\sum_{j\ge1}|a_j|,\qquad C_3=\sum_{j\ge3}|a_j|.
\]
Both sums are finite. If \(k\ge\max(1,2C_0)\), put
\(u=P_d(1/k)-1\). Then
\[
 |u|\le C_0/k\le\tfrac12,\qquad
 \left|u-\frac{a_1}k-\frac{a_2}{k^2}\right|
       \le C_3/k^3,\qquad
 \left|u-\frac{a_1}k\right|\le C_0/k^2 .
\]
The scalar series for \(\log(1+u)\), absolutely convergent on this
closed half-disc, satisfies
\[
 |\log(1+u)-u+u^2/2|
   \le\sum_{j\ge3}|u|^j/j
   \le\frac{|u|^3}{3(1-|u|)}
   \le\frac{2C_0^3}{3k^3}.
\]
Also
\[
 \frac12\left|u^2-\frac{a_1^2}{k^2}\right|
 \le\frac12\left|u-\frac{a_1}k\right|
                  \left|u+\frac{a_1}k\right|
 \le C_0^2/k^3 .
\]
Since \(A_d^{\rm cum}=a_1\) and
\(B_d^{\rm cum}=a_2-a_1^2/2\), this proves
\begin{equation}
 \boxed{\displaystyle
 |r_{d,k}^{\rm cum}|
 \le \frac{C_3+C_0^2+\tfrac23 C_0^3}{k^3},
 \qquad k\ge\max(1,2C_0).}
 \tag{FCD11}\label{a02gmf:eq:FCD11}
\end{equation}
When \(C_0=0\), \(P_d=1\) and the error is identically zero.
Thus the finite guard includes \(d=0,1\).

\subsection{The original Gamma mass and every monic norm}

The comparison density is literally
\[
 \sigma(y)=\frac{|\Gamma(1/4+iy/2)|^2}{2\pi}.
\]
Here is a derivation of its transform and norms retaining the factor
\(2\pi\). Euler's beta integral
(\href{https://dlmf.nist.gov/5.12.E1}{NIST DLMF 5.12.1})
and \(u=e^x\) give, for \(a>0\),
\[
 \Gamma(a+it)\Gamma(a-it)
   =\Gamma(2a)\int_{\mathbb R}
           e^{itx}(2\cosh(x/2))^{-2a}\,dx.
\]
Fourier inversion therefore gives, initially for real \(z\),
\[
 \int_{\mathbb R}|\Gamma(a+it)|^2e^{-itz}\,dt
        =2\pi\Gamma(2a)(2\cosh(z/2))^{-2a}.
\]
Both sides are analytic for \(|\operatorname{Im}z|<\pi\).
For the integral this follows from the Gamma bound
\(|\Gamma(a+it)|=O_a((1+|t|)^{a-1/2}e^{-\pi|t|/2})\),
recorded in
\href{https://dlmf.nist.gov/5.11.E9}{NIST DLMF 5.11.9};
on compact substrips it also justifies differentiation under the integral.
The analytic identity theorem extends the formula to that strip.
Taking \(z=2i\theta\), \(a=1/4\), and \(y=2t\), gives
\begin{equation}
 \int_{\mathbb R}e^{\theta y}\sigma(y)\,dy
       =\sqrt{2\pi}(\cos\theta)^{-1/2},\qquad
 |\theta|<\pi/2,\qquad M_\sigma=\sqrt{2\pi}.
 \tag{FCD12}\label{a02gmf:eq:FCD12}
\end{equation}

Define the monic polynomials \(p_n\) by
\[
 K(z,y)=\sum_{n\ge0}p_n(y)z^n/n!
       =(1+z^2)^{-1/4}e^{y\arctan z}.
\]
This is the \((\lambda,x,\varphi)=(1/4,y/2,\pi/2)\)
Meixner--Pollaczek generating function in
\href{https://dlmf.nist.gov/18.23.E7}{DLMF 18.23.7}.
Its source note credits M.~E.~H.~Ismail, \emph{Classical and Quantum
Orthogonal Polynomials in One Variable} (2009), equation (5.9.3).
The following calculation proves the needed orthogonality in our own
coordinate and density conventions.
For real \(z,t\) close to zero, \eqref{a02gmf:eq:FCD12} and the cosine addition
formula give
\[
 \int_{\mathbb R}K(z,y)K(t,y)\sigma(y)\,dy
    =M_\sigma(1-zt)^{-1/2}.
\]
Indeed
\(\cos(\arctan z+\arctan t)
 =(1-zt)/\sqrt{(1+z^2)(1+t^2)}\), so all other factors cancel
exactly. The Gamma exponential bound permits coefficientwise
differentiation near zero. Hence
\begin{equation}
 \int p_n(y)p_m(y)\sigma(y)\,dy
       =\mathbf1_{\{n=m\}}M_\sigma n!(1/2)_n .
 \tag{FCD13}\label{a02gmf:eq:FCD13}
\end{equation}
The polynomials \(i^np_n((S-c)/i)\) are monic in \(S\), and their
values at \(c+iy\) are \(i^np_n(y)\). Their monic triangular map
has determinant one; their squared phases have modulus one.
It follows that
\begin{equation}
 \boxed{\displaystyle
 \det H_d[\sigma]
       =\prod_{j=0}^d\left[\sqrt{2\pi}\,j!(1/2)_j\right].}
 \tag{FCD14}\label{a02gmf:eq:FCD14}
\end{equation}

\subsection{Finite source-fill error}

Retain \(b=(\log M)'(\alpha)\), \(T=kb\), and the unchanged
piecewise source
\[
 \widetilde m_k(y)=
 \begin{cases}M_\alpha^k\sigma(y),&|y|\le T,\\
 m_k(y),&|y|>T.
 \end{cases}
\]
Use \(r_*=\tanh(\pi/8)\), \(\theta_*=\pi/4\),
\(\omega_*=\alpha-\theta_*=\pi/4\), and
\(\ell_*=-\log r_*\). These are fixed estimate parameters.
By the power series for \(\arctan z\),
\(|\arctan z|\le\operatorname{artanh}r_*=\pi/8\) on \(|z|=r_*\).
Cauchy's coefficient bound gives
\[
 |p_n(y)|\le n!r_*^{-n}(1-r_*^2)^{-1/4}
                                    e^{(\pi/8)|y|}.
\]
The largest of the \(2n+1\) binomial coefficients
\(\binom{2n}{j}\) is \(\binom{2n}{n}\), so
\(n!/(1/2)_n=4^n/\binom{2n}{n}\le2n+1\).
Expanding a polynomial in the orthogonal basis \eqref{a02gmf:eq:FCD13}
and applying Cauchy--Schwarz thus proves
\[
 |P(c+iy)|^2\le A_d e^{\theta_*|y|}\|P\|_\sigma^2,\qquad
 A_d=\frac{(d+1)^2e^{2d\ell_*}}
             {M_\sigma\sqrt{1-r_*^2}},
       \qquad \deg P\le d .
\]
The sum of \(2n+1\) for \(0\le n\le d\) is \((d+1)^2\);
this proves the displayed constant including its highest-degree factor.
Evenness and the exact convolution transform give
\(\int e^{\alpha|y|}m_k(y)\,dy\le2M_\alpha^k\).
Therefore
\[
 \int_{|y|>T}|P(c+iy)|^2m_k(y)\,dy
  \le2A_dM_\alpha^ke^{-\omega_*T}\|P\|_\sigma^2 .
\]
For the original finite Gamma envelope
\(\sigma(y)\le\overline c_\Gamma
e^{-\alpha|y|}(1+|y|)^{-1/2}\), let
\[
 \eta_{k,d}=\frac{2\overline c_\Gamma A_d}
                 {\omega_*\sqrt{1+T}}e^{-\omega_*T},
 \qquad t_{k,d}=2A_de^{-\omega_*T}.
\]
Integrating the Gamma envelope separately on the two exterior half-lines
proves
\(\int_{|y|>T}|P|^2\sigma\le\eta_{k,d}\|P\|_\sigma^2\).
It follows, for \(\eta_{k,d}<1\), that
\[
 M_\alpha^k(1-\eta_{k,d})H_d[\sigma]
  \preceq H_d[\widetilde m_k]
  \preceq M_\alpha^k(1+t_{k,d})H_d[\sigma].
\]
Define the exact source-fill error
\[
 r_{d,k}^{\rm fill}
   =\log\det H_d[\widetilde m_k]
        -(d+1)k\log M_\alpha-\log\det H_d[\sigma].
\]
Taking the positive eigenvalues of the relative forms proves
\begin{equation}
 \begin{split}
 (d+1)\log(1-\eta_{k,d})
       &\le r_{d,k}^{\rm fill}
                 \le(d+1)\log(1+t_{k,d}),\\
 |r_{d,k}^{\rm fill}|
       &\le(d+1)\{-\log(1-\eta_{k,d})+\log(1+t_{k,d})\}.
 \end{split}
 \tag{FCD15}\label{a02gmf:eq:FCD15}
\end{equation}
For each fixed \(d\), this bound decays exponentially in \(k\).
Strict convexity of the even source transform gives \(b>0\):
its logarithmic second derivative is a positive tilted variance, and
its derivative at zero vanishes.

\subsection{Exact receiver for the attained marked Gram}

Define the explicit expression
\begin{equation}
 \begin{split}
 \mathcal A_{k,d}={}&
 (d+1)kD_0-\frac{d(d+1)}2\log(kV)
 +\frac{d+1}2\log(2\pi)+\sum_{j=0}^d\log(1/2)_j\\
 &-\frac{c_4d(d^2-1)}{12k}
 -\frac{d(d^2-1)}{288k^2}
          [(16-11d)c_4^2+(2d-4)c_6].
 \end{split}
 \tag{FCD16}\label{a02gmf:eq:FCD16}
\end{equation}
All factors \(j!\) cancel between \eqref{a02gmf:eq:FCD5} and
\eqref{a02gmf:eq:FCD14}, whereas \(M_0\), \(M_\alpha\), \(V\),
\(\sqrt{2\pi}\), and \((1/2)_j\) remain exactly as displayed.
The source calculation proves the exact finite identity
\begin{equation}
 \log\frac{\det H_d[\widetilde m_k]}{\det H_d[m_k]}
       =\mathcal A_{k,d}+r_{d,k}^{\rm fill}-r_{d,k}^{\rm cum}.
 \tag{FCD17}\label{a02gmf:eq:FCD17}
\end{equation}

For the original minimum-remainder metric \(G_N[\nu]\), let
\(\mathscr F_{\lambda,d}
 =\operatorname{span}\{v_\lambda,1,S,\ldots,S^d\}\), with
\(v_\lambda=b_\lambda\chi(S)/(S-\lambda)\) and the original
complete-primary coefficient \(b_\lambda\).
For each \(q-1\le N\le2q\), define the exact attained transport difference
\[
 r_{N,k,d}^{\rm att}
  =\log\frac{\det G_N[\widetilde m_k]|_{\mathscr F_{\lambda,d}}}
                  {\det G_N[m_k]|_{\mathscr F_{\lambda,d}}}
    -\log\frac{\det H_d[\widetilde m_k]}{\det H_d[m_k]} .
\]
Thus the receiver is the exact equality
\begin{equation}
 \boxed{\displaystyle
 \log\frac{\det G_N[\widetilde m_k]|_{\mathscr F_{\lambda,d}}}
                  {\det G_N[m_k]|_{\mathscr F_{\lambda,d}}}
  =\mathcal A_{k,d}+r_{N,k,d}^{\rm att}
                          +r_{d,k}^{\rm fill}-r_{d,k}^{\rm cum}.}
 \tag{FCD18}\label{a02gmf:eq:FCD18}
\end{equation}
Equation (GMF26) in the accompanying growing-flag proof evaluates
\(|r_{N,k,d}^{\rm att}|\le\Omega_{k,d}\), with uniform decay in (GMF28),
through the actual relation
minimum, terminal minimum, and their complete Schur cross block.
After that proved estimate is inserted, the complete finite error in
\eqref{a02gmf:eq:FCD18} is
\begin{equation}
 \Omega_{k,d}
 +(d+1)\{-\log(1-\eta_{k,d})+\log(1+t_{k,d})\}
 +\frac{C_3+C_0^2+\tfrac23C_0^3}{k^3}.
 \tag{FCD19}\label{a02gmf:eq:FCD19}
\end{equation}
All finite guards of that localization estimate remain in force, in
addition to \(\eta_{k,d}<1\) and \(k\ge\max(1,2C_0)\).
The quantity \(t_{k,d}\) in this section is exactly \(\zeta_{k,d}\)
in (GMF9), and \eqref{a02gmf:eq:FCD15} is the determinant consequence
of (GMF10), with the source constants independently derived above.
This cumulant calculation alone makes no replacement of an attained
metric by its unquotiented source Gram: their exact difference is
the first error term of \eqref{a02gmf:eq:FCD18}.

The original marked derivative frame has squared change-of-column
determinant
\[
 |b_\lambda|^{2(d+1)}((d+1)!)^2
                  |u|^{-(d+1)(d+2)} .
\]
It is the same coefficient map for both source metrics. Hence it cancels
in their ratio, through the exact rule
\(\det(C^*HC)=|\det C|^2\det H\); it is not assigned the value one.
Consequently \eqref{a02gmf:eq:FCD18} receives the original marked derivative
Gram with precisely the same coefficients and finite error.

\subsection{The first two nontrivial fixed flags}

For \(d=1\), the exact moment matrix is \(\operatorname{diag}(1,1)\).
Thus \(P_1=1\), every cumulant error vanishes, and
\begin{equation}
 \mathcal A_{k,1}=2kD_0-\log(kV)+\log\pi .
 \tag{FCD20}\label{a02gmf:eq:FCD20}
\end{equation}
The exact difference between the attained three-column determinant
ratio and \eqref{a02gmf:eq:FCD20} is
\(r_{N,k,1}^{\rm att}+r_{1,k}^{\rm fill}\).
The accompanying localization estimate makes it exponentially small
uniformly over the original cutoff window.

For \(d=2\), the exact moment matrix and polynomial are
\[
 \begin{pmatrix}
 1&0&1\\
 0&1&0\\
 1&0&3+c_4/k
 \end{pmatrix},
 \qquad P_2(1/k)=1+\frac{c_4}{2k}.
\]
Here \(\prod_{j=0}^2(1/2)_j=3/8\), so
\begin{equation}
 \mathcal A_{k,2}
 =3kD_0-3\log(kV)+\tfrac32\log(2\pi)+\log(3/8)
                 -\frac{c_4}{2k}+\frac{c_4^2}{8k^2}.
 \tag{FCD21}\label{a02gmf:eq:FCD21}
\end{equation}
Equations \eqref{a02gmf:eq:FCD18}--\eqref{a02gmf:eq:FCD19} give its finite remainder.
In this case \(C_0=|c_4|/2\), \(C_3=0\); the direct scalar series also
gives the sharper bound
\[
 |r_{2,k}^{\rm cum}|
    \le\frac{|c_4|^3}{24k^3(1-|c_4|/(2k))}
       \quad\text{for }k>|c_4|/2 .
\]
The sign of the quadratic term in \eqref{a02gmf:eq:FCD21} is positive because
the arithmetic determinant contains \(+\log(1+c_4/(2k))\), which is
subtracted in the source ratio.

\paragraph{Exact-algebra evidence.}
The accompanying \texttt{check\_exact\_cumulants.py} uses rational and
symbolic SymPy arithmetic. It verifies the general polynomial trace
identities, evaluates all three logarithmic coefficients directly from
the original Gaussian moment matrices at each \(d=0,\ldots,20\), and
constructs the entire finite cumulant determinant polynomial for
\(d=0,\ldots,5\). It also checks the finite Gamma norm coefficients.
The run records 95 successful assertions. These are new local checks.
The incoming author's reported Wolfram fixtures were unavailable and
were not executed or verified here. The exact finite proofs above
establish the formulas for all fixed nonnegative integers \(d\);
the finite cases are supplementary checks. None of this fixed-degree
calculation substitutes \(d=d_k\) into a nonuniform remainder or
asserts an absolute observation allocation.

% END COMPLETE FCD SOURCE



% END PRESERVED INPUT 06_GROWING_MARKED_FLAGS.tex
\endgroup

\clearpage
\begingroup
\def\im{\operatorname{im}}
\def\diag{\operatorname{diag}}

\setcounter{equation}{0}
\renewcommand{\theequation}{A02NC.\arabic{equation}}
\section*{Full-class completion, original action and boundary-response dictionary}
\addcontentsline{toc}{section}{Full-class completion, original action and boundary-response dictionary}
% BEGIN PRESERVED INPUT 07_NATIVE_CLASS_INTEGRATION.tex SHA256 7e940ecd9ba8969bfdbf6d9890b937305b585306d2395c2e7748b45ded9135b3


\begin{abstract}
The actual DNE numerator section has a residue-dual description that omits
nonzero components of every retained simple eigenclass. We construct its
full low, conductor and native completion through every original minor,
calculate its metric and action with all cross terms, and carry the
complete class through the original arithmetic realization. The complete
centered action has an exact rank-two boundary formula at every original
cutoff. Its two boundary rows give precisely the original EC response
ratios, including their signs and both directed leakages. The parity of
the full physical source evaluates the mixed boundary Gram entry exactly
and identifies the product of its two observation amplifications. The
proof retains the original full root polynomial, measure, minimum,
conductor and projection. Absolute determinant allocations and individual
projected arithmetic signs are not evaluated by these finite identities.
\end{abstract}
The originating programme and its existing public attribution are
retained \cite{a02nc:human:globalization2026,a02nc:human:splitzero2026}.
The supplied continuation is \cite{a02nc:prog:arrival02}; source editions and
complete input hashes are recorded in \texttt{SOURCE\_PINS.json}.
All original provider files remain unchanged. The complete additive
finite calculations are written below; their original theta-source
realization is the pinned FCI/FDC provider, whose original hypotheses
are retained. No new certification of that provider's earlier analytic
estimates is implied by this finite intake.

% Additive complete algebra proof; no preamble or bibliography.
% Parent bibliography keys: prog:DNE, prog:arrival02.
\section{The original residue-dual basis, every conductor chart,
and the exact action-defect rank}

\subsection{Original objects and scope}
This calculation proves equations (3)--(17) of the supplied
native-class continuation \cite{a02nc:prog:arrival02} in the original DNE
coordinates \cite{a02nc:prog:DNE}. It also evaluates the action-defect rank
for every selected set and constructs the exact transition from every
original row minor to the consecutive chart. All residue expansions and
matrix eliminations used below are proved here.

Keep the actual simple-quartet data and the actual full conductor:
\begin{equation}\tag{NCB1}\label{a02nc:ncb:data}
\begin{gathered}
a=k+4,\quad q=(a+1)^2,\quad d=q'=(a-7)^2,\quad
v=\operatorname{ord}_0E_A,\quad h=q-v,\quad m=h-d>0,\\
\beta_{a,u,w}=\frac a2+(2u-a)\delta+i(2w-a)\gamma,
\quad 0\leq u,w\leq a,\quad 0<\delta<\tfrac12,\quad\gamma>2,\\
E_A(z)=\sum_{r,t=0}^{8}a_{rt}e^{\beta_{8,r,t}z},
\qquad \mu_j=E_A^{(j)}(0),\qquad
\mu_0=\cdots=\mu_{v-1}=0,\quad\mu_v\ne0 .
\end{gathered}
\end{equation}
Every actual period, branch, coefficient, and finite guard remains
unchanged. The lower roots are
$\beta'_{\alpha'}=\beta_{a-8,u',w'}$, with their original ordering.
There are $d$ lower roots and $q$ upper roots. All roots within either
grid are distinct: equality of real and imaginary parts forces equality
of both indices. Every upper root is nonzero since its real part is
at least $a(1/2-\delta)>0$. Set
\begin{equation}\tag{NCB2}\label{a02nc:ncb:chi}
\chi(S)=\prod_{\alpha=1}^{q}(S-\beta_\alpha)
       =\sum_{\ell=0}^{q}c_\ell S^\ell,\quad c_q=1,\quad
E=\mathbb C[S]/(\chi),\quad
\mathbb V_{n,\alpha}=\beta_\alpha^n\quad(0\leq n<q).
\end{equation}
Represent elements of $E$ by their unique remainders of degree below
$q$. For an ordered subset $B$ of coordinate indices, $E_B$ inserts
coordinates in increasing index order; $E_B^{\mathsf T}$ extracts them.
These are ordinary transpose conventions, not Hermitian adjoints.

The literal DNE conductor coefficient map
$C_a^{\mathsf T}:\mathbb C^d\to\mathbb C^q$ is
\begin{equation}\tag{NCB3}\label{a02nc:ncb:conductor}
\sum_\alpha(C_a^{\mathsf T}\eta)_\alpha e^{\beta_\alpha z}
=E_A(z)\sum_{\alpha'}\eta_{\alpha'}e^{\beta'_{\alpha'}z}.
\end{equation}
Its upper $(u,w)$, lower $(u',w')$ entry is
$a_{u-u',w-w'}$ when both differences lie in $\{0,\ldots,8\}$, and
zero otherwise, because
$\beta_{8,r,t}+\beta_{a-8,u',w'}=\beta_{a,u'+r,w'+t}$.
No divided-native scalar is inserted in this map.

The map $C_a^{\mathsf T}$ is injective. If its value is zero, the
product in \eqref{a02nc:ncb:conductor} vanishes identically. Since
$\mu_v\ne0$, $E_A$ is nonzero on an open disc; the lower exponential
polynomial vanishes on that disc and hence identically. Its first $d$
derivatives at zero form the lower Vandermonde times $\eta$. The
determinant is
$\prod_{\alpha'<\beta'}(\beta'_{\beta'}-\beta'_{\alpha'})\ne0$,
so $\eta=0$.
To verify this determinant formula, its determinant is an alternating
polynomial in the $d$ root variables, hence divisible by every
$\beta'_{\beta'}-\beta'_{\alpha'}$. Its total degree is
$0+1+\cdots+(d-1)$, the degree of their product. The coefficient of
$(\beta'_1)^0(\beta'_2)^1\cdots(\beta'_d)^{d-1}$ is one in both
polynomials, so their constant quotient is one.

Let $P_v:\mathbb C^h\to\mathbb C^q$ insert $v$ leading zero
coordinates. The full product rule applied to \eqref{a02nc:ncb:conductor}
gives
\begin{equation}\tag{NCB4}\label{a02nc:ncb:moment-conductor}
\begin{gathered}
\mathbb VC_a^{\mathsf T}=P_vJ,\qquad
J=P_v^{\mathsf T}\mathbb VC_a^{\mathsf T},\\
J_{i,\alpha'}=\sum_{r=0}^{i}
\binom{v+i}{r}\mu_{v+i-r}(\beta'_{\alpha'})^r
\qquad(0\leq i<h).
\end{gathered}
\end{equation}
All derivatives below order $v$ vanish. At order $v+i$, the product
rule originally ranges over $0\leq r\leq v+i$; terms with $r>i$
vanish because their conductor derivative has order below $v$.
Thus the displayed truncation retains every nonzero conductor
derivative. Since $\mathbb V$ and $P_v$ are injective, $J$ has rank $d$.

\subsection{Complete residue-dual moment basis}
For each upper root let
$p_\alpha(S)=\chi(S)/((S-\beta_\alpha)\chi'(\beta_\alpha))$.
Define the residue isomorphism and its moment-coordinate version:
\begin{equation}\tag{NCB5}\label{a02nc:ncb:residue}
\mathcal L:\mathbb C^q\longrightarrow E,\qquad
\mathcal L(\nu)=\sum_\alpha\chi'(\beta_\alpha)\nu_\alpha p_\alpha,
\qquad F=\mathcal L\mathbb V^{-1}.
\end{equation}
Evaluation at the roots gives
$(\mathcal L^{-1}f)_\alpha=f(\beta_\alpha)/\chi'(\beta_\alpha)$,
so $\mathcal L$ is invertible. The roots are simple and no derivative
factor vanishes. The formula for a simple-pole residue gives
\[
\sum_\alpha\operatorname{Res}_{S=\beta_\alpha}
\frac{\mathcal L(\nu)(S)P(S)}{\chi(S)}\,dS
=\sum_\alpha\nu_\alpha P(\beta_\alpha).
\]

Define, for $0\leq n<q$,
\begin{equation}\tag{NCB6}\label{a02nc:ncb:pn}
\mathfrak p_n(S)=\sum_{\ell=n+1}^{q}c_\ell S^{\ell-n-1}.
\end{equation}
Their distinct monic degrees are $q-1-n$, so they form a basis of $E$.
Their complete residue duality is
\begin{equation}\tag{NCB7}\label{a02nc:ncb:duality}
\sum_\alpha\operatorname{Res}_{S=\beta_\alpha}
\frac{\mathfrak p_n(S)S^j}{\chi(S)}\,dS=\delta_{nj}
\quad(0\leq j,n<q),\qquad Fe_n=\mathfrak p_n .
\end{equation}
To prove the first equality, the sum of finite residues equals the
$S^{-1}$ coefficient at large $S$: partial fractions consist of a
polynomial and terms $b_\alpha/(S-\beta_\alpha)$, and each such term
has both residue and $S^{-1}$ coefficient $b_\alpha$.
For $j<n$, the numerator has degree at most $q-2$, so this coefficient
is zero. For $j=n$, it has degree $q-1$ and leading coefficient one.
For $j>n$, the exact identity
$S^{n+1}\mathfrak p_n=\chi-\sum_{\ell=0}^{n}c_\ell S^\ell$
implies
\[
\frac{\mathfrak p_nS^j}{\chi}
=S^{j-n-1}
-\frac{S^{j-n-1}\sum_{\ell=0}^{n}c_\ell S^\ell}{\chi}.
\]
The first term is a polynomial, and the second numerator has degree
at most $j-1\leq q-2$. Both have zero $S^{-1}$ coefficient.
This proves the first equality. The pairings of $Fe_n$ with $S^j$
are $(\mathbb V\mathbb V^{-1}e_n)_j=\delta_{nj}$, so the proved
basis duality identifies $Fe_n$ with $\mathfrak p_n$.

Fix any actual $d$-row minor
$I\subset\{0,\ldots,h-1\}$ with $\det J_I\ne0$,
where $J_I=E_I^{\mathsf T}J$. Write $J(I)=I^c$ and
\begin{equation}\tag{NCB8}\label{a02nc:ncb:native}
T_I=\{v+j:j\in J(I)\},\qquad
Z_I=\mathbb V^{-1}P_vE_{J(I)},\qquad
L_I=\mathcal LZ_I=FE_{T_I}.
\end{equation}
Hence $\operatorname{im}L_I$ is exactly the span of
$\{\mathfrak p_n:n\in T_I\}$ in the original selected coordinates.

For every upper root $\lambda$, the finite geometric sums give
\[
\sum_{n=0}^{q-1}\lambda^n\mathfrak p_n(S)
=\sum_{\ell=1}^{q}c_\ell
\frac{S^\ell-\lambda^\ell}{S-\lambda}
=\frac{\chi(S)-\chi(\lambda)}{S-\lambda}.
\]
Consequently the original eigenclass is
\begin{equation}\tag{NCB9}\label{a02nc:ncb:eigenclass}
v_\lambda=\frac{\chi(S)}{(S-\lambda)\chi'(\lambda)}
=\sum_{n=0}^{q-1}t_n\mathfrak p_n,\qquad
t_n=\frac{\lambda^n}{\chi'(\lambda)},\qquad
v_\lambda\notin\operatorname{im}L_I.
\end{equation}
Every $t_n$ is nonzero because $\lambda\ne0$ and
$\chi'(\lambda)\ne0$. A proper coordinate subset of the displayed
basis cannot contain a vector whose every coordinate is nonzero.
This proves the last assertion for every original admitted minor.

\subsection{Complete low, conductor, and native decomposition}
For an arbitrary $f\in E$, put $t=F^{-1}f$ and
$t_{\geq v}=P_v^{\mathsf T}t$. Define maps with their full types:
\begin{equation}\tag{NCB10}\label{a02nc:ncb:HQ}
\begin{gathered}
H_I=J_I^{-1}E_I^{\mathsf T}:\mathbb C^h\to\mathbb C^d,\qquad
Q_I=E_{J(I)}^{\mathsf T}-J_{J(I)}H_I:
       \mathbb C^h\to\mathbb C^m,\\
\eta_I=H_It_{\geq v},\qquad z_I=Q_It_{\geq v}.
\end{gathered}
\end{equation}
Checking first the $I$ rows and then their complement proves
\begin{equation}\tag{NCB11}\label{a02nc:ncb:split}
\begin{gathered}
JH_I+E_{J(I)}Q_I=I_h,\qquad H_IJ=I_d,\qquad Q_IJ=0,\\
H_IE_{J(I)}=0,\qquad Q_IE_{J(I)}=I_m.
\end{gathered}
\end{equation}
Together with $FP_vJ=\mathcal LC_a^{\mathsf T}$ from
\eqref{a02nc:ncb:moment-conductor}, these give the full identity
\begin{equation}\tag{NCB12}\label{a02nc:ncb:decomposition}
f=\underbrace{\sum_{n<v}t_n\mathfrak p_n}_{f^{\rm low}}
 +\underbrace{\mathcal L(C_a^{\mathsf T}\eta_I)}_{f_I^{\rm cond}}
 +\underbrace{L_Iz_I}_{f_I^{\rm nat}} .
\end{equation}
These three spaces form a direct sum. The low space has arbitrary
first $v$ moment coordinates and zero remaining coordinates; both
other spaces have their first $v$ coordinates zero. On the remaining
coordinates, \eqref{a02nc:ncb:split} proves
$\mathbb C^h=\operatorname{im}J\oplus\operatorname{im}E_{J(I)}$.
Thus the full basis map
\begin{equation}\tag{NCB13}\label{a02nc:ncb:full-chart}
B_I=[\,FE_{\{0,\ldots,v-1\}},\ \mathcal LC_a^{\mathsf T},\ L_I\,]:
\mathbb C^v\oplus\mathbb C^d\oplus\mathbb C^m\longrightarrow E
\end{equation}
is invertible, also when its first block is empty ($v=0$).

For $f=v_\lambda$, all selected coordinates of $t_{\geq v}$ are
nonzero and $d\geq1$, so $\eta_I\ne0$. Injectivity of the full
conductor map and of $\mathcal L$ gives
$v_{\lambda,I}^{\rm cond}\ne0$. If $v>0$ then
$v_\lambda^{\rm low}\ne0$ as well. The direct sum just proved implies
\begin{equation}\tag{NCB14}\label{a02nc:ncb:nonzero}
v_\lambda^{\rm low}+v_{\lambda,I}^{\rm cond}\ne0 .
\end{equation}
This is an algebraic direct sum, not an orthogonal splitting.

The exact relative quotient represented by the native section is
\begin{equation}\tag{NCB15}\label{a02nc:ncb:relative}
\frac{E}
 {\operatorname{im}(FE_{\{0,\ldots,v-1\}})
          +\operatorname{im}(\mathcal LC_a^{\mathsf T})}
\ \xrightarrow[\ [f]\mapsto Q_IP_v^{\mathsf T}F^{-1}f\ ]{\ \cong\ }
\mathbb C^m .
\end{equation}
Its inverse is $z\mapsto[L_Iz]$. Its kernel and both inverse identities
follow from \eqref{a02nc:ncb:split}. This proves the exact map relating the
full arithmetic class and its native representative; it retains their
proved nonzero difference in $E$.

\subsection{Every original minor to every other minor}
For another admitted $d$-row minor $K$, let
\begin{equation}\tag{NCB16}\label{a02nc:ncb:transition}
D_{K\leftarrow I}=H_KE_{J(I)}:\mathbb C^m\to\mathbb C^d,\qquad
M_{K\leftarrow I}=Q_KE_{J(I)}:\mathbb C^m\to\mathbb C^m .
\end{equation}
Apply \eqref{a02nc:ncb:split} for $K$ to $E_{J(I)}$ to obtain
\begin{equation}\tag{NCB17}\label{a02nc:ncb:section-transition}
\begin{gathered}
E_{J(I)}=JD_{K\leftarrow I}
              +E_{J(K)}M_{K\leftarrow I},\\
L_I=\mathcal LC_a^{\mathsf T}D_{K\leftarrow I}
              +L_KM_{K\leftarrow I}.
\end{gathered}
\end{equation}
This is the full changed-representative identity with its conductor
correction. Apply $Q_I$ to its first line to obtain
$M_{I\leftarrow K}M_{K\leftarrow I}=I_m$, and exchange the minors
to obtain the reverse product. Therefore
\begin{equation}\tag{NCB18}\label{a02nc:ncb:inverse}
M_{K\leftarrow I}^{-1}=M_{I\leftarrow K},\qquad
\eta_K=\eta_I+D_{K\leftarrow I}z_I,\quad
z_K=M_{K\leftarrow I}z_I,\quad
f_K^{\rm low}=f_I^{\rm low}.
\end{equation}
The component identities follow by applying $H_K,Q_K$ to
$t_{\geq v}=J\eta_I+E_{J(I)}z_I$. In particular
\[
f_I^{\rm nat}=f_K^{\rm nat}
       +\mathcal LC_a^{\mathsf T}D_{K\leftarrow I}z_I,\qquad
f_K^{\rm cond}=f_I^{\rm cond}
       +\mathcal LC_a^{\mathsf T}D_{K\leftarrow I}z_I.
\]
Their total is unchanged because the same correction occurs on
opposite sides when the decompositions are compared.

For a third admitted minor $H$, applying $Q_H,H_H$ to
\eqref{a02nc:ncb:section-transition} proves the oriented composition laws
\begin{equation}\tag{NCB19}\label{a02nc:ncb:cocycle}
\begin{gathered}
M_{H\leftarrow I}=M_{H\leftarrow K}M_{K\leftarrow I},\\
D_{H\leftarrow I}
 =D_{K\leftarrow I}+D_{H\leftarrow K}M_{K\leftarrow I}.
\end{gathered}
\end{equation}

The exact determinant including ordering signs is also explicit.
For $I=(i_0<\cdots<i_{d-1})$ put
$\epsilon_I=(-1)^{\sum_{r=0}^{d-1}(i_r-r)}$.
The row permutation taking full increasing order to $(I,J(I))$
has this sign, since selected row $i_r$ crosses $i_r-r$ earlier
unselected rows. Under that permutation $W_I=[J,E_{J(I)}]$ is
block lower triangular with diagonal blocks $J_I,I_m$, so
$\det W_I=\epsilon_I\det J_I$. The first line of
\eqref{a02nc:ncb:section-transition} gives
$W_I=W_K\left[\begin{smallmatrix}I_d&D\\0&M\end{smallmatrix}\right]$.
Consequently
\begin{equation}\tag{NCB20}\label{a02nc:ncb:dettransition}
\det M_{K\leftarrow I}
 =\frac{\epsilon_I\det J_I}{\epsilon_K\det J_K}.
\end{equation}

The full coordinate transformation, including the unchanged low
coordinates, is
\begin{equation}\tag{NCB21}\label{a02nc:ncb:full-transition}
T_{K\leftarrow I}=
\begin{pmatrix}I_v&0&0\\0&I_d&D_{K\leftarrow I}\\
0&0&M_{K\leftarrow I}\end{pmatrix},
\qquad B_I=B_KT_{K\leftarrow I}.
\end{equation}
For every original arithmetic Hermitian metric $G_N$ and the full
operator $A$ of multiplication by $S$ modulo $\chi$, set
$G_{I,N}=B_I^*G_NB_I$ and $A_I=B_I^{-1}AB_I$. Substitution yields
\begin{equation}\tag{NCB22}\label{a02nc:ncb:gram-action}
G_{I,N}=T_{K\leftarrow I}^*G_{K,N}T_{K\leftarrow I},\qquad
A_I=T_{K\leftarrow I}^{-1}A_KT_{K\leftarrow I}.
\end{equation}
All nine Gram blocks, including all six off-diagonal pairings, are
retained. Multiplication of these identities gives
\begin{equation}\tag{NCB23}\label{a02nc:ncb:centered-action}
\begin{split}
A_I^*G_{I,N}+G_{I,N}A_I-aG_{I,N}
={}&T_{K\leftarrow I}^*
(A_K^*G_{K,N}+G_{K,N}A_K-aG_{K,N})T_{K\leftarrow I}.
\end{split}
\end{equation}
Thus the full arithmetic action is preserved under the original
coordinate transition with its explicit conductor corrections.

\subsection{Consecutive chart with all conductor derivatives}
Let $I_0=\{0,\ldots,d-1\}$. Equation \eqref{a02nc:ncb:moment-conductor}
gives the exact factorization
\begin{equation}\tag{NCB24}\label{a02nc:ncb:triangular}
J_{I_0}=\mathsf T\mathbb V',\qquad
\mathsf T_{i,r}=
\begin{cases}\binom{v+i}{r}\mu_{v+i-r},&0\leq r\leq i<d,\\
0,&r>i,\end{cases}
\quad \mathbb V'_{r,\alpha'}=(\beta'_{\alpha'})^r .
\end{equation}
The lower-triangular diagonal entries are
$\binom{v+i}{i}\mu_v$. Taking the determinant of the actual product
therefore proves
\begin{equation}\tag{NCB25}\label{a02nc:ncb:consecutive-det}
\det J_{I_0}
=\mu_v^d\left(\prod_{i=0}^{d-1}\binom{v+i}{i}\right)
 \prod_{\alpha'<\beta'}(\beta'_{\beta'}-\beta'_{\alpha'})\ne0 .
\end{equation}
All $\mu_j$ with $j>v$ remain in the matrix. Their absence from its
determinant is the proved effect of triangularity.

Its inverse retains every such derivative. For $x\in\mathbb C^d$,
solve $\mathsf T b=x$ by the finite recursion
\begin{equation}\tag{NCB26}\label{a02nc:ncb:triangular-inverse}
b_0=\frac{x_0}{\mu_v},\qquad
b_i=\frac{x_i-\sum_{r=0}^{i-1}
          \binom{v+i}{r}\mu_{v+i-r}b_r}
          {\binom{v+i}{i}\mu_v}\quad(1\leq i<d).
\end{equation}
Write
$\chi_{\rm low}(S)=\prod_{\alpha'}(S-\beta'_{\alpha'})
=\sum_{\ell=0}^{d}c_\ell^{\rm low}S^\ell$
and
$\mathfrak p_r^{\rm low}(S)
=\sum_{\ell=r+1}^{d}c_\ell^{\rm low}S^{\ell-r-1}$.
The already proved residue-dual basis applied to $\chi_{\rm low}$
evaluates the inverse lower Vandermonde, giving
\begin{equation}\tag{NCB27}\label{a02nc:ncb:full-inverse}
(J_{I_0}^{-1}x)_{\alpha'}
=\frac{\sum_{r=0}^{d-1}b_r\mathfrak p_r^{\rm low}(\beta'_{\alpha'})}
          {\chi_{\rm low}'(\beta'_{\alpha'})}.
\end{equation}
Indeed the vector on the right has lower moments $b$ by
\eqref{a02nc:ncb:duality}, and those moments satisfy $\mathsf T b=x$.

For an arbitrary original chart $I$ and its native vector $z$, take
$x=E_{I_0}^{\mathsf T}E_{J(I)}z$ in
\eqref{a02nc:ncb:triangular-inverse}--\eqref{a02nc:ncb:full-inverse}. Their result
is exactly $D_{I_0\leftarrow I}z$. The native coordinates in the
consecutive chart are, for $d\leq j<h$,
\begin{equation}\tag{NCB28}\label{a02nc:ncb:explicit-transition}
\begin{split}
(M_{I_0\leftarrow I}z)_{j-d}
={}&(E_{J(I)}z)_j\\
&-\sum_{\alpha'}\sum_{r=0}^{j}
\binom{v+j}{r}\mu_{v+j-r}(\beta'_{\alpha'})^r
                    (D_{I_0\leftarrow I}z)_{\alpha'} .
\end{split}
\end{equation}
This evaluates the chart transition in original coefficients, roots,
and derivatives. For the full eigenclass one instead sets
$x_i=\lambda^{v+i}/\chi'(\lambda)$ to calculate $\eta_{I_0}$,
and then sets
$z_{I_0,j-d}=t_{v+j}-(J\eta_{I_0})_j$.

The selected consecutive set is
$T_{I_0}=\{v+d,\ldots,q-1\}=\{q-m,\ldots,q-1\}$.
Its residue-dual basis elements have all monic degrees
$m-1,\ldots,0$, so their triangular coefficient matrix proves
\begin{equation}\tag{NCB29}\label{a02nc:ncb:consecutive-image}
\operatorname{im}L_{I_0}=\mathbb C[S]_{\leq m-1}\subset E .
\end{equation}

\subsection{Action defect and exact rank for every selected set}
Multiplication by $S$ modulo the complete $\chi$ satisfies
\begin{equation}\tag{NCB30}\label{a02nc:ncb:action}
A\mathfrak p_0=-c_0\mathfrak p_{q-1},\qquad
A\mathfrak p_n=\mathfrak p_{n-1}-c_n\mathfrak p_{q-1}
\quad(1\leq n<q).
\end{equation}
For the first equality use $S\mathfrak p_0=\chi-c_0$;
for the second use the literal polynomial identity
$S\mathfrak p_n=\mathfrak p_{n-1}-c_n$.
Here $\mathfrak p_{q-1}=1$.
For any selected $T\subset\{0,\ldots,q-1\}$, the omitted-coordinate
action is therefore
\begin{equation}\tag{NCB31}\label{a02nc:ncb:RT}
\begin{gathered}
R_T=E_{T^c}^{\mathsf T}F^{-1}AFE_T,\\
(R_T)_{\ell n}
=\mathbf1_{\{n\geq1,\ \ell=n-1\}}
       -c_n\mathbf1_{\{\ell=q-1\}}
\qquad(\ell\in T^c,\ n\in T).
\end{gathered}
\end{equation}

Define the descending starts and their number by
$B(T)=\{n\in T:n\geq1,\ n-1\notin T\}$ and $b(T)=|B(T)|$.
The exact rank, including nongeneric vanishing coefficients, is
\begin{equation}\tag{NCB32}\label{a02nc:ncb:rank}
\begin{gathered}
\operatorname{rank}R_T=b(T)+\varepsilon(T),\\
\varepsilon(T)=
\begin{cases}
1,&q-1\notin T\ \text{and some }n\in T\setminus B(T)
                                    \text{ has }c_n\ne0,\\
0,&\text{otherwise}.
\end{cases}
\end{gathered}
\end{equation}
For each $n\in B(T)$, row $n-1$ has exactly a single one in column
$n$. These $b(T)$ rows are independent and none is row $q-1$.
All other rows vanish except possibly row $q-1$, present exactly when
$q-1\notin T$. That row is $(-c_n)_{n\in T}$.
Subtracting its coordinates on columns $B(T)$ using the independent
rows leaves its coordinates on $T\setminus B(T)$; these add one rank
precisely in the case displayed. This proves \eqref{a02nc:ncb:rank}.

For every nonempty proper $T$, the rank is positive. If $B(T)$ is
nonempty, this follows immediately. Otherwise descending from any
member of $T$ cannot leave $T$ before reaching zero, so $T$ is an
initial interval containing zero. Properness implies $q-1\notin T$,
and $c_0=(-1)^q\prod_\alpha\beta_\alpha\ne0$ gives
$\varepsilon(T)=1$. For the consecutive set $T_{I_0}$ there is just
one descending start $q-m=v+d>0$, and $q-1\in T_{I_0}$.
Thus its rank is exactly one.

\subsection{Original metric residual and its full Gram}
Let $G_N$ be any original positive Hermitian arithmetic quotient metric
on $E$, fix $T=T_I$, and set $L=FE_T$ and $\Gamma=L^*G_NL$.
The metric compression and residual, both with domain $\mathbb C^m$,
are
\begin{equation}\tag{NCB33}\label{a02nc:ncb:compression}
A_{\rm nat,N}=\Gamma^{-1}L^*G_NAL,\qquad
\mathcal R_N=AL-LA_{\rm nat,N}.
\end{equation}
Write $\widetilde G=F^*G_NF$, and use $C=T^c$ solely as a block
index in the next formulas. Set
\begin{equation}\tag{NCB34}\label{a02nc:ncb:lift}
\begin{gathered}
W_G=F(E_C-E_T\widetilde G_{TT}^{-1}\widetilde G_{TC}),\\
S_G=\widetilde G_{CC}
 -\widetilde G_{CT}\widetilde G_{TT}^{-1}\widetilde G_{TC}.
\end{gathered}
\end{equation}
This leaves the original $G_N$ unchanged. Direct multiplication gives
$L^*G_NW_G=0$ and $E_C^{\mathsf T}F^{-1}W_G=I_{q-m}$.
Hence $W_G$ is an injective map onto the $G_N$-orthogonal complement
of $\operatorname{im}L$, whose dimension is $q-m$; its inverse there
is literal omitted-coordinate extraction.
By \eqref{a02nc:ncb:compression}, $\mathcal R_N$ has image in that
complement and omitted coordinates $R_T$. Uniqueness of this exact
lift proves
\begin{equation}\tag{NCB35}\label{a02nc:ncb:residual}
\mathcal R_N=W_GR_T,\qquad
\mathcal R_N^*G_N\mathcal R_N=R_T^*S_GR_T,\qquad
\operatorname{rank}\mathcal R_N=\operatorname{rank}R_T .
\end{equation}
For the Gram identity, multiply \eqref{a02nc:ncb:lift} to obtain
$W_G^*G_NW_G=S_G$, retaining its cross terms.
It is positive definite since $W_G$ is injective and $G_N$ is
positive definite. This proves the residual rank at every original
cutoff, with its exact value \eqref{a02nc:ncb:rank}, and proves nonvanishing
for every original nonempty proper chart. No size estimate follows
from rank alone.

\subsection{What this finite calculation establishes}
The residue-dual basis, numerator image, every eigenclass
decomposition, nonzero omitted component, original companion action,
and complete consecutive determinant prove equations (3)--(17) of
\cite{a02nc:prog:arrival02}. Equations
\eqref{a02nc:ncb:transition}--\eqref{a02nc:ncb:centered-action} provide the exact
comparison between every original minor and the consecutive chart.
The derivative recursion \eqref{a02nc:ncb:triangular-inverse} and
\eqref{a02nc:ncb:full-inverse} evaluate this comparison without deleting
conductor derivatives. The rank formula \eqref{a02nc:ncb:rank} and residual
Gram \eqref{a02nc:ncb:residual} strengthen the finite action calculation.

Under the programme's injective arithmetic realization, the nonzero
omitted vector \eqref{a02nc:ncb:nonzero} remains nonzero. Its map from the
full remainder space to the native relative quotient is precisely
\eqref{a02nc:ncb:relative}; changing the section is precisely
\eqref{a02nc:ncb:section-transition}. The arithmetic pairing of the full
class consequently uses the complete block congruence
\eqref{a02nc:ncb:gram-action} and \eqref{a02nc:ncb:centered-action}.
This calculation does not identify the native section with the
original observation kernel or evaluate an absolute allocation or
individual projected sign.


% Complete actual-class completion, 19 September 2026.
% Original DNE and FCI sources are unchanged.
\section{The full arithmetic class, its joint Gram, and its action}

Retain the simple quartet, the original period and branch, the full
nonzero conductor $E_A$, and the original full arithmetic unit.
Write $a=k+4$, $k=4\ell+1\ge9$,
$q=(a+1)^2$, $q'=(a-7)^2$, $v=\operatorname{ord}_0E_A\le80$,
$m=q-q'-v$, $c=a/2$, and
\[
 \beta_{a,r,t}=c+(2r-a)\delta+i(2t-a)\gamma,\qquad
 0<\delta<\tfrac12,\quad \gamma>2,\qquad
 \chi(S)=\prod_{r,t=0}^{a}(S-\beta_{a,r,t}).
 \tag{NCRP1}
\]
In particular $q'>0$, $m>0$, every root is simple, and
$\Re\beta\ge a(1/2-\delta)>0$. No root or conductor factor is
removed. The residue-dual polynomial frame and original DNE minor
are those proved in the preceding section. Let
$\mathscr P:\mathbb C^q\to E=\mathbb C[S]/(\chi)$ send the standard
column $e_n$ to $\mathfrak p_n$.
Thus $\mathscr P=\mathcal L\mathbb V^{-1}$, exactly in FCI33--34,
and multiplication by $S$ has matrix $C$ in this frame:
\[
 C_{\ell n}=\mathbf1_{\{n\ge1,\ell=n-1\}}
              -c_n\mathbf1_{\{\ell=q-1\}},\qquad
 \chi(S)=\sum_{n=0}^{q}c_nS^n,\quad c_q=1.
 \tag{NCRP2}
\]
This use of $C$ denotes the full packet action matrix; the original
conductor coefficient matrix continues to be written $C_a$.

\subsection{A square completion using the actual minor}

Let $E_{<v}:\mathbb C^v\to\mathbb C^q$ insert the first $v$
coordinates and retain the original $P_v,J_v,\mathcal I,\mathcal J$.
Define the square coefficient map
\[
 K_{\mathcal I}=
 [E_{<v},\,P_vJ_v,\,P_vE_{\mathcal J}],\qquad
 F_{\mathcal I}=\mathscr P K_{\mathcal I}=[F_0,F_C,L_a].
 \tag{NCRP3}
\]
Here $F_0$ has $v$ columns, $F_C=\mathcal L C_a^{\mathsf T}$
has $q'$ columns, and $L_a=\mathcal L Z_{\exp}$ has $m$ columns.
All three maps have codomain the original polynomial packet $E$.
For a moment column $t=(t_{<v},t_{\ge v})$ the inverse is explicitly
\[
 K_{\mathcal I}^{-1}t=
 \begin{pmatrix}
 t_{<v}\\
 (J_v)_{\mathcal I}^{-1}(t_{\ge v})_{\mathcal I}\\
 (t_{\ge v})_{\mathcal J}
       -(J_v)_{\mathcal J}(J_v)_{\mathcal I}^{-1}
                                (t_{\ge v})_{\mathcal I}
 \end{pmatrix}.
 \tag{NCRP4}
\]
Indeed substitution in the $\mathcal I$ and $\mathcal J$ rows,
and then in the first $v$ rows, gives exactly $t$.
Thus $K_{\mathcal I}$ and $F_{\mathcal I}$ are isomorphisms.
This proves a full arithmetic representative map, not only a
relative quotient map of dimension $m$.

For a retained simple eigenclass $v_\lambda$ let
$t_{\lambda,n}=\lambda^n/\chi'(\lambda)$ and
\[
 w_\lambda=(u_0,\eta,z)^{\mathsf T}
                  =K_{\mathcal I}^{-1}t_\lambda .
 \tag{NCRP5}
\]
The preceding residue calculation gives the exact equality
$v_\lambda=F_0u_0+F_C\eta+L_az$.
The vector $\eta$ is nonzero, because $q'>0$, all entries of
$(t_{\lambda,\ge v})_{\mathcal I}$ are nonzero, and the retained
minor is invertible. If $v>0$, $u_0$ is also nonzero.
Since $F_{\mathcal I}$ is injective, the omitted vector
$F_0u_0+F_C\eta$ is nonzero. These assertions hold for the actual
minor selected in DNE; they do not replace it by the consecutive chart.

\subsection{The original full-unit arithmetic injection}

The map into theta cohomology retains the source $F_h$ of FCI35:
$\mathcal M F_h=v_h=2\xi/h$ and $h(D)F_h=\Theta\phi_*$.
On the simple-quartet branch,
$h(s)=\prod_{\alpha\in\{\rho,\bar\rho,1-\rho,1-\bar\rho\}}(s-\alpha)$.
For $P\in E$ set, with the original Euler variables,
\[
 \iota_a(P)=
 [r_\chi(P)(D_1+\cdots+D_a)F_h^{\otimes a}]
            \in Q_\theta^{\otimes a}.
 \tag{NCRP6}
\]
Here $r_\chi$ is the increasing-degree representative below $q$.
We recall the complete finite reason that this packet map is injective.
In the tensor algebra
$\mathbb C[s_1,\ldots,s_a]/(h(s_1),\ldots,h(s_a))$,
evaluation at all ordered root tuples is an isomorphism: apply the
four-root Vandermonde separately in every coordinate.
Every pair of counts $(r,t)$, $0\le r,t\le a$, occurs among such
tuples, since the real-sign and imaginary-sign choices can be made
independently in each tensor leg. The resulting distinct sums are
exactly the roots in NCRP1. Hence $P(s_1+\cdots+s_a)$ vanishes
in that tensor algebra precisely when $\chi$ divides $P$.
The cyclic annihilator is therefore the full $\chi$.

The original Taylor map on the theta quotient sends NCRP6 to
$\upsilon_h^{\otimes a}P(A_h^{(a)})1$.
Every component of the original primary unit $\upsilon_h=j_h(v_h)$
is nonzero, so multiplication by it is invertible. It cannot turn
a nonzero packet class into zero.
For completeness, multiples of $\chi(s_1+\cdots+s_a)$ lie in the
ideal $(h(s_1),\ldots,h(s_a))$ by the preceding tensor evaluation.
Successive monic division in $s_1,\ldots,s_a$ writes such a polynomial
as $\sum_i h(s_i)R_i$. Substitution of the original differential
operators supplies its theta primitive, leg by leg, with the original
$(-1)^{i-1}$ tensor-differential signs. Thus the representative
independence and the Taylor injectivity apply to the same map.
This is the FCI35/FDC construction \cite{a02nc:prog:FCI}, with its full
unit and all ordered tuples retained.

Consequently
\[
 \iota_a(v_\lambda)
 =\iota_a(F_0u_0)+\iota_a(F_C\eta)+\iota_a(L_az),
 \qquad
 \iota_a(F_0u_0+F_C\eta)\ne0.
 \tag{NCRP7}
\]
Taking four independent copies of these maps preserves every
equation. The omitted summands cannot be removed as theta boundaries.

\subsection{The full original arithmetic metric}

Let $dm_a^{\rm ar}(y)=m_a^{\rm ar}(y)\,dy$,
$m_a^{\rm ar}=w_h^{*a}$, and
$w_h(y)=|v_h(1/2+iy)|^2/(2\pi)$.
The source has its complete mass $(\int_{\mathbb R}w_h)^a$.
Let $H_N^{\rm ar}$ be its moment Gram in the increasing powers of
$S=c+iy$. For the literal remainder map
$J_N^{\chi}:\mathbb C[S]_{\le N}\to E$ and $N\ge q-1$, set
\[
 G_N=(J_N^\chi(H_N^{\rm ar})^{-1}(J_N^\chi)^*)^{-1},
 \qquad
 T_N=(H_N^{\rm ar})^{-1}(J_N^\chi)^*G_N.
 \tag{NCRP8}
\]
The norm of a class is the minimum over its full relation fibre.
To verify the formula, $J_N^\chi T_N=I$ and
$(T_Nx)^*H_N^{\rm ar}r=0$ whenever $J_N^\chi r=0$.
Every other lift is $T_Nx+r$, so its squared norm is
$x^*G_Nx+r^*H_N^{\rm ar}r$. This proves both the minimum and
$T_N^*H_N^{\rm ar}T_N=G_N$ with the whole relation space
$\chi\mathbb C[S]_{\le N-q}$ retained. It is the classical positive
Schur calculation \cite{a02nc:human:boyd-vandenberghe}, proved here in
its complex Hermitian coordinates.

All polynomial moments are finite for the admitted arithmetic source.
Its density is positive almost everywhere: the nonzero analytic
function $v_h$ on the physical line has isolated zeros; hence
$w_h>0$ almost everywhere, and its convolution powers are positive.
It follows that every nonzero polynomial has positive squared norm.
Thus all displayed moment and quotient forms are positive definite.
The original physical coordinate, both integration half-lines, and
the full root polynomial have remained in NCRP8.

Write $\mathbf F_{\mathcal I}$ for the increasing-power coefficient
matrix of $F_{\mathcal I}$. Its complete Gram and action are
\[
 \mathsf H_N=\mathbf F_{\mathcal I}^*G_N\mathbf F_{\mathcal I},
 \qquad
 \mathsf A=K_{\mathcal I}^{-1}CK_{\mathcal I}.
 \tag{NCRP9}
\]
The map $F_{\mathcal I}$ intertwines $\mathsf A$ and multiplication
by $S$ exactly, so
$\mathsf A w_\lambda=\lambda w_\lambda$.
Partition the blocks by $0,C,Z$, of sizes $v,q',m$. The actual
squared norm of the retained class is
\[
 \begin{split}
 \|v_\lambda\|_{G_N}^2={}&
 u_0^*\mathsf H_{00}u_0+\eta^*\mathsf H_{CC}\eta
       +z^*\mathsf H_{ZZ}z\\
 &+2\Re\{u_0^*\mathsf H_{0C}\eta
              +u_0^*\mathsf H_{0Z}z
              +\eta^*\mathsf H_{CZ}z\}.
 \end{split}                                                     \tag{NCRP10}
\]
These are the three diagonal terms and all six directed cross
pairings. There is no orthogonality claim about the DNE decomposition.

The full centered action has blocks
\[
 \mathsf D_{ij}
 =\sum_{r\in\{0,C,Z\}}\mathsf A_{ri}^*\mathsf H_{rj}
  +\sum_{r\in\{0,C,Z\}}\mathsf H_{ir}\mathsf A_{rj}
  -a\mathsf H_{ij}.
 \tag{NCRP11}
\]
Direct multiplication gives
$\mathsf D=\mathsf A^*\mathsf H+\mathsf H\mathsf A-a\mathsf H$.
For $\lambda=a\rho=c+a\delta+ia\gamma$, the original same-class
Hermitian action therefore obeys the exact full receiver
\[
 \begin{split}
 &u_0^*\mathsf D_{00}u_0+\eta^*\mathsf D_{CC}\eta
       +z^*\mathsf D_{ZZ}z\\
 &\qquad+
 2\Re\{u_0^*\mathsf D_{0C}\eta+
                  u_0^*\mathsf D_{0Z}z+
                  \eta^*\mathsf D_{CZ}z\}
       =2a\delta\, w_\lambda^*\mathsf H_Nw_\lambda>0.
 \end{split}                                                     \tag{NCRP12}
\]
The proof substitutes $\mathsf Aw_\lambda=\lambda w_\lambda$;
both terms contribute their full $\lambda,\bar\lambda$.
This is a statement on the complete retained class. It has not
assigned a sign to one of its summands.

\subsection{Orthogonal completion with every original cross block}

It is useful to obtain an exact minimum-distance formula without
changing the original minor. Put
$F_\perp^{\,0}=[F_0,F_C]$ and $u=(u_0,\eta)^{\mathsf T}$.
In the coefficient formulas below the maps are expressed in the
same increasing-power frame used for $G_N$. Define
\[
 \begin{gathered}
 \Gamma=L_a^*G_NL_a,\qquad
 H_{ZC}=L_a^*G_NF_\perp^{\,0},\qquad
 H_{CC}=(F_\perp^{\,0})^*G_NF_\perp^{\,0},\\
 P_Z=L_a\Gamma^{-1}L_a^*G_N,\qquad
 E_C=(I-P_Z)F_\perp^{\,0},\\
 \Sigma=H_{CC}-H_{ZC}^*\Gamma^{-1}H_{ZC},\qquad
 \widehat z=z+\Gamma^{-1}H_{ZC}u .
 \end{gathered}                                                   \tag{NCRP13}
\]
Completing the positive square proves that $\Sigma\succ0$:
for $u\ne0$, $E_Cu=0$ would place $F_\perp^{\,0}u$ in
$\operatorname{im}L_a$, contradicting the invertibility of
the full frame in NCRP3. Since $P_Z$ is the exact $G_N$-orthogonal
projector, $L_a^*G_NE_C=0$ and $E_C^*G_NE_C=\Sigma$.
Substitution now gives
\[
 \begin{gathered}
 v_\lambda=L_a\widehat z+E_Cu,\qquad
 \|v_\lambda\|_{G_N}^2
       =\widehat z^*\Gamma\widehat z+u^*\Sigma u,\\
 \operatorname{dist}_{G_N}(v_\lambda,\operatorname{im}L_a)^2
       =u^*\Sigma u>0.
 \end{gathered}                                                   \tag{NCRP14}
\]
This is a quantitative exact repair of the omitted-class issue.
The cross blocks in NCRP10 have contributed to both $\widehat z$
and $\Sigma$; they were not set to zero.

Write the action in the original algebraic frame, now ordered by
$(C,Z)$, as
\[
 A F_\perp^{\,0}=F_\perp^{\,0}A_{CC}+L_aA_{ZC},\qquad
 A L_a=F_\perp^{\,0}A_{CZ}+L_aA_{ZZ}.
 \tag{NCRP15}
\]
These matrices are the corresponding blocks of NCRP9 and therefore
are explicitly calculated from $K_{\mathcal I}^{-1}CK_{\mathcal I}$.
In the metric-orthogonal full frame $[L_a,E_C]$ they give
\[
 \begin{gathered}
 A_{\rm nat}=\Gamma^{-1}L_a^*G_NAL_a
              =A_{ZZ}+\Gamma^{-1}H_{ZC}A_{CZ},\\
 R_N=AL_a-L_aA_{\rm nat}=E_CA_{CZ},\\
 B_N=\Gamma^{-1}L_a^*G_NAE_C,\qquad
 D_N^{C}=\Sigma^{-1}E_C^*G_NAE_C,\\
 [A]_{[L_a,E_C]}
       =\begin{pmatrix}A_{\rm nat}&B_N\\ A_{CZ}&D_N^{C}\end{pmatrix}.
 \end{gathered}                                                   \tag{NCRP16}
\]
The $D_N^{C}$ here is a complement action block, not the native
Gamma metric named $D_N$ in earlier source modules.
For the lower-left block, subtracting $P_ZAL_a$ gives
$E_CA_{CZ}$, and multiplication by the inverse of $\Sigma$
proves its coefficient is exactly $A_{CZ}$.
Thus the full eigenclass has the two coupled equations
\[
 (A_{\rm nat}-\lambda I)\widehat z+B_Nu=0,\qquad
 A_{CZ}\widehat z+(D_N^{C}-\lambda I)u=0 .
 \tag{NCRP17}
\]
These equations require no inverse of $D_N^C-\lambda I$;
they hold on every original rank stratum.
The original FCI39 residual is exactly the matrix $R_N$ in
NCRP16. In particular
\[
 R_N^*G_NR_N=A_{CZ}^*\Sigma A_{CZ},\qquad
 \operatorname{rank}R_N=\operatorname{rank}A_{CZ}
                         =\operatorname{rank}R_{\mathcal T}>0 .
 \tag{NCRP18}
\]
The last equality follows by taking the quotient of the two full
coefficient frames: $F_\perp^{\,0}$ identifies its coefficient space
with $E/\operatorname{im}L_a$, while the omitted residue-dual
coordinates give another isomorphism of that same quotient.
Their transition is invertible, so both represent the same action
map into that quotient and have the same rank.

Define
\[
 \begin{split}
 D_Z&=A_{\rm nat}^*\Gamma+\Gamma A_{\rm nat}-a\Gamma,\\
 D_C&=(D_N^{C})^*\Sigma+\Sigma D_N^{C}-a\Sigma,\\
 X_N&=\Gamma B_N+A_{CZ}^*\Sigma.
 \end{split}                                                     \tag{NCRP19}
\]
The complete centered same-class equation becomes
\[
 \boxed{\quad
 \widehat z^*D_Z\widehat z+
    2\Re(\widehat z^*X_Nu)+u^*D_Cu
 =2a\delta(\widehat z^*\Gamma\widehat z+u^*\Sigma u).
 \quad}                                                         \tag{NCRP20}
\]
It follows by multiplying the two full blocks in NCRP16 against
the block diagonal metric $\operatorname{diag}(\Gamma,\Sigma)$.
Both leakage directions, $\Gamma B_N$ and $A_{CZ}^*\Sigma$,
remain separately in $X_N$. No vanishing or sign is inferred from
the positivity of $\Gamma$ or $\Sigma$.

The original native coefficient metric is still
$D_{a,N-v}^{(s)}$, at centre $a/2-4$ and order $s=1,a$.
Its exact comparison with the arithmetic section is the positive
matrix
\[
 (D_{a,N-v}^{(s)})^{-1/2}\Gamma
                         (D_{a,N-v}^{(s)})^{-1/2}.
 \tag{NCRP21}
\]
Equations NCRP13--20 do not assign this comparison an isometry or
a bounded condition number. The full conductor square in DNE and
XIM remains the actual morphism connecting those two presentations.

\subsection{The consecutive chart and the first physical metric}

The consecutive chart is used only through the proved chart
transition in the preceding section. In that chart
$\operatorname{im}L_a=\mathbb C[S]_{\le m-1}$.
At $N=q-1$ the remainder quotient has no relation to minimize,
so $G_{q-1}=H_{q-1}^{\rm ar}$ and $m<q$ implies that $Sf$
is still the literal polynomial for each $f$ in this section.
For two such polynomials, the physical identity
$S+\bar S=a$ gives
\[
 \langle f,Sg\rangle_{\rm ar}+
 \langle Sf,g\rangle_{\rm ar}
       =a\langle f,g\rangle_{\rm ar}.
 \tag{NCRP22}
\]
Projecting $Sf,Sg$ onto the same section changes neither of these
pairings. Therefore $D_Z=0$ at this cutoff in this chart.
The full receiver is consequently the exact evaluated identity
\[
 2\Re(\widehat z^*X_{q-1}u)+u^*D_Cu
       =2a\delta\,\|v_\lambda\|_{G_{q-1}}^2>0.
 \tag{NCRP23}
\]
It identifies the entire centered action of the retained class with
the complement and both cross terms. A centered compression by itself
does not bound these terms.

For its residual, let $p_n(y)$ be the monic real orthogonal
polynomials of the original arithmetic measure, with actual norms
$\omega_n=\int_{\mathbb R}|p_n(y)|^2dm_a^{\rm ar}(y)>0$.
Set $\phi_n(S)=p_n((S-c)/i)/\sqrt{\omega_n}$.
Multiplication by the real variable $y$ in this orthonormal frame is
Hermitian. Its coefficient from $\phi_{m-1}$ to $\phi_m$ is
$\sqrt{\omega_m/\omega_{m-1}}$, by comparing the leading coefficients.
For $j<m-1$, $y\phi_j$ has degree below $m$ and has no omitted
component. This proves the full formula
\[
 R_{q-1}
  =i\sqrt{\omega_m/\omega_{m-1}}\,
                       \phi_m e_{m-1}^*,\qquad
 \|R_{q-1}\|=\sqrt{\omega_m/\omega_{m-1}} .
 \tag{NCRP24}
\]
This is the original orthogonal-polynomial recurrence coefficient,
whose conventional formulation is DLMF18.2(iv)
\cite{a02nc:human:dlmf18}. The proof above fixes its phase and scale.
For the two Gamma measures, and only for them,
$\omega_{n,s}=(2\pi)^{s/2}n!(s/2)_n$, giving
\[
 \|R_{q-1}^{(s)}\|=\sqrt{m(m-1+s/2)},\quad s=1,a,\qquad
 \frac{\|R_{q-1}^{(1)}\|}{a}\to16,\quad
 \frac{\|R_{q-1}^{(a)}\|}{a}\to\sqrt{264}.
 \tag{NCRP25}
\]
The complete arithmetic ratio in NCRP24 has not been replaced by
either Gamma ratio.

For the original arithmetic metric and nonzero
$f\in\mathbb C[S]_{\le m-1}$,
$(A-\lambda)f=(S-\lambda)f$ is literal. Thus
\[
 \begin{split}
 \|(A-\lambda)f\|_{\rm ar}^2
 &=(a\delta)^2\|f\|_{\rm ar}^2\\
 &\quad+\int_{\mathbb R}(y-a\gamma)^2
                      |f(c+iy)|^2dm_a^{\rm ar}(y)
 >(a\delta)^2\|f\|_{\rm ar}^2 .
 \end{split}                                                     \tag{NCRP26}
\]
The strict inequality follows because the integrand can vanish
almost everywhere only if the nonzero polynomial $f$ does.
This is a statement in the consecutive comparison chart at the
first cutoff; no estimate is transported to another chart by
dropping its conductor transition.

\subsection{The original observation projection on the completed class}

Let $\Lambda:E\to B_{\rm obs}$ be any one of the original packet
observation maps, and let $I_K$ be its original full kernel frame.
Its arithmetic orthogonal projector and complement are
\[
 P_K=I_K(I_K^*G_NI_K)^{-1}I_K^*G_N,\qquad P_B=I-P_K.
 \tag{NCRP27}
\]
For a zero-dimensional kernel the first map is zero. For the full
kernel it is the identity. In all cases
$P_K^2=P_K$, $P_K^*G_N=G_NP_K$, and
$P_Kv+P_Bv=v$.
The literal projected vector is
\[
 P_Kv_\lambda
  =P_KF_0u_0+P_KF_C\eta+P_KL_az.
 \tag{NCRP28}
\]
This equation, with the same original $P_K$, is the receiving
correction for any subsequent kernel response. It never replaces
$\ker\Lambda$ by $\operatorname{im}L_a$.

Put $x_K=P_Kv_\lambda$, $x_B=P_Bv_\lambda$ and
$D=A^*G_N+G_NA-aG_N$. Exact substitution of
$Av_\lambda=\lambda v_\lambda$ gives
\[
 \begin{split}
 x_K^*Dx_K
   &=2a\delta\|x_K\|_{G_N}^2
                    -2\Re(x_K^*G_NAx_B),\\
 x_B^*Dx_B
   &=2a\delta\|x_B\|_{G_N}^2
                    -2\Re(x_B^*G_NAx_K),\\
 2\Re(x_K^*Dx_B)
   &=2\Re\{x_K^*G_NAx_B+x_B^*G_NAx_K\}.
 \end{split}                                                     \tag{NCRP29}
\]
For example, $Ax_K=\lambda(x_K+x_B)-Ax_B$ and
$x_K^*G_Nx_B=0$ prove the first line. The second follows by
the other decomposition; in the last line expand $D$ and use
orthogonality to remove only its actual $aG_N$ term.
Adding all three equations returns
$2a\delta\|v_\lambda\|_{G_N}^2$ exactly.
Thus both projected signs still depend on their actual leakage
pairings, now evaluated on the completed class NCRP28.

\subsection{The literal extension of the relative comparison cone}

For four copies let $\mathcal H=\mathbb C^{4m}$ and let
$L=I_4\otimes L_a$. The original coefficient cone is
$B\to W\oplus S\to\mathcal H$ with differentials
$b\mapsto(b,-b)$ and $(w,s)\mapsto w+s$.
The enlarged arithmetic packet cone is obtained by replacing its
last term with $E^{\oplus4}$ and using $L(w+s)$ there.
Since $L$ is injective, its middle cohomology remains
$(S\cap W)/B$. If $T=S+W$, its final cohomology is
$E^{\oplus4}/L(T)$, and the exact sequence is
\[
 0\longrightarrow L(\mathcal H)/L(T)
 \longrightarrow E^{\oplus4}/L(T)
 \longrightarrow E^{\oplus4}/L(\mathcal H)
 \longrightarrow0.
 \tag{NCRP30}
\]
Each map is induced by the identity on representatives. Its exactness
follows by testing whether a representative lies in $L(\mathcal H)$.
The full complement from NCRP3 provides explicit representatives for
the last quotient, of dimension $4(v+q')$.
All quotient metrics are the attained quotients of the full original
$G_N^{\oplus4}$, so their Schur blocks retain the cross terms in
NCRP10 and NCRP13. A retained eigenclass in one indicated copy lies
outside $L(\mathcal H)$ and therefore has a nonzero class in this
enlarged final cohomology.

Multiplication by $S$ on the full packet is the exact action
NCRP9/NCRP16. The native subspace has the nonzero defect NCRP18,
so the action is carried by the complete ambient block matrices.
No unproved preservation of $B,S,W,T$ is used to create an action
on their quotients. In particular NCRP30 is an explicit extension
of the old relative cohomology, with its additional support and
metric; it does not identify a relative denominator with a theta
boundary. The original native determinant allocations remain valid
on their own maps and metrics, and their arithmetic use proceeds
through the completed class and the complete receiving equations above.


% Additive proof fragment.  All matrices use literal powers of S.
% Provider bibliography key: prog:arithmetic.
\section{The full arithmetic action boundary in the original quotient}

The polynomial quotient and positive arithmetic moments in this section are
the original ones of the arithmetic construction
\cite{a02nc:prog:arithmetic}.  The calculation below uses the full relation ideal,
the original physical line, and the attained quotient metric.  Every identity
needed for the boundary calculation is proved here.

\subsection{Spaces, moments, and the attained section}

In the actual packet \(a=k+4\), with \(k\) in its original guarded domain;
retain that domain and its original \(0<\delta<1/2\) and \(\gamma>2\).
The finite identities below in fact use only that \(a\) is a positive
integer and \(\delta,\gamma>0\).  Put
\begin{equation}
 \beta_{rt}=\frac a2+(2r-a)\delta+i(2t-a)\gamma,\qquad
 \chi(S)=\prod_{r=0}^{a}\prod_{t=0}^{a}(S-\beta_{rt}),\qquad
 q=(a+1)^2 .
 \tag{NAB1}\label{a02nc:eq:NAB1}
\end{equation}
These roots are pairwise distinct: equality of real parts gives equality
of \(r\), and equality of imaginary parts then gives equality of \(t\).
No root or coefficient of \(\chi\) is discarded.  Write
\(\mathcal P_j=\{p\in\mathbb C[S]:\deg p\leq j\}\), with
\(\mathcal P_{-1}=\{0\}\), and
\(\mathcal E=\mathbb C[S]/(\chi)\).  Coordinates in
\(\mathcal P_j\) are the literal coefficients of
\(1,S,\ldots,S^j\); coordinates in \(\mathcal E\) are the corresponding
classes of \(1,S,\ldots,S^{q-1}\).

Fix an integer \(N\geq q-1\).  Use the original positive measure
\(d\mu_N(y)\) on the physical line \(S=a/2+iy\), with all of its original
components and both integration regions.  The subscript allows the measure
to depend on the fixed endpoint \(N\).  Within this calculation the measure
is held fixed.  Its moment matrix through degree \(N+1\) is
\begin{equation}
 H_j^{[N]}=
 \left(
 \int_{\mathbb R}
    \overline{(a/2+iy)^r}(a/2+iy)^t\,d\mu_N(y)
 \right)_{0\leq r,t\leq j},
 \qquad 0\leq j\leq N+1 .
 \tag{NAB2}\label{a02nc:eq:NAB2}
\end{equation}
All moments displayed here are finite and \(H_{N+1}^{[N]}\) is positive
definite in the original arithmetic construction.  Equivalently, the
integral of \(|p(a/2+iy)|^2\) is strictly positive for every nonzero
\(p\in\mathcal P_{N+1}\).  The superscript records that
\(H_N^{[N]}\) and \(H_{N+1}^{[N]}\) use the same measure; an endpoint change
of measure is not implicit in this notation.  In the remainder of this
section write these matrices as \(H_N\) and \(H_{N+1}\).

The Hermitian product is conjugate linear in its first entry:
\(\langle f,g\rangle_N=\int\overline{f(a/2+iy)}g(a/2+iy)d\mu_N(y)\).
For \(f,g\in\mathcal P_N\), the identity
\(\overline S+S=a\) on this actual line gives
\begin{equation}
 \langle Sf,g\rangle_N+\langle f,Sg\rangle_N
       =a\langle f,g\rangle_N .
 \tag{NAB3}\label{a02nc:eq:NAB3}
\end{equation}
The two left-hand products use the moment matrix through degree \(N+1\).
Thus (NAB3) retains the physical coordinate rather than replacing the
action by an abstract self-adjoint multiplication operator.

Let \(\rho_N:\mathcal P_N\longrightarrow\mathcal E\) be reduction modulo
\(\chi\).  Polynomial division by the original monic \(\chi\) proves
\begin{equation}
 \ker\rho_N=\mathcal R_N:=\chi\mathcal P_{N-q},\qquad
 \mathcal P_N/\mathcal R_N \xrightarrow[\rho_N]{\ \cong\ }\mathcal E .
 \tag{NAB4}\label{a02nc:eq:NAB4}
\end{equation}
For \(N=q-1\), \(\mathcal R_N=\{0\}\).
Let \(E_N:\mathbb C^q\to\mathbb C^{N+1}\) be the coefficient inclusion of
\(\mathcal P_{q-1}\) into \(\mathcal P_N\).  If \(N\geq q\), let \(W_N\)
have columns the coefficient vectors of
\(\chi,S\chi,\ldots,S^{N-q}\chi\).  Its columns are linearly independent
because multiplication by a nonzero polynomial is injective.
Consequently \(W_N^*H_NW_N\) is positive definite.  The attained section is
\begin{equation}
 T_N=E_N-W_N(W_N^*H_NW_N)^{-1}W_N^*H_NE_N,\qquad
 G_N=T_N^*H_NT_N ;
 \quad T_{q-1}=E_{q-1}.
 \tag{NAB5}\label{a02nc:eq:NAB5}
\end{equation}
Here \(T_N:\mathcal E\to\mathcal P_N\) uses the fixed coefficient
coordinates, and \(G_N\) is its \(q\)-by-\(q\) attained metric.
Indeed, \(\rho_NT_N=I_{\mathcal E}\) by (NAB4), and direct multiplication
gives \(W_N^*H_NT_N=0\).  Every representative of \(v\in\mathcal E\)
is \(T_Nv+W_Nz\) for a unique \(z\).  Orthogonality therefore proves
\[
 \|T_Nv+W_Nz\|_N^2
 =v^*G_Nv+z^*W_N^*H_NW_Nz.
\]
This proves attainment, uniqueness, and the original minimum property
without a change of relation space.  It also proves \(G_N>0\), since
\(T_Nv\ne0\) whenever \(v\ne0\).

\subsection{The exact one-column boundary and its Hermitian form}

Let \(A:\mathcal E\to\mathcal E\) be multiplication by the original \(S\)
modulo \(\chi\).  Write
\(\chi(S)=S^q+\sum_{j=0}^{q-1}c_jS^j\).
In the fixed ordered basis,
\begin{equation}
 A e_j=e_{j+1}\quad(0\leq j<q-1),\qquad
 A e_{q-1}=-\sum_{j=0}^{q-1}c_je_j .
 \tag{NAB6}\label{a02nc:eq:NAB6}
\end{equation}
Embed \(T_N\) into \(\mathcal P_{N+1}\) by appending a zero top
coefficient, writing this embedding as \(\widetilde T_N\).
Define the row, polynomial column, and vector
\begin{equation}
 \ell_N=[S^N]T_N:\mathcal E\longrightarrow\mathbb C,\qquad
 b_N=\operatorname{coeff}_{0,\ldots,N+1}
             \bigl(\chi(S)S^{N+1-q}\bigr),\qquad
 \xi_N=\widetilde T_N^*H_{N+1}b_N\in\mathbb C^q .
 \tag{NAB7}\label{a02nc:eq:NAB7}
\end{equation}
The notation \([S^N]\) means the literal coefficient functional, with no
factorial, root evaluation, or rescaling.

Let \(M_N:\mathcal P_N\to\mathcal P_{N+1}\) denote multiplication by \(S\)
in those same coordinates.  There is a linear map
\(Z_N:\mathcal E\to\mathcal R_N\), embedded in \(\mathcal P_{N+1}\),
such that
\begin{equation}
 M_NT_N-\widetilde T_N A=b_N\ell_N+Z_N,\qquad
 \widetilde T_N^*H_{N+1}Z_N=0.
 \tag{NAB8}\label{a02nc:eq:NAB8}
\end{equation}
To prove the first identity on \(v\in\mathcal E\), reduce
\(S(T_Nv)-T_N(Av)\) modulo \(\chi\).
The result is \(Av-Av=0\), so the polynomial is divisible by \(\chi\).
Its degree is at most \(N+1\), and its coefficient of \(S^{N+1}\) is
exactly \(\ell_Nv\).  Monicity of the full \(\chi\) now shows that
subtracting \(\chi S^{N+1-q}\ell_Nv\) leaves an element of
\(\chi\mathcal P_{N-q}\).
This subtraction is linear in \(v\), giving \(Z_N\).
If \(N=q-1\), the remainder has degree less than \(q\) and is divisible
by \(\chi\), hence is zero.
Otherwise the second identity is exactly the orthogonality in (NAB5):
both factors of this product have degree at most \(N\), so the use of
\(H_{N+1}\) restricts to the same \(H_N\).

Define the full arithmetic action form
\begin{equation}
 D_N=A^*G_N+G_NA-aG_N .
 \tag{NAB9}\label{a02nc:eq:NAB9}
\end{equation}
It has the exact boundary formula
\begin{equation}
 \boxed{\displaystyle
 D_N=-\xi_N\ell_N-\ell_N^*\xi_N^*.}
 \tag{NAB10}\label{a02nc:eq:NAB10}
\end{equation}
For a direct proof, (NAB8) and its orthogonality give
\[
 \widetilde T_N^*H_{N+1}M_NT_N
       =G_NA+\xi_N\ell_N .
\]
Taking conjugate transposes gives the other term.
Applying (NAB3) to every pair of columns of \(T_N\) gives
\[
 aG_N=
 \widetilde T_N^*H_{N+1}M_NT_N+
 (M_NT_N)^*H_{N+1}\widetilde T_N
 =G_NA+A^*G_N+\xi_N\ell_N+\ell_N^*\xi_N^*.
\]
Rearranging proves (NAB10), including both minus signs.
In particular \(\operatorname{rank}D_N\leq2\).
No estimate for a trial section or different denominator was used.

\subsection{Exact inertia, trace, and boundary magnitude}

For every root \(\beta\) in (NAB1), define the original Lagrange
polynomial and its class
\begin{equation}
 L_\beta(S)=\frac{\chi(S)}{(S-\beta)\chi'(\beta)}
       \in\mathcal P_{q-1},\qquad
 v_\beta=[L_\beta]\in\mathcal E .
 \tag{NAB11}\label{a02nc:eq:NAB11}
\end{equation}
The denominator \(\chi'(\beta)\) is nonzero by the distinctness proved
after (NAB1).  At the original roots \(L_\beta(\beta)=1\) and
\(L_\beta(\beta')=0\) for \(\beta'\ne\beta\).
Thus \(v_\beta\ne0\), and the \(q\) such classes form a basis: a linear
relation among them evaluates at each root to its corresponding
coefficient.  The literal polynomial identity
\((S-\beta)L_\beta=\chi/\chi'(\beta)\) proves the exact morphism
\begin{equation}
 Av_\beta=\beta v_\beta,\qquad
 v_{\beta_{rt}}^*D_Nv_{\beta_{rt}}
      =2(2r-a)\delta\,v_{\beta_{rt}}^*G_Nv_{\beta_{rt}} .
 \tag{NAB12}\label{a02nc:eq:NAB12}
\end{equation}
In particular the expression is strictly negative at \(r=0\) and
strictly positive at \(r=a\).
A Hermitian form of rank at most two with both a positive and a negative
value must have rank exactly two and one sign of each kind:
\begin{equation}
 \operatorname{inertia}(D_N)=(1,1,q-2).
 \tag{NAB13}\label{a02nc:eq:NAB13}
\end{equation}
For completeness, this elementary inference can be seen by diagonalizing
the Hermitian form using successive orthogonal eigenvectors.  Such a
basis is obtained by maximizing its real Rayleigh quotient on the compact
unit sphere; differentiation in every real and imaginary tangent
direction makes a maximizing vector an eigenvector.  Its orthogonal
complement is invariant because the matrix is Hermitian, so induction
gives the diagonalization.  A positive value requires a positive
diagonal entry and a negative value requires a negative diagonal entry.
The rank bound allows precisely those two nonzero entries.

The full original grid also gives an exact trace cancellation:
\begin{equation}
 \operatorname{tr}(G_N^{-1}D_N)
 =\overline{\operatorname{tr}A}+\operatorname{tr}A-aq
 =2\operatorname{Re}\sum_{r,t=0}^{a}\beta_{rt}-aq=0 .
 \tag{NAB14}\label{a02nc:eq:NAB14}
\end{equation}
The first equality uses only cyclicity of finite matrix trace.
For the last equality,
\(\sum_{r=0}^a(2r-a)=0=\sum_{t=0}^a(2t-a)\), so the root sum
is exactly \(aq/2\).
The Lagrange basis in (NAB11) proves that this root sum is the trace of
the actual \(A\); no approximate eigenvalue calculation is involved.

Take the positive Hermitian square root of \(G_N\), and define
\begin{equation}
 B_N=G_N^{-1/2}D_NG_N^{-1/2},\qquad
 u_N=G_N^{-1/2}\xi_N,\qquad
 w_N=G_N^{-1/2}\ell_N^*,\qquad
 c_N=w_N^*u_N=\ell_NG_N^{-1}\xi_N .
 \tag{NAB15}\label{a02nc:eq:NAB15}
\end{equation}
The square root is obtained by the same finite Hermitian
diagonalization, taking positive square roots of the positive entries of
\(G_N\).  Equations (NAB10) and (NAB14) prove
\begin{equation}
 B_N=-u_Nw_N^*-w_Nu_N^*,\qquad
 -2\operatorname{Re}c_N=\operatorname{tr}B_N=0.
 \tag{NAB16}\label{a02nc:eq:NAB16}
\end{equation}
Congruence by the invertible \(G_N^{-1/2}\) preserves positive and
negative subspaces, hence (NAB13) applies to \(B_N\) too.
In particular \(u_N,w_N\) are linearly independent.
Set \(b=\|w_N\|>0\), \(e_1=w_N/b\), and write
\[
 u_N=(c_N/b)e_1+\eta e_2,\qquad
 \eta=\left(\|u_N\|^2-|c_N|^2/b^2\right)^{1/2}>0,
\]
where \(e_2\) is the resulting unit vector orthogonal to \(e_1\).
On their span the matrix of \(B_N\) is exactly
\[
 \begin{pmatrix}-2\operatorname{Re}c_N&-b\eta\\
                 -b\eta&0\end{pmatrix},
\]
and on its orthogonal complement it is zero.  Since
\(\operatorname{Re}c_N=0\), its nonzero eigenvalues are
\begin{equation}
 \boxed{\displaystyle
 \lambda_\pm(B_N)=\pm\kappa_N,\qquad
 \kappa_N=
 \sqrt{\|u_N\|^2\|w_N\|^2-
                    (\operatorname{Im}(w_N^*u_N))^2}>0 .}
 \tag{NAB17}\label{a02nc:eq:NAB17}
\end{equation}
The equivalent formula using only the original coefficient matrices is
\begin{equation}
 \kappa_N^2=
  (\xi_N^*G_N^{-1}\xi_N)
  (\ell_NG_N^{-1}\ell_N^*)
  -|\ell_NG_N^{-1}\xi_N|^2
  =\tfrac12\operatorname{tr}\bigl((G_N^{-1}D_N)^2\bigr).
 \tag{NAB18}\label{a02nc:eq:NAB18}
\end{equation}
The first equality follows from (NAB15)--(NAB17), and the second follows
because \(G_N^{-1}D_N\) is similar to \(B_N\).
The exact grid eigenclasses in (NAB12) furthermore give the finite bound
\begin{equation}
 \kappa_N\geq2a\delta .
 \tag{NAB19}\label{a02nc:eq:NAB19}
\end{equation}
Indeed, put \(z=G_N^{1/2}v_{\beta_{0t}}\).  It is nonzero and
\(z^*B_Nz/(z^*z)=-2a\delta\).
Every Rayleigh quotient of the diagonal matrix with entries
\(-\kappa_N,\kappa_N,0,\ldots,0\) lies in
\([-\kappa_N,\kappa_N]\); this proves (NAB19).
The choice \(r=a\) gives the matching positive quotient \(2a\delta\).

If the complete grid trace cancellation is not available, the boundary
identity (NAB10) still has its stated sign whenever (NAB2)--(NAB8) hold,
but the two eigenvalues are instead
\[
 -\operatorname{Re}c_N
 \ \pm\sqrt{\|u_N\|^2\|w_N\|^2-(\operatorname{Im}c_N)^2}.
\]
This follows by taking the characteristic polynomial of the displayed
two-by-two matrix.  In the actual full grid the shift vanishes exactly
by (NAB14).  Thus omission of a root subfamily cannot silently retain
the full-grid trace cancellation.

\subsection{An exact original-coordinate sign diagnostic}

The following small example is an exact algebraic diagnostic only.  It
is not an actual packet or period, and is not an evaluation of the
arithmetic source.  It neither replaces nor relaxes any original
large-degree finite guard.
Take \(a=1\), \(\delta=1/4\), \(\gamma=3\), \(q=4\), \(N=4\),
\(S=1/2+iy\), and the standard real Gaussian measure.  The original
root polynomial and relation space are
\begin{equation}
 \begin{aligned}
 \chi(S)&=
 \bigl((S-\tfrac12-\tfrac14)^2+9\bigr)
 \bigl((S-\tfrac12+\tfrac14)^2+9\bigr)\\
 &=S^4-2S^3+\tfrac{155}{8}S^2-\tfrac{147}{8}S+\tfrac{22185}{256},
 \qquad \mathcal R_4=\operatorname{span}\{\chi\}.
 \end{aligned}
 \tag{NAB20}\label{a02nc:eq:NAB20}
\end{equation}
The roots are exactly \(1/2\pm1/4\pm3i\), all in the right half-plane.
For \(m\geq0\), put
\(\mu_{2m}=(2m-1)!!\), with \((-1)!!=1\), and
\(\mu_{2m+1}=0\).  These are the Gaussian moments: the odd integral
vanishes under \(y\mapsto-y\), and integration by parts gives
\(\mu_{2m}=(2m-1)\mu_{2m-2}\).
The moment entries, entirely in the original \(S\)-coordinates, are
\begin{equation}
 (H_5)_{mn}
   =\sum_{r=0}^{m}\sum_{t=0}^{n}
    \binom mr\binom nt\,2^{-(m+n-r-t)}
            (-\mathrm i)^r\mathrm i^t\mu_{r+t}.
 \tag{NAB21}\label{a02nc:eq:NAB21}
\end{equation}
Here \(\mathrm i^2=-1\), and \(0\leq m,n\leq5\).
Every nonzero polynomial has positive squared integral because a
nonzero polynomial has finitely many zeros and the Gaussian density is
strictly positive.  Thus the moment and quotient matrices are positive
definite.

Write
\(c=(22185/256,-147/8,155/8,-2,1)^{\mathsf T}\), and put
\(d=316481153\).
The finite sums (NAB21) give \(c^*H_4c=d/65536\).
Substitution into (NAB5) and (NAB7) gives, exactly,
\begin{equation}
 \begin{aligned}
 \ell_4&=\frac1d(-4407552,-2203776,1749184,3725664),\\
 \xi_4&=\frac1{1024}(0,44548,44548,-51849)^{\mathsf T}.
 \end{aligned}
 \tag{NAB22}\label{a02nc:eq:NAB22}
\end{equation}
and hence
\begin{equation}
 D_4=-\xi_4\ell_4-\ell_4^{\mathsf T}\xi_4^{\mathsf T}.
 \tag{NAB23}\label{a02nc:eq:NAB23}
\end{equation}
To specify the remaining finite computation without omitting its
inputs, take \(E_4\) to be the first four columns of the five-by-five
identity, \(T_4=E_4-(65536/d)c(c^*H_4E_4)\), and
\(G_4=T_4^*H_4T_4\), with \(H_4\) the leading block of (NAB21).
Direct substitution and ordinary exact matrix multiplication give
\begin{equation}
 \begin{aligned}
 \operatorname{tr}(G_4^{-1}D_4)&=0,\\
 \frac{\det(zG_4-D_4)}{\det G_4}
      &=z^2\left(z^2-\frac{46738154483830145}{497781812232192}\right),\\
 \kappa_4^2&=\frac{46738154483830145}{497781812232192}.
 \end{aligned}
 \tag{NAB24}\label{a02nc:eq:NAB24}
\end{equation}
All the matrices and exact rational substitutions for this diagnostic
are recorded by the accompanying
\texttt{exact\_boundary\_example.py} and
\texttt{EXACT\_EXAMPLE.json}.  The general proof is (NAB3)--(NAB19);
the example supplies an independent check of its signs and its use of
the nontrivial attained relation space.

\paragraph{Scope of the calculation.}
Equations (NAB10), (NAB13), and (NAB17)--(NAB19) determine the full
action form and its full-space inertia.  For a specified map
\(J:\mathbb C^d\to\mathcal E\), their exact receiver is
\[
 J^*D_NJ=-(J^*\xi_N)(\ell_NJ)
                 -(\ell_NJ)^*(\xi_N^*J).
\]
This formula carries the original class-position data into the
restricted form.  Evaluating the sign on an individual projected class
requires its actual \(J\), or the corresponding attained projection,
inside that formula.  No individual projected sign, absolute
allocation, or value of the responses \(U,V\) is asserted here.


% Exact FCI/DNE-to-EC dictionary; new actual class completion.
\section{The original EC three-column response after full-class completion}

We use exactly the arithmetic packet, physical source and attained
section of NCRP8. The letter $a$ is the tensor degree used in this
manuscript; it replaces EC's tensor-degree letter $k$, and does not
change the original source order or packet. Let $p_n(S)$ be the
original arithmetic monic orthogonal polynomials at $S=c+iy$ and
let $\omega_n=\|p_n\|_{\rm ar}^2$. Explicitly, if
$p_n^{y}$ is the monic real-variable polynomial in NCRP24, then
$p_n(S)=i^n p_n^{y}((S-c)/i)$; this is monic in the original $S$
and has exactly the same norm $\omega_n$. Put $b_n=[p_n]\in E$.
Here $b_n$ retains EC's polynomial-class notation; the relation
column called $b_N$ in NAB7 is called $r_N$ below.
At each original cutoff $N$ the three columns and full Grams are
precisely
\[
 \begin{gathered}
 W=[b_N,b_{N+1},v_\lambda],\qquad F=W^*G_NW,\\
 K=W^*G_N I_K(I_K^*G_NI_K)^{-1}I_K^*G_NW,\qquad B=F-K,\\
 F_{13}=b_N^*G_Nv_\lambda,\quad
 F_{23}=b_{N+1}^*G_Nv_\lambda,\quad F_{33}=\|v_\lambda\|_{G_N}^2 .
 \end{gathered}                                                   \tag{NCE1}
\]
These are EC8 and EC37 \cite{a02nc:prog:EC}; the kernel projector is
the original $P_K$ of NCRP27, with its complete inverse.

\subsection{Identify both boundary rows, with their orientation}

Let $\ell_N,\xi_N$ be the full action-boundary rows proved in the
preceding section, and put $r_N(S)=\chi(S)S^{N+1-q}$.
For every class $x$, its attained lift $T_Nx$ has leading coefficient
$\ell_Nx$. Because $p_N$ is monic and orthogonal to every lower
degree polynomial,
\[
 \langle p_N,T_Nx\rangle_{\rm ar}=\omega_N\ell_Nx.
 \tag{NCE2}
\]
The polynomial $p_N-T_Nb_N$ is an old relation, orthogonal to
$T_Nx$. Thus the left side is $b_N^*G_Nx$, and consequently
\[
 \boxed{\ell_N=\omega_N^{-1}b_N^*G_N.}
 \tag{NCE3}
\]
Since $p_{N+1}$ and $r_N$ are both monic of degree $N+1$,
$p_{N+1}-r_N$ has degree at most $N$ and represents $b_{N+1}$.
The orthogonality of $p_{N+1}$ to every polynomial of degree
at most $N$ gives
\[
 \langle T_Nx,r_N\rangle_{\rm ar}
   =-\langle T_Nx,p_{N+1}-r_N\rangle_{\rm ar}
   =-x^*G_Nb_{N+1}.
 \tag{NCE4}
\]
The definition of $\xi_N$ therefore proves the second exact row:
\[
 \boxed{\xi_N=-G_Nb_{N+1}.}
 \tag{NCE5}
\]
In particular the original centered action matrix is
\[
 \boxed{\quad
 D_N=A^*G_N+G_NA-aG_N
 =\frac{G_Nb_{N+1}b_N^*G_N+G_Nb_Nb_{N+1}^*G_N}{\omega_N}.
 \quad}                                                         \tag{NCE6}
\]
This is a direct coefficient proof of the rank-two marked-pairing
matrix on the original arithmetic packet. Its sign agrees with EC10.

To identify the complex coefficient, put $v=v_\lambda$,
$f=T_Nv$, $E=F_{33}$ and
\[
 \mu_{N,\lambda}
   =\frac{\int_{\mathbb R}y|f(c+iy)|^2dm_a^{\rm ar}(y)}{E}.
 \tag{NCE7}
\]
The polynomial $(S-\lambda)f$ is a relation through degree $N+1$.
Its top coefficient is $\ell_Nv$, and subtracting
$(\ell_Nv)r_N$ leaves an old relation. Pairing with the attained
section eliminates exactly that old relation and gives
\[
 (v^*\xi_N)(\ell_Nv)
   =\langle f,(S-\lambda)f\rangle_{\rm ar}
   =E\{-a\delta+i(\mu_{N,\lambda}-a\gamma)\}.
 \tag{NCE8}
\]
Combining NCE3 and NCE5, the left side is
$-\overline{F_{23}}F_{13}/\omega_N$. Taking its conjugate proves
the literal EC current, including its phase:
\[
 \boxed{\quad
 \mathfrak c_{N,\lambda}
 :=\frac{\overline{F_{13}}F_{23}}{\omega_NE}
 =a\delta+i(\mu_{N,\lambda}-a\gamma).
 \quad}                                                         \tag{NCE9}
\]
Its positive real part proves $F_{13}\ne0$ and $F_{23}\ne0$.
No asymptotic estimate or period minor is needed to justify these
two denominators.

\subsection{The exact finite current bound in the original source}

Here is a direct finite proof of the useful strict current interval.
The arithmetic source is even in $y$: conjugation preserves the
complete quartet and $2\xi/h$, so $w_h(-y)=w_h(y)$, and its
convolution powers remain even. The full physical root polynomial
$i^{-q}\chi(c+iy)$ is real and has parity $(-1)^q$, because its
roots in the $y$-coordinate are invariant under conjugation and negation.
In particular $\chi(c-iy)=(-1)^q\chi(c+iy)$; no assertion that
$q$ is even is used.
Hence reflection preserves the relation space and its orthogonal
complement $\mathcal V_N=\operatorname{im}T_N$.

Let $\Pi_N$ be the source-orthogonal projector onto $\mathcal V_N$,
let $J=\Pi_NM_y|_{\mathcal V_N}$, and retain
$\widehat b_n=T_Nb_n$. The operator $J$ is Hermitian and
reflection anticommutes with it. The two vectors
$\widehat b_N,\widehat b_{N+1}$ have opposite reflection parities,
so their actual inner product is zero.
Multiplying an attained representative by $S$, subtracting its
leading multiple of $p_{N+1}$, and projecting gives the full
action identity
\[
 T_NAT_N^{-1}
   =cI+iJ+
       \frac{\widehat b_{N+1}\widehat b_N^*}{\omega_N}.
 \tag{NCE10}
\]
The proof of the last coefficient is NCE2; the class of the
subtracted polynomial is $b_{N+1}$, which is added back after
projection. Thus the term has the displayed positive sign.
The old relations are orthogonal to $\mathcal V_N$ and no
additional term is removed.

Put $d=a\delta>0$, $t=a\gamma>0$ and
$R=(d+it-iJ)^{-1}$. For every vector $x$,
$\|(d+it-iJ)x\|\,\|x\|\ge
\Re\langle x,(d+it-iJ)x\rangle=d\|x\|^2$.
The finite square matrix is injective, hence invertible, with
$\|R\|\le d^{-1}$.
The eigenclass equation and NCE10 imply
\[
 f=\frac{F_{13}}{\omega_N}R\widehat b_{N+1},\qquad
 \mathfrak c_{N,\lambda}
  =\frac{\langle\widehat b_{N+1},R\widehat b_{N+1}\rangle}
         {\|R\widehat b_{N+1}\|^2}.
 \tag{NCE11}
\]
The second equation follows by substituting the first into
$F_{23}$ and $E=\|f\|^2$. The original vector and source mass
have not been rescaled.

The spectral measure of the Hermitian $J$ at
$\widehat b_{N+1}$ is even, since reflection anticommutes with
$J$ and preserves that vector up to sign. Under the resolvent
squared norm, the paired spectral points $(x,-x)$ have mean
\[
 \frac{x\{[d^2+(t-x)^2]^{-1}-[d^2+(t+x)^2]^{-1}\}}
 {[d^2+(t-x)^2]^{-1}+[d^2+(t+x)^2]^{-1}}
       =\frac{2tx^2}{d^2+t^2+x^2}\in[0,2t).
 \tag{NCE12}
\]
There is some nonzero spectral mass away from zero. Otherwise
$R\widehat b_{N+1}$ would be a scalar multiple of
$\widehat b_{N+1}$, orthogonal to $\widehat b_N$; taking the
$\widehat b_N$ inner product in NCE11 would contradict
$F_{13}\ne0$.
Thus its total weighted mean is strictly between zero and $2t$.
This total mean is $\mu_{N,\lambda}$ by NCE11 and the definition
of $J$. Consequently
\[
 0<\mu_{N,\lambda}<2a\gamma,\qquad
 \Re\mathfrak c_{N,\lambda}=a\delta,\qquad
 |\Im\mathfrak c_{N,\lambda}|<a\gamma,\qquad
 |\mathfrak c_{N,\lambda}|<a\sqrt{\delta^2+\gamma^2}.
 \tag{NCE13}
\]
The finite spectral argument is the one behind EC11--13. It is
reproduced here to keep the physical moment and orientation explicit.
It supplies no estimate for the actual observation responses.

\subsection{The same original ratios and all three projected terms}

By NCE3--5 the exact dictionary, with no row exchanged, is
\[
 \boxed{\quad
 U=\frac{K_{13}}{F_{13}}
     =\frac{\ell_NP_Kv}{\ell_Nv},\qquad
 V=\frac{K_{23}}{F_{23}}
     =\frac{\xi_N^*P_Kv}{\xi_N^*v}.
 \quad}                                                         \tag{NCE14}
\]
Both occurrences of the minus sign in NCE5 cancel in the second
ratio. These are precisely EC37's $U,V$.
Set
$P_K^{\rm resp}=\bar UV$,
$P_B^{\rm resp}=(1-\bar U)(1-V)$ and
$P_\times^{\rm resp}=\bar U+V-2\bar UV$.
For $x_K=P_Kv$ and $x_B=(I-P_K)v$, define the original three
pairings by
$\Phi_K=x_K^*D_Nx_K$,
$\Phi_B=x_B^*D_Nx_B$ and
$\Phi_\times=2\Re(x_K^*D_Nx_B)$.
Substitution in NCE6 gives
\[
 \boxed{\frac{\Phi_X}{2E}
       =\Re(\mathfrak c_{N,\lambda}P_X^{\rm resp}),
           \qquad X=K,B,\times.}
 \tag{NCE15}
\]
For example the kernel term is
$2\Re(\overline{K_{13}}K_{23})/\omega_N$.
The mixed term is
$2\Re(\overline{K_{13}}B_{23}+\overline{B_{13}}K_{23})/\omega_N$;
substitution of the two ratios gives exactly the displayed
$P_\times^{\rm resp}$. Their sum is one, so the sum of the three
pairings is exactly $2a\delta E$.

Let $w_\lambda=(u_0,\eta,z)$ be NCRP5 and let the two exact rows be
$r_1=\ell_N$, $r_2=\xi_N^*$.
For $j=1,2$, the numerator of the corresponding response is
\[
 r_jP_Kv
    =r_jP_KF_0u_0+r_jP_KF_C\eta+r_jP_KL_az,
 \qquad
 r_jv=r_jF_0u_0+r_jF_C\eta+r_jL_az .
 \tag{NCE16}
\]
Thus NCE14 evaluates the full coefficient map needed to pass from
the native construction to the actual EC responses. Dropping the
first two terms in either expression changes the class and the
response. The individual products and their imaginary parts are
still those of these complete ratios.

Finally let $\theta_K=\|P_Kv\|_{G_N}^2/E$. The two separate action
leakages obtained in NCRP29 are
\[
 \begin{split}
 \frac{2\Re(x_K^*G_NAx_B)}{2E}
       &=a\delta\theta_K-\Re(\mathfrak c P_K^{\rm resp}),\\
 \frac{2\Re(x_B^*G_NAx_K)}{2E}
       &=a\delta(1-\theta_K)-\Re(\mathfrak c P_B^{\rm resp}).
 \end{split}                                                     \tag{NCE17}
\]
These are exactly EC41 with its two block directions retained.
Their sum is $\Phi_\times/(2E)$.
The more precise parity-dependent EC current, and its original
eventual guards and analytic error radius, retain their existing
verification status. This source proves the actual row dictionary,
the complete-class input, the finite current interval, and the
exact three-term and two-leakage receiving maps; it does not
reassign their period-dependent response products or re-prove
the later EC asymptotic radius.

\subsection{The exact parity reduction of the three-column Gram}

The same original reflection used in NCE10 proves a further evaluated
entry of EC's actual Gram. Since the two attained columns have opposite
parities, reflection is a unitary map and changes the sign of their
inner product. Thus
\[
 F_{12}=0,\qquad
 \ell_NG_N^{-1}\xi_N=-F_{12}/\omega_N=0,\qquad
 \kappa_N=\frac{\sqrt{F_{11}F_{22}}}{\omega_N}.
 \tag{NCE18}
\]
The last equality follows by substituting NCE3 and NCE5 in the complete
boundary magnitude NAB18. Both $F_{11},F_{22}$ are strictly positive,
since NCE9 has already proved $F_{13},F_{23}\ne0$.

Put $f=(F_{13},F_{23})^{\mathsf T}$, $E=F_{33}$ and retain the exact
diagonal $D_f=\operatorname{diag}(F_{13},F_{23})$. The literal Schur
complement of the third column and its response-coordinate form are
\[
 \begin{split}
 S_F&=\operatorname{diag}(F_{11},F_{22})-ff^*/E,\\
 T_F&=D_f^{-1}S_FD_f^{-*}
  =\begin{pmatrix}
   F_{11}/|F_{13}|^2-1/E&-1/E\\
   -1/E&F_{22}/|F_{23}|^2-1/E
   \end{pmatrix}\succeq0.
 \end{split}                                                     \tag{NCE19}
\]
Indeed $S_F$ is the Gram of the two columns after subtracting their
exact orthogonal components along $v$, proving positivity even if
those two residual columns are dependent. The displayed congruence
proves every entry, including its negative real off-diagonal value.

Define the two scalar observation amplifications without changing any
source or vector by
\[
 \alpha=\frac{E F_{11}}{|F_{13}|^2},\qquad
 \beta=\frac{E F_{22}}{|F_{23}|^2}.
 \tag{NCE20}
\]
The squared norm of the orthogonal projection of $v$ onto the two
mutually orthogonal original columns is
$|F_{13}|^2/F_{11}+|F_{23}|^2/F_{22}$. It cannot exceed $E$:
subtracting that projection leaves an orthogonal square. Consequently
\[
 \alpha,\beta\ge1,\qquad
 \frac1\alpha+\frac1\beta\le1,\qquad
 \alpha\beta\ge\alpha+\beta,
 \qquad
 \boxed{\alpha\beta=
       \frac{\kappa_N^2}{|\mathfrak c_{N,\lambda}|^2}.}
 \tag{NCE21}
\]
The final equality follows by multiplying NCE20 and using NCE9 and
NCE18. This evaluates the product of the two original denominator
amplifications in the complete action magnitude and the actual current.
It does not assign either factor from the other or from its rank.





% END PRESERVED INPUT 07_NATIVE_CLASS_INTEGRATION.tex
\endgroup

\clearpage
\begingroup


\setcounter{equation}{0}
\renewcommand{\theequation}{A02JRE.\arabic{equation}}
\section*{Joint response ellipse and original projection-energy bounds}
\addcontentsline{toc}{section}{Joint response ellipse and original projection-energy bounds}
% BEGIN PRESERVED INPUT 08_JOINT_RESPONSE_ELLIPSE.tex SHA256 331c0411e0a35d9e4f005a45af318255917f0da1f34efa7ba57d8fe11657f30b
% Root derivation: original EC three-column response, 19 September 2026.
\section{The joint response ellipse in the original three-column Gram}

This calculation starts from the same full retained eigenclass and the same
observation kernel as EC37--41. Its source is the supplied complete EC
continuation, preserved in \texttt{02\_Untitled 1775.md}, lines9401--10815;
the full-class completion is NCRP1--30. Every Gram below uses the original
arithmetic quotient metric $G_N$ and its full attained polynomial lift.
No native-section representative is substituted for the retained class.
The positive-square argument is the Hermitian Schur argument of
Boyd, El Ghaoui, Feron and Balakrishnan, \emph{Linear Matrix Inequalities
in System and Control Theory} (SIAM,1994), section2.1, equations(2.3)--(2.4),
doi:10.1137/1.9781611970777; its rank-deficient version is proved below
on the actual image, so no nonsingularity of a three-column Gram is assumed.

\subsection{Exact map, image, minimum and remainder}
Let $v=v_\lambda$ be the full retained class and let
$b_1=b_N$, $b_2=b_{N+1}$ be EC's original columns. Retain its original
kernel frame $I$, and write
\[
 P=I(I^*G_NI)^{-1}I^*G_N,\quad
 W=[b_1,b_2,v],\quad F=W^*G_NW,\quad K=W^*G_NPW .                 \tag{JRE1}
\]
The zero-dimensional kernel gives $P=0$ and the full kernel gives $P=1$.
In every case $P^2=P$ and $P^*G_N=G_NP$. In particular
$K=W^*P^*G_NPW$, so $0\preceq K\preceq F$.
Set
\[
 E=F_{33}>0,\quad f=\binom{F_{13}}{F_{23}},\quad
 \theta=K_{33}/E\in[0,1],\quad
 x=Pv,\quad z=x-\theta v .                                    \tag{JRE2}
\]
The two entries of $f$ are nonzero on EC's stated domain, by its original
Cauchy-transform proof EC20 and EC27--29. We keep their exact values.
Since $v^*G_Nx=x^*G_Nx=\theta E$, direct expansion gives
\[
 v^*G_Nz=0,\qquad z^*G_Nz=E\theta(1-\theta).                   \tag{JRE3}
\]
Let $B=[b_1,b_2]$ and define the exact residual-column map
\[
 A=B-vf^*/E:\mathbb C^2\longrightarrow E_{\rm packet},\quad
 S=A^*G_NA=F_{[1,2],[1,2]}-ff^*/E\succeq0 .                    \tag{JRE4}
\]
Here $E_{\rm packet}$ denotes the original polynomial quotient, while
$E$ in JRE2 is its retained-class squared norm. Thus their types differ
explicitly. The response residual is
\[
 r=A^*G_Nz=\binom{K_{13}}{K_{23}}-\theta f .                    \tag{JRE5}
\]
For every $c\in\ker S$ positivity implies $Ac=0$; hence $c^*r=0$.
Since $S$ is Hermitian, $r\in\operatorname{im}S$.
Define $S^\dagger$ as the inverse of $S$ on its positive eigenspaces,
and zero on its kernel. This specifies it at every rank without a
regularization parameter. The vector
$z_0=AS^\dagger r$ has $A^*G_Nz_0=r$, and $z-z_0$ is orthogonal to
$\operatorname{im}A$. Expanding the orthogonal square proves
\[
 \boxed{r^*S^\dagger r+
       \|z-AS^\dagger r\|_{G_N}^{2}=E\theta(1-\theta).}       \tag{JRE6}
\]
Thus $AS^\dagger r$ is the attained minimum-norm vector with the actual
two observations $A^*G_Nz=r$. Every other solution is its sum with a
vector in $\ker(A^*G_N)$, whose squared norm adds. This proves the
complete map, its range, its minimum, and its exact residual term.

To state the same identity in the original dimensionless responses,
retain EC37's $U=K_{13}/F_{13}$ and $V=K_{23}/F_{23}$ and put
\[
 D=\operatorname{diag}(F_{13},F_{23}),\quad
 C=D^{-1}A^*G_N,\quad T=D^{-1}S D^{-*},\quad
 w=Cz=\binom{U-\theta}{V-\theta}.                              \tag{JRE7}
\]
The codomain of $C$ is $\mathbb C^2$ with its standard Hermitian
metric; its adjoint from that space to the original $G_N$-space is
$C^\sharp=AD^{-*}$. In particular $CC^\sharp=T$.
The same orthogonal argument applied to $C$ proves
\[
 w\in\operatorname{im}T,\qquad
 w^*T^\dagger w+\|z-C^\sharp T^\dagger w\|_{G_N}^{2}
             =E\theta(1-\theta).                              \tag{JRE8}
\]
The two minimum vectors in JRE6 and JRE8 agree, since they solve
equivalent observations and both lie in $\operatorname{im}A$.
This proves the exact congruence of the image metrics even when
$S$ has rank zero or one. It does not use an invalid unrestricted
congruence formula for a Moore--Penrose inverse.

\subsection{A correlated finite bound for the three projected errors}
Write $t=T_{12}$, retain
\[
 L=E\theta(1-\theta),\quad
 d=\sqrt{L(T_{11}+T_{22}-2\Re t)},\quad
 h=\frac L2\left(|\Im t|+
                   \sqrt{T_{11}T_{22}-(\Re t)^2}\right).       \tag{JRE9}
\]
The quantities under square roots are nonnegative by $T\succeq0$.
All full source masses remain in $E$ and $T$; the products in JRE9
are independent of a common rescaling of the inner product only
because their displayed factors cancel exactly.
For $w=(u_1,u_2)^\mathsf T$, JRE8 and Cauchy--Schwarz give
\[
 |u_2-u_1|\le d,\qquad
 |u_j|^2\le L T_{jj}.                                        \tag{JRE10}
\]
Here is a sharper joint estimate for the quadratic term, including
its complex phase. Let
\[
 J=\begin{pmatrix}0&-i/2\\i/2&0\end{pmatrix}.
\]
Then $w^*Jw=\Im(\bar u_1u_2)$.
Write $w=T^{1/2}\eta$ with $\eta$ in the image of $T$;
JRE8 gives $\|\eta\|^2\le L$. The Hermitian matrix
$T^{1/2}JT^{1/2}$ has trace $-\Im t$ and determinant
$-\det(T)/4$. Its two eigenvalues are consequently
\[
 \frac{-\Im t\ \pm\sqrt{T_{11}T_{22}-(\Re t)^2}}2 .           \tag{JRE11}
\]
The formula remains true at rank zero or one, since trace and
determinant still determine its characteristic polynomial.
The spectral theorem now proves
\[
 |\Im(\bar u_1u_2)|\le h.                                   \tag{JRE12}
\]
This is the exact largest absolute quadratic value on the complete
ellipse JRE8 with its residual term omitted: an eigenvector of the
larger absolute eigenvalue attains it. The actual observation projection
need not attain that ellipse boundary; no such attainability is claimed.

Use EC38's three original products
\[
 P_K=\bar U V,\qquad P_B=(1-\bar U)(1-V),\qquad
 P_\times=\bar U+V-2\bar U V .                                \tag{JRE13}
\]
Substitute $U=\theta+u_1$, $V=\theta+u_2$. Their imaginary parts are
exactly
\[
 \begin{split}
 \Im P_K&=\theta\Im(u_2-u_1)+\Im(\bar u_1u_2),\\
 \Im P_B&=-(1-\theta)\Im(u_2-u_1)+\Im(\bar u_1u_2),\\
 \Im P_\times&=(1-2\theta)\Im(u_2-u_1)-2\Im(\bar u_1u_2).
 \end{split}                                                  \tag{JRE14}
\]
Their sum is zero. JRE10 and JRE12 prove the three simultaneous bounds
\[
 \begin{split}
 |\Im P_K|&\le\theta d+h,\\
 |\Im P_B|&\le(1-\theta)d+h,\\
 |\Im P_\times|&\le|1-2\theta|d+2h .
 \end{split}                                                  \tag{JRE15}
\]
All three use the same $T$, the same $\theta$, and the same retained
class. In particular a small eigenvalue or a small original $F_{j3}$
is visible through the exact entries of $T$; it is never removed
by replacing the Gram with an identity or a rank fraction.

In EC39 let $\varepsilon_{N,\lambda}>0$ be its single shared phase
remainder and $E_{\rm cur}$ its proved finite bound. The exact
response-error vector, in the order $(K,B,\times)$, is
\[
 \mathcal R_X=2E(-1)^r\varepsilon_{N,\lambda}\Im P_X,
 \qquad \sum_X\mathcal R_X=0 .                               \tag{JRE16}
\]
Thus its individual absolute values are bounded, respectively, by
\[
 2EE_{\rm cur}(\theta d+h),\quad
 2EE_{\rm cur}((1-\theta)d+h),\quad
 2EE_{\rm cur}(|1-2\theta|d+2h).                              \tag{JRE17}
\]
This is a finite amplification estimate on the actual attained
subquotient metrics. It leaves the actual two-column Schur complement
and the original period-dependent kernel projection to be evaluated.
It does not assign a sign to an individual projected term from the
positive full-class sum. The next analytic task is now the joint
profile of $T$ and $\theta$ in these exact original coordinates.

\subsection{The actual parity map and its two boundary energies}
For the original even arithmetic measure on $S=a/2+iy$, reflection
$\mathcal R P(S)=P(a-S)$ is an isometry. The full root polynomial
satisfies $\chi(a-S)=(-1)^q\chi(S)$, so reflection preserves each
complete relation space $\chi\mathbb C[S]_{\le N-q}$. It therefore
induces an isometry of the attained quotient, commuting with its
minimum lift. The two original boundary columns have opposite
parities $(-1)^N$ and $(-1)^{N+1}$: their source monic orthogonal
polynomials do, by uniqueness of the monic minimum in the even
measure, and the remainder map commutes with reflection. Consequently
\[
 F_{12}=\langle b_N,b_{N+1}\rangle_{G_N}
       =-\langle b_N,b_{N+1}\rangle_{G_N}=0 .                   \tag{JRE18}
\]
This is the full quotient identity also calculated in NCE18--21.
It preserves the actual source, period, conductor and class. Inserting
it in JRE4 and JRE7 gives the exact original response matrix
\[
 T=\begin{pmatrix}F_{11}/|F_{13}|^2&0\\0&F_{22}/|F_{23}|^2\end{pmatrix}
       -\frac1E\begin{pmatrix}1&1\\1&1\end{pmatrix},\qquad
 T_{12}=-1/E .                                                \tag{JRE19}
\]
Define the two actual boundary energies
\[
 p_1=\frac{|F_{13}|^2}{EF_{11}},\qquad
 p_2=\frac{|F_{23}|^2}{EF_{22}}.                              \tag{JRE20}
\]
Both are positive. The two columns $b_j/\sqrt{F_{jj}}$ are
orthonormal by JRE18. The orthogonal square of their projection
of $v/\sqrt E$ is $p_1+p_2$, so $p_1+p_2\le1$ with the exact
remainder equal to the squared norm of its orthogonal complement.
In particular these are values of specified original projections,
not fractions assigned from ranks.
Substitute JRE19--20 into JRE9. Every original factor cancels
only inside its displayed matched ratio, yielding
\[
 \boxed{d^2=\theta(1-\theta)(p_1^{-1}+p_2^{-1}),\qquad
 h=\frac{\theta(1-\theta)}{2\sqrt{p_1p_2}}
                    \sqrt{1-p_1-p_2}.}                       \tag{JRE21}
\]
Together JRE17 and JRE21 reduce the actual amplification estimate
to the three original scalar projection energies
$p_1,p_2,\theta$. In particular the exact total allowance obeys
\[
 \sum_{X=K,B,\times}|\mathcal R_X|
 \le2EE_{\rm cur}\{(1+|1-2\theta|)d+4h\},\qquad
 \sum_X\mathcal R_X=0.                                      \tag{JRE22}
\]
The first statement follows by adding the three proved absolute
bounds, while the second is the unchanged signed identity JRE16.
No lower bound for either $p_j$ is inserted. Their possible smallness
remains explicitly in JRE21 and is a concrete target for the next
original-source calculation.

% END PRESERVED INPUT 08_JOINT_RESPONSE_ELLIPSE.tex
\endgroup

\clearpage
\begingroup


\setcounter{equation}{0}
\renewcommand{\theequation}{A02ARU.\arabic{equation}}
\section*{Exact backward propagation into the earlier programme}
\addcontentsline{toc}{section}{Exact backward propagation into the earlier programme}
% BEGIN PRESERVED INPUT 09_RECEIVER_UPDATES.tex SHA256 a5c2d9d83b8c08206e68305caca7026de7a1e9162d735823330bdd4d38716425
\section{Successor receiving identities on the original programme}

The preceding edition's MRI1--11 remain preserved. This section supplies
their more precise successors. The scalar statements below retain the
stipulated actual quartet, the full multiplicity $m_0$, $k=4\ell+1$, and
$q=[1+k(m_0-1)](k+1)^2$. The correlated native-row statements retain their
own simple-quartet degree $a=k+4$, dimension $q_a=(a+1)^2$, conductor
$\Delta=16a-48$, and the two original source orders. These are the
original domains of SR and CRS; no degree or source-order substitution
is made between them.

\subsection{The scalar, source order and complete action}
Write $d=4m_0-2$, $c_\zeta=2\gamma_E-\log(2\pi)$ and retain the physical
equilibrium constant $C_\Gamma$ and the complete-source pressure integral
$I_h$. The full SR1--30 proof gives, on its explicit finite guards,
\[
 \delta_{k,1}^{\rm ar}=-dC_\Gamma q+
 \frac{\log2}{\pi}I_hk^2+
 \frac{C_\Gamma q}{2(\log q+c_\zeta)}+R_{h,k}.
 \tag{ARU1}
\]
Here $R_{h,k}$ is exactly the remainder in SR29, with the SR30 bound
$o_h(q/\log q)$. Inserting ARU1 in the unchanged MRI3 identity proves
\[
\begin{split}
 J_k={}&\mathcal B_k^{(1)}+\mathcal W_k^{\rm mono}
 -\mathcal P_{k,-}-\mathcal P_{k,+}-dC_\Gamma q\\
 &+\frac{\log2}{\pi}I_hk^2+
 \frac{C_\Gamma q}{2(\log q+c_\zeta)}+R_{h,k}.
\end{split}\tag{ARU2}
\]
Subtracting the original source-order-one and source-order-$k$ identities
for this same $J_k$ gives the exact map
\[
 \delta_{k,k}^{\rm ar}=\delta_{k,1}^{\rm ar}
                   +\mathcal B_k^{(1)}-\mathcal B_k^{(k)}.
 \tag{ARU3}
\]
All source masses in the two baselines stay in this difference. On the
simple stratum define, as before,
$\mathfrak R_{h,k}=J_k-\delta_{k,1}^{\rm ar}-C_Bq^2-8q\log k$.
The previously proved estimate is $\mathfrak R_{h,k}=O_h(q)$.
Using this definition and ARU1, without absorbing either remainder,
proves
\[
\begin{split}
 J_k={}&C_Bq^2+8q\log k-2C_\Gamma q+
 \frac{\log2}{\pi}I_hk^2+
 \frac{C_\Gamma q}{2(\log q+c_\zeta)}\\
 &+\mathfrak R_{h,k}+R_{h,k}.
\end{split}\tag{ARU4}
\]
Consequently ARU1's improved arithmetic remainder does not improve
the separate analytic remainder $\mathfrak R_{h,k}$ by substitution.
Their exact relation is ARU4.

The full TS1--41 proof establishes $0<C_\Gamma<2/3$ and
$I_h>1681/(50\pi)$ from the actual theta source with its complete
denominator and outer mass. Its simple-stratum coefficient satisfies
\[
 d_0(h)=-2C_\Gamma+(\log2/\pi)I_h
       >7370383/7260000>1.
 \tag{ARU5}
\]
The finite guard TS40 therefore proves $\delta_{k,1}^{\rm ar}>q$.
For every fixed $m_0\ge2$, divide ARU1 by its literal $q$.
Then $k^2/q\to0$, $(\log q+c_\zeta)^{-1}\to0$ and
$R_{h,k}/q\to0$, proving the negative limit
$-(4m_0-2)C_\Gamma$. These are signs of the specified scalar on the
specified multiplicity strata. The exact projected receiver appears
below and retains its additional variables.

\subsection{The two path remainders and every marked frame}
Use the full two-parameter path
$F(t,r)=\mathcal B_k[\sigma e^{tH_C+rH_O}]$ of SR41--46.
Taylor's integral identity in $t$ and the fundamental theorem in $r$
give respectively
\[
\begin{split}
 F(1,1)-F(0,1)&=F_t(0,1)+
                  \int_0^1(1-t)F_{tt}(t,1)\,dt,\\
 F_t(0,1)-F_t(0,0)&=\int_0^1F_{tr}(0,r)\,dr.
\end{split}\tag{ARU6}
\]
Adding proves the original mixed and nonlinear receiver
\[
 \mathcal C_k^{\rm mixed}+\mathcal N_k^{\rm nonlinear}
 =F(1,1)-F(0,1)-F_t(0,0),\qquad
 |\mathcal C_k^{\rm mixed}+\mathcal N_k^{\rm nonlinear}|
 \le E_{\rm cen}+E_{\rm lin}=o_h(q/\log q).
 \tag{ARU7}
\]
The finite expressions for both errors, rather than just their order,
are SR5 and SR43. The matrices inside $F_{tr}$ retain the full
mixed trace SR44; the matrices inside $F_{tt}$ retain the full
projection variance SR45. Thus this bound applies to their displayed
sum with its original signs.

For an invertible original marked frame $\mathcal J$,
$\det(\mathcal J^*G_N[\nu]\mathcal J)=
 |\det\mathcal J|^2\det G_N[\nu]$. Apply this twice, with the same
$\mathcal J$ and the original arithmetic and Gamma measures. Division
proves
\[
 \log\frac{\det(\mathcal J^*G_N[m_k]\mathcal J)}
                 {\det(\mathcal J^*G_N[\sigma]\mathcal J)}
 =\log\frac{\det G_N[m_k]}{\det G_N[\sigma]}.
 \tag{ARU8}
\]
The same equality holds in each signed endpoint occurrence; summing
transports ARU1, with its unchanged error, into the marked relative
determinant. MAR supplies the exact regularized frame when the naive
frame degenerates. Its explicit inclusion and the attained metric are
used before ARU8; a singular determinant is never cancelled.
GMF supplies the growing-flag high/low comparison and retains its
terminal--low Schur block. FCD's fixed-degree expansion is used only
on its stated fixed-degree domain.

\subsection{Corrected native rows and their backward cancellation}
In the original CRS domain write $\gamma_r$ for the attained row
logarithm. The complete CRP and CRS derivations prove
\[
\begin{split}
 c_0&=\frac\Delta2\{1+\log(q_a/\Delta)\},\\
 c_r&=\Delta f(r/q_a)-\frac12
       \{(r+\Delta)\log(1+\Delta/r)-\Delta\}\quad(r\ge1),\\
 \sum_r|\gamma_r-c_r|&\le E_c=O_{\rm actual}(q_a).
\end{split}\tag{ARU9}
\]
The summation range and every finite guard are those of CRS17,
including the extension through $q_a+32a$. For weights $w_r$ equal
to the original endpoint weights $1,2,\ldots,2,1$, the triangle
inequality gives $|\sum w_r(\gamma_r-c_r)|\le2E_c$.
CRS23 retains the complete equilibrium integral for each prefix and
tail. At $R_*=\lfloor a^{3/2}/\log a\rfloor$, the deterministic
corrections to those two integrals are respectively
\[
 -64q_a\log a+128q_a\log\log a+O(q_a),\qquad
 -64q_a\log a-128q_a\log\log a+O(q_a).
 \tag{ARU10}
\]
Their sum is the existing $-128q_a\log a+O(q_a)$. Adding the
original lower-root $+128q_a\log a$ cancels it. This proves that the
more precise prefix/tail calculation introduces no extra
$q_a\log a$ term into ARU4. Retaining the old uncorrected
$\Delta f$ while merely reducing its error would omit exactly
the correction calculated in ARU10.

For the actual row projection $p_r$ and
$\kappa_r=-\log p_r$, retain $\kappa_r=+\infty$ when $p_r=0$.
CRP33 gives the two clipped expressions
\[
\begin{gathered}
 \widetilde S_K=\sum_r w_r\min(c_r^+,\kappa_r),\qquad
 \widetilde S_R=\sum_r w_r(c_r^+-\kappa_r)_+,\\
 |S_X-\widetilde S_X|\le2E_c+(\log2)\sum_r w_r\quad(X=K,R).
\end{gathered}
 \tag{ARU11}
\]
The pointwise equality
$\min(u,v)+(u-v)_+=u$ for $u,v\ge0$, also valid for $v=+\infty$
under the stated convention, proves
$\widetilde S_K+\widetilde S_R=\sum_rw_rc_r^+$.
The actual determinant identity remains $S_K+S_R=\mathcal T$.
The eight original signed flags retain their two ordered chains:
CRP28 gives the radial Lipschitz constant $4$ and row clipping
allowance $8\log2$; CRP29 gives $4\log2$ for each complementary
CPQ pair. Their original projection spectra are still present
in these formulas. This is the exact receiving map for their
next evaluation; no allocation is supplied by rank alone.
For the independent proper source, CRS27--30 likewise update
the displacement on its original four intervals, while retaining
its standard-window determinant.

\subsection{Complete retained class and projected arithmetic current}
NCB and NCRP prove the original-minor invertible coordinate completion
$F_I=[F_0,F_C,L_a]$ and the exact decomposition
$v_\lambda=F_0u_0+F_C\eta+L_az$.
For its full coordinate vector $w=(u_0,\eta,z)$, define
$H=F_I^*G_NF_I$ and $A_I=F_I^{-1}AF_I$. Matrix multiplication gives
\[
 w^*(A_I^*H+HA_I-aH)w
 =v_\lambda^*(A^*G_N+G_NA-aG_N)v_\lambda.
 \tag{ARU12}
\]
This identity includes all nine blocks of $H$. Expanding either side
after inserting the three summands proves equality term by term,
including both directed cross terms. It propagates the completed
class to the original arithmetic receiver without altering the native
section or its quotient maps.

For the literal observation projection $P$ and original boundary rows,
NCE proves
\[
 \ell_N=b_N^*G_N/\omega_N,\quad \xi_N=-G_Nb_{N+1},\qquad
 U=\frac{\ell_NPv_\lambda}{\ell_Nv_\lambda},\quad
 V=\frac{\xi_N^*Pv_\lambda}{\xi_N^*v_\lambda}.
 \tag{ARU13}
\]
Substituting the full three-component decomposition in each numerator
and denominator is exact because the two rows are linear. The
nonzero complementary components therefore remain in both ratios.
With the original current $c$ and $E=\|v_\lambda\|_{G_N}^2$, the
complete projected current is
\[
 \frac{\Phi_X}{2E}=\Re(cP_X),\quad
 (P_K,P_B,P_\times)=
 (\bar UV,(1-\bar U)(1-V),\bar U+V-2\bar UV),\quad
 P_K+P_B+P_\times=1.
 \tag{ARU14}
\]
Expanding the three products proves the last equality and hence
$\sum_X\Phi_X=2E\Re c$. This also proves the same-class cancellation
after the full-class repair, including both leakage directions.

JRE1--22 prove the exact joint image metric and its attained remainder.
Parity gives $F_{12}=0$ on this same quotient. Put
$p_j=|F_{j3}|^2/(EF_{jj})$, $\theta=K_{33}/E$ and retain the
original finite current allowance $E_{\rm cur}$. The resulting bound is
\[
\begin{gathered}
 d_*^2=\theta(1-\theta)(p_1^{-1}+p_2^{-1}),\qquad
 h_*={\theta(1-\theta)\sqrt{1-p_1-p_2}\over2\sqrt{p_1p_2}},\\
 (|\mathcal R_K|,|\mathcal R_B|,|\mathcal R_\times|)
 \le2EE_{\rm cur}
 (\theta d_*+h_*,(1-\theta)d_*+h_*,|1-2\theta|d_*+2h_*),
 \qquad \sum_X\mathcal R_X=0.
\end{gathered}\tag{ARU15}
\]
The componentwise bounds follow by inserting JRE21 in JRE17;
the cancellation follows from the imaginary part of ARU14.
The factors $p_1,p_2,\theta$ are original projection energies with
$p_1,p_2>0$, $p_1+p_2\le1$, $0\le\theta\le1$, proved in JRE20.
Their actual values and the actual directional spectra in ARU11
are the next analytic calculation. All earlier completed scalar,
period, marked-frame and finite-field calculations remain available
through the explicit maps above.

% END PRESERVED INPUT 09_RECEIVER_UPDATES.tex
\endgroup

\endgroup
